\documentclass[11pt]{article}

\usepackage[margin=1.1in]{geometry}

\usepackage{amssymb,mathtools,natbib}

% Anchored formal environments for causalsmith papers.
% First mandatory argument = obj_id (machine anchor joining the paper to the
% Lean crosswalk; invisible in the rendered PDF). Optional argument = name.
% Each environment also defines \label{obj:<obj_id>} so prose can use
% \cref{obj:T-1}; cleveref derives the printed kind and number from the target.
\usepackage{amsmath}
\usepackage{mathrsfs}
\usepackage{amsthm}
\usepackage{booktabs}
% Long displays may break across pages; let TeX stretch a little harder before an
% overfull \hbox so wide formulas/tables are less likely to run off the page.
\allowdisplaybreaks
\setlength{\emergencystretch}{3em}
\newtheorem{theorem}{Theorem}
\newtheorem{lemma}{Lemma}
\newtheorem{assumption}{Assumption}
\newtheorem{definition}{Definition}
\newtheorem{proposition}{Proposition}
\newtheorem{remark}{Remark}
\newtheorem{citedresult}{Cited result}
% The amsthm counter is `algorithmbox` (not `algorithm`) so a later `\usepackage{algorithm}` cannot clash.
\newcounter{algorithmbox}
\renewcommand{\thealgorithmbox}{\arabic{algorithmbox}}
\NewDocumentEnvironment{theoremv}{m o s}{\begin{theorem}\label{obj:#1}{\normalfont[\texttt{\detokenize{#1}}]\IfValueT{#2}{\ (#2)}\IfBooleanT{#3}{\verificationfootnotemark}\ }}{\end{theorem}}
\NewDocumentEnvironment{lemmav}{m o s}{\begin{lemma}\label{obj:#1}{\normalfont[\texttt{\detokenize{#1}}]\IfValueT{#2}{\ (#2)}\IfBooleanT{#3}{\verificationfootnotemark}\ }}{\end{lemma}}
\NewDocumentEnvironment{assumptionv}{m o s}{\begin{assumption}\label{obj:#1}{\normalfont[\texttt{\detokenize{#1}}]\IfValueT{#2}{\ (#2)}\IfBooleanT{#3}{\verificationfootnotemark}\ }}{\end{assumption}}
\NewDocumentEnvironment{definitionv}{m o s}{\begin{definition}\label{obj:#1}{\normalfont[\texttt{\detokenize{#1}}]\IfValueT{#2}{\ (#2)}\IfBooleanT{#3}{\verificationfootnotemark}\ }}{\end{definition}}
% A `propositionv` is a proved result tiered as a Proposition (theorem-grade, machine-verified).
\NewDocumentEnvironment{propositionv}{m o s}{\begin{proposition}\label{obj:#1}{\normalfont[\texttt{\detokenize{#1}}]\IfValueT{#2}{\ (#2)}\IfBooleanT{#3}{\verificationfootnotemark}\ }}{\end{proposition}}
% A `remarkv` holds NON-load-bearing interpretive / open-question content (e.g. a causalsmith `oeq:`
% open question). It is neither proved nor assumed; no theorem may depend on it.
\NewDocumentEnvironment{remarkv}{m o s}{\begin{remark}\label{obj:#1}{\normalfont[\texttt{\detokenize{#1}}]\IfValueT{#2}{\ (#2)}\IfBooleanT{#3}{\verificationfootnotemark}\ }}{\end{remark}}
% An `algorithmv` is a DEFINITION-kind object presented as a boxed, numbered-step procedure (an
% estimator / selection rule built in stages). Same anchor contract as `definitionv`; its body is
% ordinary LaTeX whose steps are `\begin{enumerate}` items, so tex2html needs no extra machinery.
% Presented in the CONVENTIONAL algorithm style readers expect from `algorithm`+`algorithmic`:
% a rule above, an "Algorithm N (Title)" caption line, a rule under the caption, the numbered
% steps, and a closing rule. It is NOT a float — the anchor contract requires the object to stay
% where the outline places it, and floats drift. The obj-id tag is suppressed here (it is
% bookkeeping, and a caption line carrying it reads as debug output in a finished paper).
\NewDocumentEnvironment{algorithmv}{m o s}%
  {\par\addvspace{\medskipamount}\noindent\rule{\linewidth}{0.8pt}\par\nobreak\vspace{2pt}%
   \refstepcounter{algorithmbox}\label{obj:#1}%
   \noindent\textbf{Algorithm \thealgorithmbox}\IfValueT{#2}{ \textnormal{(#2)}}%
   \IfBooleanT{#3}{\verificationfootnotemark}\par\nobreak\vspace{2pt}%
   \noindent\rule{\linewidth}{0.4pt}\par\nobreak\vspace{2pt}\normalfont}%
  {\par\nobreak\vspace{2pt}\noindent\rule{\linewidth}{0.8pt}\par\addvspace{\medskipamount}}
% A `citedv` is an IMPORTED external result the paper relies on but does NOT prove
% (a causalsmith `gate_class:"cited"` node). It is machine-checked against the cited
% source at F2.5, never proved here; the fixed provenance note makes that explicit
% so the environment can never be read as a contribution of this paper. The \citep
% attribution is supplied by the surrounding prose (P2), keyed to references.bib.
\NewDocumentEnvironment{citedv}{m o s}{\begin{citedresult}\label{obj:#1}{\normalfont[\texttt{\detokenize{#1}}]\IfValueT{#2}{\ (#2)}\IfBooleanT{#3}{\verificationfootnotemark}\ \textit{(imported result; machine-checked against the cited source, not proved here.)}\ }}{\end{citedresult}}
% \leanref{obj_id}{display text}: an inline first-mention link to a from-note object's
% verified Lean code. In the PDF it renders the display text plainly (the Lean link is a
% web-only affordance); tex2html turns it into a drawer-opening span.
\newcommand{\leanref}[2]{#2}
% For a theorem/lemma whose Lean declaration accepts a source-matched published proposition as an
% input, the starred anchored environment puts the mark in its heading; the generated text remains
% outside the frozen mathematical environment. Keep the combined macro for older paper bundles.
\newcommand{\verificationfootnotemark}{\footnotemark}
\newcommand{\verificationfootnotetext}[1]{\footnotetext{#1}}
\newcommand{\verificationfootnote}[1]{\verificationfootnotemark\verificationfootnotetext{#1}}
% xcolor before hyperref (hyperref would otherwise pull in plain `color` for colorlinks, and a
% later xcolor load clashes). The page-1 identity stamp at the end of this file needs a grey.
\usepackage{xcolor}
\usepackage{graphicx}
\usepackage{eso-pic}
% hyperref follows natbib and the environment macros; cleveref must follow hyperref.
\usepackage[colorlinks=true,linkcolor=blue,citecolor=blue,urlcolor=blue]{hyperref}
\usepackage[capitalise,nameinlink,noabbrev]{cleveref}
% Pin the names explicitly: labels are placed inside the underlying amsthm environments, and these
% declarations make the PDF contract independent of cleveref's theorem-name inference. House style:
% reference names always print with a capital first letter ("Theorem 1", "Definition 2"), in any
% sentence position — hence `capitalise` above and capitalized \crefname forms below.
\crefname{theorem}{Theorem}{Theorems}
\Crefname{theorem}{Theorem}{Theorems}
\crefname{lemma}{Lemma}{Lemmas}
\Crefname{lemma}{Lemma}{Lemmas}
\crefname{assumption}{Assumption}{Assumptions}
\Crefname{assumption}{Assumption}{Assumptions}
\crefname{definition}{Definition}{Definitions}
\Crefname{definition}{Definition}{Definitions}
\crefname{proposition}{Proposition}{Propositions}
\Crefname{proposition}{Proposition}{Propositions}
\crefname{remark}{Remark}{Remarks}
\Crefname{remark}{Remark}{Remarks}
\crefname{citedresult}{Cited result}{Cited results}
\Crefname{citedresult}{Cited result}{Cited results}
\crefname{algorithmbox}{Algorithm}{Algorithms}
\Crefname{algorithmbox}{Algorithm}{Algorithms}
% ---------------------------------------------------------------------------
% Page-1 identity stamp (arXiv style) and the pinned title-block date.
%
% Split of responsibility: THIS file owns the FORMAT (placement, font, colour, punctuation),
% the generated `paper_stamp.tex` owns only the VALUES (number+version, area, revised date,
% DOI, link target). P4 refreshes this template on every emit, so a layout fix reaches every
% bundle from a recompile; and because `paper_stamp.tex` is not one of the P2 assembly
% digest's authored inputs, a DOI that arrives after publication needs only that one small
% file rewritten plus a recompile — never a redraft, never a reassembly.
%
% Defaults first, so a bundle WITHOUT a stamp file compiles exactly as it always did:
% `\cspaperdate` falls back to `\today` and no stamp is drawn.
\providecommand*{\cspaperdate}{\today}  % the \date{} target
\providecommand*{\csstampid}{}          % e.g. CSWP-2026-008v3 (empty ⇒ no stamp at all)
\providecommand*{\csstamparea}{}        % e.g. Stat
\providecommand*{\csstampdate}{}        % e.g. 27 Aug 2026
\providecommand*{\csstampdoi}{}         % e.g. 10.5281/zenodo.123 (empty ⇒ DOI part omitted)
\providecommand*{\csstamplink}{}        % href target (DOI url, else the paper's page; may be empty)
\IfFileExists{paper_stamp.tex}{% GENERATED by CausalSmith P4 — paper identity stamp values. Do not hand-edit.
%
% Field order on the stamp (owned by paper_macros.tex):
%   <number>v<version> · doi:<doi> · [<Area>] <date>   — the DOI is never the last token, so a
%   text extractor cannot run it into whatever follows. Without a DOI that part is omitted.
% Format lives in paper_macros.tex; only the values are here, and this file is NOT one of
% the P2 assembly digest's authored inputs. A DOI arriving after publication therefore needs
% this file rewritten plus a `latexmk` recompile — never a redraft or a reassembly.
\renewcommand*{\csstampid}{CSWP-2026-028v4}
\renewcommand*{\csstamparea}{Stat}
\renewcommand*{\csstampdate}{1 Oct 2026}
\renewcommand*{\csstampdoi}{}
\renewcommand*{\csstamplink}{https://causalsmith.org/papers/stat_private_cate_roughdesign_frontier_v1}
\renewcommand*{\cspaperdate}{1 October 2026}
}{}
\makeatletter
% Field order is deliberate: the DOI is NEVER the last token on the line. A trailing identifier
% runs into whatever a text extractor emits next, which makes it unsafe to match mechanically
% (the Zenodo stamp matcher reads exactly this line); the date is a safe terminator.
%   <number>v<version> · doi:<doi> · [<Area>] <date>      — with a DOI
%   <number>v<version> · [<Area>] <date>                  — without
\newcommand{\cs@stampsep}{\,\textperiodcentered\,}
\newcommand{\cs@stamptext}{%
  \csstampid
  \ifx\csstampdoi\@empty\else\cs@stampsep doi:\csstampdoi\fi
  \ifx\csstamparea\@empty\else\cs@stampsep[\csstamparea]\fi
  \ifx\csstampdate\@empty\else\quad\csstampdate\fi
}
\ifx\csstampid\@empty\else
  \definecolor{csstampgrey}{gray}{0.45}
  % Background shipout material: it occupies no space, so the text block and the \thanks
  % footnote are byte-identical to an unstamped compile. The starred form applies to the
  % NEXT page shipped only — i.e. page 1. Placed in the left margin (the text block starts
  % at 1.1in), reading bottom-to-top, in a Times-like face at 8pt.
  \AddToShipoutPictureBG*{%
    \AtPageLowerLeft{%
      \put(\LenToUnit{0.42in},\LenToUnit{1.1in}){%
        \rotatebox{90}{%
          \fontfamily{ptm}\fontsize{8}{10}\selectfont
          \ifx\csstamplink\@empty
            \textcolor{csstampgrey}{\cs@stamptext}%
          \else
            \expandafter\href\expandafter{\csstamplink}{\textcolor{csstampgrey}{\cs@stamptext}}%
          \fi
        }%
      }%
    }%
  }
\fi
\makeatother


\title{Private estimation and honest inference for conditional treatment effects with rough nuisance functions}

\author{CausalSmith\thanks{Machine-generated by the CausalSmith research pipeline.}}

\date{\cspaperdate}

\begin{document}

\maketitle

\begin{abstract}
We characterize the optimal accuracy of pure differentially private pointwise treatment-effect estimation and honest interval inference from independent binary observational records. The model has a scalar covariate on \([0,1]\), an unknown measurable density between \(1/2\) and \(3/2\), treatment probabilities between \(1/4\) and \(3/4\), propensity and control regressions of H\"older order \(1/10\), and a Lipschitz treatment contrast, all with regularity radius \(3\). Within this model, consistency and conditional exchangeability identify the causal contrast almost surely; continuity and positive design density identify its value at the fixed covariate location \(1/2\). For sample size \(n\ge2\) and pure replacement privacy budget \(0<\epsilon\le1\), both minimax expected absolute error and minimax maximal expected length of connected intervals with uniform finite-sample \(90\%\) coverage have order
\[
\min\left\{1,\max\left\{
n^{-1/4},(n^2\epsilon)^{-1/7},(n\epsilon)^{-1/2}
\right\}\right\}.
\]
A count-weighted within-cell pair ratio attains these orders through one joint Laplace release and public tuning. Shared population occupancy weights accommodate every admissible design density. Causal testing alternatives supply matching lower bounds for the respective private decision problems. Along admissible privacy-budget sequences, vanishing optimal risk and honest expected length are equivalent to \(n\epsilon\to\infty\).
\end{abstract}

\section{Introduction}

This paper establishes matching finite-sample orders for private estimation and honest interval inference for a treatment effect at a fixed covariate value. The experiment consists of independent binary observational records with a scalar covariate, fixed treatment overlap, a bounded positive design density, rough nuisance regressions, and a Lipschitz treatment-effect contrast. Within this experiment, a single joint private release attains both the minimax expected absolute-error order and the optimal order of maximal expected length of connected intervals with uniform \(90\%\) coverage. The matching conclusions concern the dependence on sample size and privacy budget, with explicit numerical constants.

Pointwise treatment effects describe how the expected response to treatment varies across observed characteristics. Under consistency and conditional exchangeability, observational treatment-arm regressions identify the conditional potential-outcome contrast. Pointwise evaluation additionally requires a choice of regression version. In the class studied here, contrast continuity and positive covariate density determine a unique continuous representative, making its value at the public evaluation location an identified feature of the observed-record law, as established in \cref{obj:lem:point-version}. Privacy then constrains how that feature can be estimated and reported from confidential records.

The statistical class in \cref{obj:def:complete-model} fixes the covariate domain to \([0,1]\), the evaluation location to \(1/2\), and a common regularity radius of \(3\). It combines a measurable density bounded above and away from zero, treatment overlap, one-tenth H\"older propensity and control regressions, and a Lipschitz contrast for binary outcomes. Releases satisfy pure central differential privacy under replacement adjacency on the full ordered-dataset space.

The smooth contrast permits the two treatment-arm means to share rough local variation that cancels in their difference, while rough treatment selection remains an additional nuisance. This scalar binary experiment makes the interaction between nuisance roughness, local occupancy, and privacy explicit.

Let \(n\) denote sample size and \(\epsilon\) the privacy budget. For every integer \(n\ge2\) and \(0<\epsilon\le1\), define the accuracy benchmark
\[
r(n,\epsilon)=
\min\left\{1,\max\left\{
n^{-1/4},(n^2\epsilon)^{-1/7},(n\epsilon)^{-1/2}
\right\}\right\}.
\]
\Cref{obj:thm:matched-risk-frontier} establishes that minimax maximal expected absolute error is bounded above and below by numerical multiples of this benchmark. \Cref{obj:thm:sharp-interval-frontier} establishes the same order for minimax maximal expected length among private connected-interval procedures satisfying uniform finite-sample \(90\%\) coverage. Both criteria integrate sampling variation and release randomization. The lower bounds range over all everywhere-defined private randomized procedures admitted by the respective criteria.

The attaining construction uses a public window around the evaluation location, partitioned into equal cells. Within each cell, treatment--outcome pair differences form a numerator and treatment pair differences form a denominator. Count weighting gives their population means the same occupancy weights, so the population ratio preserves the cellwise approximation to the target across every admissible measurable design density. The joint statistic has uniformly bounded replacement sensitivity, including changes involving cells with fewer than two records. Independent Laplace perturbations, a public denominator floor, and clipping therefore yield an everywhere-defined private ratio. Public tuning balances localization bias, within-cell nuisance approximation, pair variation, and privacy noise. The uniform squared-error envelope in \cref{obj:thm:uniform-private-upper} supplies an honest interval radius, and candidate-value inversion in \cref{obj:def:interval-handle} reproduces that interval from the same joint release.

The converse uses causal alternatives satisfying the same model restrictions. Shared-sign alternatives vary the nuisance functions while maintaining a common treatment-effect separation; their mixture permits likelihood and component-coupling comparisons. Direct localized alternatives provide the complementary privacy comparisons. The finite testing priors in \cref{obj:lem:frontier-testing-priors} combine these experiments across budget regimes. For scalar estimation, accurate recovery of separated targets would distinguish their released-output laws. For honest intervals, transferred coverage probabilities force simultaneous containment of separated targets with positive probability, and connectedness converts that containment into a length lower bound. This interval argument establishes the expected-length converse directly.

The closest pointwise comparison is \citet[][Theorems 1--2 and Remarks 2 and 11]{KennedyBalakrishnanRobinsWasserman2024}. At dimension one, contrast smoothness one, and average nuisance smoothness one tenth, their rate specializes to \(n^{-1/4}\), with the nuisance average below the relevant one-sixth threshold. For the binary bounded-positive-density experiment here, shared occupancy weighting attains this statistical scale together with the privacy scales \((n^2\epsilon)^{-1/7}\) and \((n\epsilon)^{-1/2}\). The reusable contribution is the combination of density-robust occupancy-weighted attainment with matching private testing comparisons; the connected-interval converse then transfers that frontier directly to honest expected length. \Cref{sec:related-work} gives the detailed comparison of experiments, assumptions, and adjacent private and inferential literature.

The benchmark gives an order-based guide to sample-size and budget choices. As the budget decreases, its dominant term moves from statistical variation to sparse-pair privacy noise, then dense-occupancy privacy noise, and finally bounded-range saturation. \Cref{obj:prop:regime-and-sanity} identifies the transition budgets and establishes that vanishing private minimax absolute risk along admissible budget sequences is equivalent to \(n\epsilon\to\infty\). The interval comparison gives the same condition for vanishing optimal honest expected length.



The paper proceeds through related work, the observation model and decision criteria, and the within-cell private release. It then presents the matching estimation and honest interval frontiers, followed by discussion and limitations. The appendices develop identification and upper-bound ingredients, lower-bound experiments and component coupling, testing reductions, and proofs of the main results.


\section{Related work}\label{sec:related-work}

The closest comparison is with \citet[][Theorems 1--2 and Remarks 2 and 11]{KennedyBalakrishnanRobinsWasserman2024}, who characterize pointwise conditional average treatment effect estimation under separate H\"older restrictions on the propensity score, control regression, and treatment contrast. For their formula, dimension one, contrast smoothness one, and average nuisance smoothness one tenth give the rate \(n^{-1/4}\); one tenth is below the corresponding one-sixth threshold. The present choice therefore isolates a simple rough-nuisance regime in which higher-order bias control is relevant. Their lower bound applies to a broader experiment allowing real outcomes; its construction uses a binary-outcome submodel and a covariate density bounded above that can vanish locally. Their higher-order local-polynomial estimator attains the stated rate under additional approximation, eigenvalue, boundedness, and covariate-distribution estimation conditions. Remark 11 expresses the last requirement through a third-order product involving design estimation and nuisance errors. The present experiment in \cref{obj:def:complete-model} combines binary outcomes, a scalar covariate, a density bounded above and away from zero, fixed overlap, one-tenth H\"older nuisance functions, and a Lipschitz treatment contrast. Within this class, \cref{obj:thm:matched-risk-frontier} gives matching orders for private pointwise expected absolute error, with additional scales \((n^2\epsilon)^{-1/7}\) and \((n\epsilon)^{-1/2}\). The comparison therefore concerns both the statistical experiment and the attainment argument: here, occupancy-weighted within-cell pairs support private attainment under the stated density bounds.
% Source: [Minimax rates for heterogeneous causal effect estimation](https://arxiv.org/pdf/2203.00837).

This construction connects to the literature on private U-statistics. \citet[][Theorem 2]{ChaudhuriLohPandeySarkar2024} give a pure differentially private estimator for the expectation of a bounded symmetric statistic, with a probability bound involving its variance, concentration of local H\'ajek projections, and privacy terms. The release in \cref{obj:def:private-ratio} combines two count-weighted pair statistics within a public partition, adds Laplace noise jointly, and forms a clipped ratio with a public denominator floor. Its analysis addresses the random number of observations contributing pairs and the normalization needed for a local treatment contrast. The shared occupancy weights in the population numerator and denominator are central to \cref{obj:lem:occupancy-cancellation}, while \cref{obj:lem:all-count-stability} controls replacement sensitivity across cell counts, including transitions involving cells with fewer than two observations. Together, these features support the uniform guarantee in \cref{obj:thm:uniform-private-upper}.
% Source: [On differentially private U statistics](https://proceedings.neurips.cc/paper_files/paper/2024/file/290848c892a85a6936b473d4daa6d0ad-Paper-Conference.pdf).

The lower-bound analysis draws on coupling-based private testing. \citet[][Section 2]{AcharyaSunZhang2021} relate private distinguishability to the expected Hamming distance between datasets under a coupling, including comparisons between mixtures of product laws. The present analysis specializes this approach to causal alternatives satisfying the binary experiment and its smoothness restrictions. The common-mass construction in \cref{obj:def:common-mass-coupling} operates on conditional components of the shared-sign experiment, and \cref{obj:lem:measurable-component-contraction,obj:lem:private-kernel-comparison} connect these comparisons to private output laws. The separated testing priors in \cref{obj:lem:frontier-testing-priors} then supply the experiments used for the estimation and interval lower bounds. The contribution lies in constructing and controlling these comparisons within the specified causal class.
% Source: [Differentially private Assouad, Fano, and Le Cam](https://proceedings.mlr.press/v132/acharya21a/acharya21a.pdf).

Private heterogeneous-effect learners provide a complementary perspective. \citet{NiuEtAl2022} develop a general procedure for conditional average treatment effect learning under approximate differential privacy, using sample splitting and private base learners. Their implementations include single-stage and multistage estimators, and their empirical performance criterion is test-set mean squared error, with a bias--variance decomposition. \citet[][Theorem 1]{SchroderMelnychukFeuerriegel2025} instead establish approximate differential privacy for a prespecified finite vector of conditional treatment effect queries through Gaussian output perturbation calibrated to the second-stage learner's sensitivity. Their accompanying performance analysis concerns Neyman orthogonality and quasi-oracle estimation under nuisance estimation conditions \citep[][Theorem 2]{SchroderMelnychukFeuerriegel2025}. The present estimand is the treatment contrast at a fixed covariate location, privacy is pure differential privacy under replacement adjacency as defined in \cref{obj:def:private-kernels}, and performance is measured by maximal expected absolute error and honest maximal expected interval length in \cref{obj:def:risk-criterion,obj:def:interval-criterion}. These choices make the comparisons specific to the estimand, privacy requirement, and decision criterion.
% Sources: [Differentially private estimation of heterogeneous causal effects](https://proceedings.mlr.press/v177/niu22a/niu22a.pdf), [Differentially private learners for heterogeneous treatment effects](https://arxiv.org/html/2503.03486v1).

The methodological background includes residualization in semiparametric regression \citep{Robinson1988}, orthogonal scores and cross-fitting \citep{Chernozhukov2018}, and two-stage conditional treatment effect estimation \citep{Nie2021,Kennedy2023}. Higher-order influence functions supply a general framework for controlling nonlinear nuisance bias \citep{Robins2008}. Causal forests and generalized random forests provide adaptive local estimation and asymptotic inference under their respective regularity conditions \citep{WagerAthey2018,Athey2019}. These approaches explain the importance of separating contrast regularity from nuisance complexity; the present pair construction gives a concrete private implementation for the rough-nuisance experiment in \cref{obj:def:complete-model}.
% Sources: [Root-N-consistent semiparametric regression](https://doi.org/10.2307/1912705), [Double/debiased machine learning for treatment and structural parameters](https://doi.org/10.1111/ectj.12097), [Quasi-oracle estimation of heterogeneous treatment effects](https://arxiv.org/abs/1712.04912), [Towards optimal doubly robust estimation of heterogeneous causal effects](https://arxiv.org/abs/2004.14497), [Higher order influence functions and minimax estimation of nonlinear functionals](https://arxiv.org/abs/0805.3040), [Estimation and inference of heterogeneous treatment effects using random forests](https://arxiv.org/abs/1510.04342), [Generalized random forests](https://imstat.org/publications/aos/aos_47_2/aos_47_2.pdf).

Smooth-nuisance kernel results form a separate model comparison. \citet[][Sections 2, 5, and 6]{Kim2026} study conditional treatment effects under reproducing kernel Hilbert space restrictions on response functions, with additional spectral, source, or dimension structure for the contrast. Their kernel ridge regression analysis relates estimation rates to contrast complexity and overlap within those models. Here, regularity is specified directly through the H\"older and Lipschitz conditions in \cref{obj:def:complete-model}, and the private risk comparison is established for that class. The distinction between these function spaces and performance criteria governs the interpretation of their respective optimality results.
% Source: [Optimal and structure-adaptive CATE estimation with kernel ridge regression](https://arxiv.org/html/2602.18958v1).

Finally, the interval criterion connects private causal inference to honest nonparametric inference. \citet{ArmstrongKolesar2018} derive finite-sample optimal intervals and efficiency bounds for linear regression functionals under Gaussian errors with known variance, while \citet{ArmstrongKolesar2020} develop intervals whose critical values account explicitly for worst-case bias. The present construction likewise incorporates a public error envelope, and \cref{obj:thm:sharp-interval-frontier} establishes matching orders, up to constant factors, for maximal expected length subject to uniform finite-sample \(90\%\) coverage over the causal class and pure replacement privacy. A complementary private causal inference result is provided by \citet[][Sections 4.2--4.3]{SchroderHartensteinFeuerriegel2025}, who construct approximate differentially private intervals for the average treatment effect using asymptotic normality and accounting for uncertainty from estimation and privatization. The present expected-length result places pointwise inference for a Lipschitz contrast with rough nuisance functions within a uniform finite-sample decision framework.
% Sources: [Optimal inference in a class of regression models](https://www.princeton.edu/~mkolesar/papers/optimal.pdf), [Simple and honest confidence intervals in nonparametric regression](https://onlinelibrary.wiley.com/doi/10.3982/QE1199), [PrivATE: Differentially private confidence intervals for average treatment effects](https://arxiv.org/html/2505.21641v1).


\section{Observation model, identification, and decision criteria}

Generic constants \(c\) and \(C\) denote positive numerical constants whose values may change across occurrences. Absolute values denote scalar magnitudes, and \(a\asymp b\) means that \(a\) is bounded above and below by positive numerical multiples of \(b\). For sample size \leanref{sym:n}{\ensuremath{n}}, record indices run from \(1\) to \(n\); sums over empty index sets equal zero. Expectations and probabilities with procedure subscripts integrate the procedure's randomization as well as the observed sample.

\subsection{Causal records and observational sampling}

The full causal record is \((X,A,(Y(0),Y(1)),Y)\), where \(X\in[0,1]\) is the scalar covariate, \(A\in\{0,1\}\) is treatment, and the potential and observed outcomes are binary. Its observed marginal retains \((X,A,Y)\), and an ordered sample consists of \(n\) independent records from that marginal. The anchored definitions immediately below record the coordinate, marginal-law, and dataset conventions used by the statistical class and decision problems.

The model uses a common public regularity radius \leanref{sym:L}{\ensuremath{L}} for the nuisance functions and the treatment-effect contrast. Its numerical value is fixed throughout the analysis.

\begin{definitionv}{synth_4}[Fixed public regularity radius]
The public regularity radius \(L\) is the real constant
\[
L = 3.
\]
\end{definitionv}

The causal record contains the covariate \leanref{sym:X}{\ensuremath{X}}, treatment indicator \leanref{sym:A}{\ensuremath{A}}, potential-outcome pair \leanref{sym:Ypot}{\ensuremath{Y^{\mathrm{pot}}}}, and observed outcome \leanref{sym:Y}{\ensuremath{Y}}. The following coordinate conventions distinguish the observed variables from the potential outcomes.

\begin{definitionv}{synth_9}[Observed records and ordered datasets]
The observed-record space is
\[
\mathcal O=[0,1]\times\{0,1\}^2,
\]
equipped with its Borel sigma-algebra, with a record written as \(z=(x,a,y)\). For a nonnegative integer \(n\), define \(\mathcal O^n\) to be its ordered \(n\)-fold product, equipped with the product sigma-algebra. A dataset is an element
\[
D=(z_1,\ldots,z_n)\in\mathcal O^n.
\]
For a random dataset \(D\), write \(\mathcal L(D)\) for its probability law on \(\mathcal O^n\); explicitly,
\[
\mathcal L(D)(E)=\Pr\{D\in E\}
\]
for every measurable \(E\subseteq\mathcal O^n\). For \(n=0\), the product consists of the empty dataset.
\end{definitionv}

\begin{definitionv}{synth_3}[Observed binary outcome coordinate]
The observed binary outcome \(Y\) is the final coordinate of a full causal record, whose coordinates are, in order, a covariate, a binary treatment indicator, an ordered pair of binary potential outcomes, and a binary observed outcome.
\end{definitionv}

\begin{definitionv}{synth_2}[Binary treatment coordinate]
The binary treatment indicator \(A\) maps each full causal record to its treatment coordinate, taking values in \(\{0,1\}\). A full causal record is ordered as covariate, binary treatment, ordered pair of binary potential outcomes, and binary observed outcome.
\end{definitionv}

A causal law \leanref{sym:P}{\ensuremath{P}} specifies the joint distribution of the full record. Its observational marginal \leanref{sym:Pobs}{\ensuremath{P^{\mathrm{obs}}}} specifies the distribution available to an estimator; the covariate marginal \(P_X\) describes the random design.

\begin{definitionv}{synth_1}[Scalar covariate coordinate]
The scalar covariate \(X\) is the first coordinate of the causal record, with values in the closed unit interval:
\[
X \in [0,1].
\]
\end{definitionv}

Write \(Y(A)=(1-A)Y(0)+AY(1)\) for the potential outcome selected by the realized treatment. The observed-record space \leanref{sym:O}{\ensuremath{\mathcal O}} retains the covariate, treatment, and observed outcome, and an ordered dataset \leanref{sym:D}{\ensuremath{D}} collects these records.

\begin{definitionv}{synth_8}[Causal coordinates and marginal laws]
A causal law \(P\) is a Borel probability law of
\[
(X,A,Y^{\mathrm{pot}},Y)
\in[0,1]\times\{0,1\}\times\{0,1\}^2\times\{0,1\},
\]
where \(X\) is the covariate, \(A\) is the treatment indicator, \(Y\) is the observed outcome, and
\[
Y^{\mathrm{pot}}=(Y(0),Y(1))
\]
is the ordered pair of control and treatment potential outcomes. Define the observational marginal and the covariate marginal by
\[
P^{\mathrm{obs}}=\mathcal L_P(X,A,Y),
\qquad
P_X=\mathcal L_P(X).
\]
Here \(\mathcal L_P\) denotes the probability law of the indicated coordinates under \(P\).
\end{definitionv}

Sampling and causal identification rest on the following three conditions. The first specifies how observed records enter the dataset.

\begin{assumptionv}{ass:iid}[Independent and identically distributed observations]
The dataset \(D\) has distribution
\[
\mathcal L(D)=(P^{\mathrm{obs}})^{\otimes n}.
\]
\end{assumptionv}

Independent and identically distributed observational sampling makes the common observed-record law the statistical experiment for estimation and inference \citep{KennedyBalakrishnanRobinsWasserman2024}. The next condition connects each observed outcome to the corresponding potential outcome.

\begin{assumptionv}{ass:consistency}[Consistency of observed outcomes]
The observed outcome \(Y\) equals the potential outcome \(Y(A)\) selected by the realized treatment with probability one:
\[
P\{Y=Y(A)\}=1.
\]
\end{assumptionv}

Consistency assigns the realized outcome to the treatment actually received, giving the observed outcome regression its potential-outcome interpretation \citep{KennedyBalakrishnanRobinsWasserman2024}. Identification further requires that the covariate account for treatment selection in the following sense.

\begin{assumptionv}{ass:exchangeability}[Conditional exchangeability of treatment]
The potential-outcome pair \(Y^{\mathrm{pot}}\) is conditionally independent of the treatment indicator \(A\) given the covariate \(X\) under the causal law \(P\):
\[
Y^{\mathrm{pot}}\perp\!\!\!\perp A\mid X\quad[P].
\]
\end{assumptionv}

Conditional exchangeability permits comparison of the two treatment arms at a common covariate value \citep{KennedyBalakrishnanRobinsWasserman2024}. Together with consistency and the overlap restriction below, it identifies the conditional mean potential-outcome contrast from observational regressions, initially up to covariate-null sets.

\subsection{Selected regressions and the statistical class}

Pointwise estimation requires explicit regression versions. We therefore specify the selected design density \leanref{sym:f}{\ensuremath{f_P}}, propensity regression \leanref{sym:e}{\ensuremath{e_P}}, and control outcome regression \leanref{sym:mu0}{\ensuremath{\mu_{0,P}}} before imposing their regularity conditions.

\begin{definitionv}{synth_10}[Selected design density and observational regressions]
For a causal law \(P\) whose covariate marginal \(P_X\) is absolutely continuous with respect to Lebesgue measure \(dx\) on \([0,1]\), let \(f_P\) be the selected measurable version of its density:
\[
f_P=\frac{dP_X}{dx},
\qquad
P_X(B)=\int_B f_P(x)\,dx
\]
for every Borel set \(B\subseteq[0,1]\). Let \(e_P\) and \(\mu_{0,P}\) be the selected versions of the propensity regression and the control outcome regression, respectively:
\[
e_P(x)=P(A=1\mid X=x),
\qquad
\mu_{0,P}(x)=E_P[Y\mid X=x,A=0],
\qquad x\in[0,1].
\]
These conditional expressions denote the selected regression functions on \([0,1]\). The model's overlap and smoothness conditions apply to these selected versions, and its density bounds apply to \(f_P\) Lebesgue-almost everywhere.
\end{definitionv}

The selected treated outcome regression \leanref{sym:mu1}{\ensuremath{\mu_{1,P}}} completes the observational contrast \leanref{sym:tau}{\ensuremath{\tau_P}}. We use the selected conditional versions throughout and define
\[
\mu_{1,P}(x)=E_P[Y\mid X=x,A=1],
\qquad
\tau_P(x)=\mu_{1,P}(x)-\mu_{0,P}(x),
\qquad
x_0=\frac12,
\qquad
\theta(P)=\tau_P(x_0).
\]
Thus the fixed public evaluation location \leanref{sym:x0}{\ensuremath{x_0}} determines the scalar target \leanref{sym:theta}{\ensuremath{\theta(P)}}.



Binary outcomes place both the contrast and the target in \([-1,1]\). We next organize the random-design, overlap, and smoothness restrictions that determine the statistical class.

\begin{assumptionv}{ass:density}[Bounded covariate density]
The covariate density \(f_P\) satisfies
\[
\frac12\le f_P(x)\le\frac32
\quad\text{for Lebesgue-almost every }x\in[0,1].
\]
\end{assumptionv}

The bounded positive random-design density ensures that covariate neighborhoods receive probability comparable to their Lebesgue length. These fixed numerical bounds instantiate the density condition in \citep[Theorem 1, footnote 6]{WagerAthey2018}. Treatment probabilities are likewise bounded away from their endpoints.

\begin{assumptionv}{ass:overlap}[Uniform treatment overlap]
The conditional treatment probability \(e_P\) satisfies
\[
\frac14\le e_P(x)\le\frac34
\qquad\text{for every }x\in[0,1].
\]
\end{assumptionv}

Uniform treatment overlap ensures that both treatment arms remain represented throughout the covariate domain; the constants give a fixed numerical version of the standard overlap condition \citep{KennedyBalakrishnanRobinsWasserman2024}. The nuisance regressions may vary substantially over short distances, subject to the following common Hölder order.

\begin{assumptionv}{ass:propensity-holder}[Hölder continuity of propensity]
The conditional treatment probability \(e_P\) satisfies the Hölder bound with common smoothness radius \(L\):
\[
|e_P(x)-e_P(x^{\prime})|
\le L|x-x^{\prime}|^{1/10}
\qquad\text{for every }x,x^{\prime}\in[0,1].
\]
\end{assumptionv}

Hölder propensity regularity controls local changes in treatment selection while allowing a rough propensity function of order \(1/10\) \citep{KennedyBalakrishnanRobinsWasserman2024}. The control regression obeys the corresponding restriction.

\begin{assumptionv}{ass:control-holder}[Hölder continuity of control means]
For every \(x,x'\in[0,1]\),
\[
|\mu_{0,P}(x)-\mu_{0,P}(x')|
\le L|x-x'|^{1/10}.
\]
\end{assumptionv}

Hölder control-regression regularity bounds local changes in the baseline outcome mean at the same rough order \citep{KennedyBalakrishnanRobinsWasserman2024}. The treatment-effect contrast has greater regularity than these nuisance functions.

\begin{assumptionv}{ass:contrast-lipschitz}[Lipschitz continuity of contrasts]
For every \(x,x'\in[0,1]\),
\[
|\tau_P(x)-\tau_P(x')|
\le L|x-x'|.
\]
\end{assumptionv}

Lipschitz contrast regularity controls the change in the conditional treatment effect over a covariate neighborhood. It is the Hölder order-one condition under the convention used in \citep{KennedyBalakrishnanRobinsWasserman2024}. Thus the model combines a Lipschitz target function with propensity and control regressions of Hölder order \(1/10\), all within the fixed radius in \cref{obj:synth_4}.

We collect these statistical restrictions into a class of causal laws. The private-release requirement is introduced afterward, once its output-space notation is in place.

\begin{definitionv}{synth_16}[Borel probability laws]
For a space \(\mathsf X\) equipped with its Borel \(\sigma\)-algebra \(\mathcal B(\mathsf X)\), define
\[
\operatorname{Prob}(\mathsf X)
=
\{P:P\text{ is a probability measure on }(\mathsf X,\mathcal B(\mathsf X))\}.
\]
\end{definitionv}

The statistical class includes the identifying conditions and the design and function restrictions just stated.

\begin{definitionv}{def:complete-model}[Causal law class \(\mathcal P\)]
\(\mathcal P\) is the class of causal probability laws
\[
P\in\operatorname{Prob}
\bigl([0,1]\times\{0,1\}\times\{0,1\}^2\times\{0,1\}\bigr)
\]
satisfying all of the following conditions:
\begin{itemize}
\item \textbf{(Consistency.)} \cref{obj:ass:consistency}.
\item \textbf{(Exchangeability.)} \cref{obj:ass:exchangeability}.
\item \textbf{(Covariate density.)} \cref{obj:ass:density}.
\item \textbf{(Overlap.)} \cref{obj:ass:overlap}.
\item \textbf{(Propensity smoothness.)} \cref{obj:ass:propensity-holder}.
\item \textbf{(Control-mean smoothness.)} \cref{obj:ass:control-holder}.
\item \textbf{(Contrast smoothness.)} \cref{obj:ass:contrast-lipschitz}.
\end{itemize}
The function conditions apply to the selected versions of the observational regressions. The covariate density is an unknown measurable function. Both arm means range over their full binary-outcome domain, subject to the stated conditions. Membership in \(\mathcal P\) is determined by these properties of the causal law.
\end{definitionv}

For laws in \cref{obj:def:complete-model}, consistency, exchangeability, and overlap identify the observational contrast with the conditional mean of \(Y(1)-Y(0)\) almost everywhere under the covariate law. Contrast continuity and the positive design density then determine its pointwise interpretation: continuous versions agreeing almost everywhere agree throughout \([0,1]\). Consequently, the target \(\theta(P)\) is the value of the continuous causal contrast at the fixed evaluation location. The argument is given in \cref{obj:lem:point-version}.

\subsection{Private releases and decision criteria}

Privacy constrains a release on the full ordered dataset space in \cref{obj:synth_9}. Its comparison concerns an arbitrary dataset and a comparison dataset \leanref{sym:Dprime}{\ensuremath{D'}}, using replacement Hamming distance \leanref{sym:dHam}{\ensuremath{d_{\mathrm H}(D,D')}}. A randomized release \leanref{sym:M}{\ensuremath{M}} takes values in a standard Borel output space \leanref{sym:B}{\ensuremath{\mathcal B}}; the public privacy budget \leanref{sym:epsilon}{\ensuremath{\epsilon}} bounds the change in its output probabilities.



The privacy requirement applies to every replacement-adjacent pair, including datasets outside any particular sampling law's support.

\begin{assumptionv}{ass:pure-privacy}[Pure replacement differential privacy]
For every pair of datasets \(D,D'\in\mathcal O^n\) with replacement distance \(d_{\mathrm H}(D,D')=1\), and every measurable output event \(E\in\mathcal B(\mathcal B)\),
\[
M(D,E)\le \exp(\epsilon)M(D',E).
\]
\end{assumptionv}

Pure central differential privacy under replacement adjacency bounds the multiplicative change in every output-event probability when one observed record is replaced \citep{AcharyaSunZhang2021}. Requiring this property everywhere makes privacy a condition on the release itself, uniformly across the full dataset space. The admissible class is therefore defined independently of the causal law.

\begin{definitionv}{def:private-kernels}[Private release class \(\mathfrak M_{n,\epsilon}(\mathcal B)\)]
\(\mathfrak M_{n,\epsilon}(\mathcal B)\) is the class of all everywhere-defined Markov kernels
\[
M:\mathcal O^n\rightsquigarrow\mathcal B
\]
satisfying \cref{obj:ass:pure-privacy}. Here \(\mathcal B\) is a standard Borel output space, and \(M(D,E)\) is the conditional probability of an output in the Borel event \(E\) given dataset \(D\).
\end{definitionv}

We use separate decision criteria for scalar estimation and interval inference. For a randomized scalar estimator \leanref{sym:Tgeneric}{\ensuremath{T}}, the minimax absolute-error risk \leanref{sym:R}{\ensuremath{R_{n,\epsilon}}} measures the smallest achievable worst-case expected distance from the pointwise target.

\begin{definitionv}{def:risk-criterion}[Minimax absolute-error risk \(R_{n,\epsilon}\)]
\[
R_{n,\epsilon}
=
\inf_{T\in\mathfrak M_{n,\epsilon}(\mathbb R)}
\sup_{P\in\mathcal P}
E_{(P^{\mathrm{obs}})^{\otimes n},T}
\bigl|T(D)-\theta(P)\bigr|.
\]
Here \(T(D)\) is the scalar output of the randomized estimator, \(\theta(P)\) is the pointwise treatment-effect target, and the expectation integrates both the observed sample with law \((P^{\mathrm{obs}})^{\otimes n}\) and the estimator randomization.
\end{definitionv}

The supremum evaluates accuracy uniformly over the causal class in \cref{obj:def:complete-model}, while the infimum ranges over all private scalar releases in \cref{obj:def:private-kernels}. Interval inference instead evaluates expected length subject to a uniform coverage requirement. To express that criterion for randomized procedures, we first specify the measurable interval output space \leanref{sym:Ispace}{\ensuremath{\mathcal I}}.

\begin{definitionv}{synth_13}[Measurable interval output space]
Define
\[
\mathcal I=\{I\subseteq\mathbb R:I\text{ is a connected Borel set}\}.
\]
Code the empty set by a distinguished symbol \(\varnothing\). Code each nonempty interval by its lower and upper endpoints \(\ell,u\) and endpoint-inclusion flags \(\eta_\ell,\eta_u\in\{0,1\}\), where
\[
\ell\in\mathbb R\cup\{-\infty\},\qquad
u\in\mathbb R\cup\{+\infty\},\qquad \ell\le u.
\]
For \(\ell<u\), the code represents
\[
(\ell,u)\cap\mathbb R
\;\cup\;
\begin{cases}\{\ell\},&\ell\in\mathbb R,\ \eta_\ell=1,\\ \varnothing,&\text{otherwise},\end{cases}
\;\cup\;
\begin{cases}\{u\},&u\in\mathbb R,\ \eta_u=1,\\ \varnothing,&\text{otherwise}.\end{cases}
\]
An infinite endpoint has inclusion flag zero. For \(\ell=u\), require \(\ell=u\in\mathbb R\) and \(\eta_\ell=\eta_u=1\), so the code represents the singleton \(\{\ell\}\). Equip the valid nonempty codes with the Borel sigma-algebra inherited from
\[
(\mathbb R\cup\{-\infty\})\times
(\mathbb R\cup\{+\infty\})\times\{0,1\}^2,
\]
using the usual extended-real Borel structures and discrete flags, and adjoin the empty-set code as an isolated point. Transport this sigma-algebra to \(\mathcal I\) through the coding bijection. This makes \(\mathcal I\) a standard Borel measurable space, including empty, singleton, bounded, and unbounded interval outputs.
\end{definitionv}

This coding accommodates all connected Borel intervals and makes coverage events and interval lengths measurable. For a randomized interval procedure \leanref{sym:Igeneric}{\ensuremath{I}}, the honest expected-length criterion \leanref{sym:H}{\ensuremath{H_{n,\epsilon}}} minimizes maximal expected length among private procedures attaining at least \(90\%\) coverage for every law in the class.

\begin{definitionv}{def:interval-criterion}[Honest expected-length criterion \(H_{n,\epsilon}\)]
The minimax expected length under uniform \(90\%\) coverage is
\[
H_{n,\epsilon}
=
\inf_{\substack{
I\in\mathfrak M_{n,\epsilon}(\mathcal I)\\
\inf_{P\in\mathcal P}
\Pr_{(P^{\mathrm{obs}})^{\otimes n},I}
\{\theta(P)\in I(D)\}\ge 9/10
}}
\ \sup_{P\in\mathcal P}
E_{(P^{\mathrm{obs}})^{\otimes n},I}
\operatorname{Leb}(I(D)).
\]
Here \(\mathcal I\) is the interval output space, \(I(D)\) is the released connected interval, and \(\operatorname{Leb}(I(D))\) is its Lebesgue length. The probability and expectation integrate both the observed sample and the procedure's randomization.
\end{definitionv}

Honesty here means uniform coverage over the entire causal class at the given sample size and privacy budget. Connectedness gives length its usual interpretation as the span of a confidence interval, and both coverage and expected length account for sampling and release randomization. Since the target lies in \([-1,1]\), the constant release of that interval is private, has coverage one, and supplies a finite-length admissible procedure. The two criteria in \cref{obj:def:risk-criterion,obj:def:interval-criterion} thus evaluate estimation accuracy and honest interval precision within the same observational experiment and privacy model.


\section{A private release based on within-cell pairs}

A single release of local pair statistics supplies both an estimator for the pointwise treatment effect and an honest confidence interval under the criteria in \cref{obj:def:risk-criterion,obj:def:interval-criterion}. The construction uses public cell geometry and gives the numerator and denominator the same occupancy weights. This shared weighting accommodates every measurable design density satisfying \cref{obj:ass:density}. We first specify the clipping operation used to keep the reported estimate within the target range.
\begin{definitionv}{synth_12}[Clipping to a closed interval]
For real numbers \(a\le b\) and \(u\in\mathbb R\), define
\[
\operatorname{clip}_{[a,b]}(u)=\min\{b,\max\{a,u\}\}.
\]
\end{definitionv}

Localization uses a public radius \leanref{sym:h}{\ensuremath{h}} around \(x_0=1/2\). A public cell count \leanref{sym:k}{\ensuremath{k}} determines the cell width \leanref{sym:delta}{\ensuremath{\delta}} and the equal-cell partition \leanref{sym:cells}{\ensuremath{(C_j)_{j=1}^k}}. The cell occupancy \leanref{sym:Nj}{\ensuremath{N_j(D)}} counts records with their dataset multiplicities. The following geometry fixes these quantities before the release is computed.
\begin{definitionv}{synth_11}[Public local cells]
Let \(x_0=1/2\), let \(h\in(0,1/4]\) be the public localization radius, and let \(k\ge2\) be the integer public cell count. Set \(\delta=2h/k\). The public local window is \([x_0-h,x_0+h]\subseteq[0,1]\), with indexed partition \((C_j)_{j=1}^k\) defined by
\[
C_j=
\begin{cases}
[x_0-h+(j-1)\delta,\ x_0-h+j\delta),&1\le j<k,\\
[x_0-h+(k-1)\delta,\ x_0+h],&j=k.
\end{cases}
\]
Public cell \(j\) means \(C_j\). A record \((X_i,A_i,Y_i)\) belongs to public cell \(j\) precisely when \(X_i\in C_j\); its count is
\[
N_j(D)=\sum_{i=1}^n\mathbf1\{X_i\in C_j\}.
\]
Records are counted with their dataset multiplicities.
\end{definitionv}

The public occupancy scale \leanref{sym:info}{\ensuremath{\mathcal A}} measures the effective amount of information available from cells containing pairs:
\[
\mathcal A=nh\min\{1,n\delta\}.
\]
The factor \(nh\) describes the local sample size, while \(\min\{1,n\delta\}\) accounts for the availability of a second record in the same cell. Thus the scale accommodates both sparse and densely occupied cells.

For the procedures below, define the public mean-squared-error envelope \leanref{sym:Vbound}{\ensuremath{V(h,\delta)}} and its conservative bias component \leanref{sym:bias}{\ensuremath{b_{\mathrm{env}}}} by
\[
V(h,\delta)=2^{20}\left(
h^2+\delta^{2/5}+\mathcal A^{-1}+(\epsilon\mathcal A)^{-2}
\right),
\qquad
b_{\mathrm{env}}=2^{10}\sqrt{h^2+\delta^{2/5}}.
\]
Here \(b_{\mathrm{env}}^2\) collects the localization and within-cell bias terms in the envelope. The appendix separately writes \(b_{\mathrm{cell}}=3h+24\delta^{1/5}\) for the sharper cellwise population certificate. The numerical benchmark \leanref{sym:r}{\ensuremath{r(n,\epsilon)}}, which will determine the public tuning, is
\[
r(n,\epsilon)=
\min\left\{1,\,
\max\left\{
n^{-1/4},\,
(n^2\epsilon)^{-1/7},\,
(n\epsilon)^{-1/2}
\right\}\right\}.
\]
All of these quantities depend on public parameters.

The numerator statistic \leanref{sym:SN}{\ensuremath{S_N(D)}} aggregates within-cell treatment--outcome differences, and the denominator statistic \leanref{sym:SD}{\ensuremath{S_D(D)}} aggregates within-cell treatment differences. Their common occupancy weighting explains the construction's treatment of the design density. To describe that weighting, let \((X,A,Y)\) and \((X',A',Y')\) be independent records drawn from the observed law conditional on membership in cell \(C_j\). Define the conditional pair moments \leanref{sym:nu}{\ensuremath{\nu_j}} and \leanref{sym:dj}{\ensuremath{d_j}}, and the occupancy weight \leanref{sym:w}{\ensuremath{w_j}}, by
\[
\nu_j=E\{(A-A')(Y-Y')\mid X,X'\in C_j\},
\qquad
d_j=E\{(A-A')^2\mid X,X'\in C_j\},
\qquad
w_j=E\{N_j(D)\mathbf1\{N_j(D)\ge2\}\}.
\]
The population means \leanref{sym:Nbar}{\ensuremath{\overline N}} and \leanref{sym:Dbar}{\ensuremath{\overline D}} of the two statistics consequently take the form
\[
\overline N=\sum_{j=1}^k w_j\nu_j,
\qquad
\overline D=\sum_{j=1}^k w_jd_j.
\]
Both sums use the same weights, including the effect of cells containing fewer than two records. Their ratio therefore averages the cell ratios with weights proportional to \(w_jd_j\). The smoothness restrictions in \cref{obj:ass:propensity-holder,obj:ass:control-holder,obj:ass:contrast-lipschitz} control the cellwise approximation to the target, while the measurable density determines the common weighting. The detailed occupancy argument is given in \cref{obj:lem:occupancy-cancellation}.

The first procedure computes a single joint noisy release \leanref{sym:release}{\ensuremath{\widetilde S}}. A positive public denominator floor \leanref{sym:d0}{\ensuremath{d_0}} then stabilizes the ratio. The local estimator \leanref{sym:Thk}{\ensuremath{\widehat T_{h,k}}} and local interval \leanref{sym:Ihk}{\ensuremath{\widehat I_{h,k}}} are both postprocessings of that release.
\begin{algorithmv}{def:private-ratio}[Private ratio \(\widehat T_{h,k}\) and interval \(\widehat I_{h,k}\)]
The inputs are an ordered dataset \(D\in\mathcal O^n\), the public localization radius \(h\), the public cell count \(k\), the privacy parameter \(\epsilon\), the public occupancy scale \(\mathcal A\), and the public error envelope \(V\).
\begin{enumerate}
\item Set the cell width and positive denominator floor to
\[
\delta=\frac{2h}{k},
\qquad
d_0=\frac{3\mathcal A}{64}>0.
\]
For public cell \(j\), let its records and record count be
\[
D_j=\{(X_i,A_i,Y_i):i\text{ indexes a record of }D
\text{ in public cell }j\},
\qquad
s_j=|D_j|,
\]
where records are counted with their dataset multiplicities.

\item Form the numerator and denominator statistics by summing the cell contributions:
\[
S_N(D)=
\sum_j
\begin{cases}
0,&s_j<2,\\[2pt]
\displaystyle
\frac{2\left\{
s_j\sum_{(X,A,Y)\in D_j}AY
-\left(\sum_{(X,A,Y)\in D_j}A\right)
 \left(\sum_{(X,A,Y)\in D_j}Y\right)
\right\}}{s_j-1},&s_j\ge2,
\end{cases}
\]
\[
S_D(D)=
\sum_j
\begin{cases}
0,&s_j<2,\\[2pt]
\displaystyle
\frac{2\left\{
s_j\sum_{(X,A,Y)\in D_j}A
-\left(\sum_{(X,A,Y)\in D_j}A\right)^2
\right\}}{s_j-1},&s_j\ge2.
\end{cases}
\]
Only occupied cells need to be included in these sums.

\item Draw independent Laplace noises once:
\[
\xi_N,\xi_D\ \stackrel{\mathrm{ind}}{\sim}\
\operatorname{Laplace}\!\left(0,\frac{12}{\epsilon}\right).
\]
Release the joint statistic
\[
\widetilde S
=
(\widetilde S_N,\widetilde S_D)
=
\bigl(S_N(D)+\xi_N,\ S_D(D)+\xi_D\bigr).
\]

\item Output the following postprocessing of this single release:
\[
\widehat T_{h,k}
=
\operatorname{clip}_{[-1,1]}
\left\{
\frac{\widetilde S_N}
{\max\{d_0,\widetilde S_D\}}
\right\},
\]
\[
\widehat I_{h,k}
=
\left[
\max\left\{-1,\widehat T_{h,k}-\sqrt{10V(h,\delta)}\right\},
\min\left\{1,\widehat T_{h,k}+\sqrt{10V(h,\delta)}\right\}
\right].
\]
The map is defined on every \(D\in\mathcal O^n\).
\end{enumerate}
\end{algorithmv}

The cell formulas are count-weighted sums of symmetric pair contributions. They can be computed from cell totals, preserving the pair construction while retaining a direct implementation. Cells with zero or one record contribute zero to both coordinates. Replacement privacy requires control of changes across all cell counts, including replacements that move a record between cells; \cref{obj:lem:all-count-stability} supplies this sensitivity argument. The independent Laplace perturbations implement the joint release, and the estimator and interval inherit its privacy by postprocessing \citep{Dwork2006,Dwork2014}.

Public tuning selects the estimator \leanref{sym:Tstar}{\ensuremath{\widehat T^*}} and interval \leanref{sym:Istar}{\ensuremath{\widehat I^*}} using the benchmark alone. In the local branch, the dyadic choice of \(h\) makes \(k\) an integer and sets \(\delta=h^5\), balancing the two squared-bias terms. The following procedure also specifies the constant outputs for the branch \(r(n,\epsilon)\ge1/8\).
\begin{algorithmv}{def:public-tuning}[Public tuning for \(\widehat T^*\) and \(\widehat I^*\)]
The inputs are the dataset, the sample size \(n\), the privacy parameter \(\epsilon\), and the numerical benchmark \(r(n,\epsilon)\).
\begin{enumerate}
\item If \(r(n,\epsilon)\ge1/8\), output
\[
\widehat T^*=0,
\qquad
\widehat I^*=[-1,1].
\]

\item If \(r(n,\epsilon)<1/8\), choose the public tuning parameters
\[
h=2^{-\lceil\log_2(1/r(n,\epsilon))\rceil},
\qquad
k=2h^{-4},
\qquad
\delta=\frac{2h}{k}=h^5,
\]
and output the local private ratio and interval:
\[
(\widehat T^*,\widehat I^*)
=
(\widehat T_{h,k},\widehat I_{h,k}).
\]
\end{enumerate}
This defines a single sample map whose public tuning depends only on \(n\) and \(\epsilon\).
\end{algorithmv}

The stated interval is deliberately conservative at moderate sample sizes. In the local branch, \(\delta=h^5\), so \(V(h,\delta)\ge 2^{21}h^2\). The dyadic rule gives \(h>r(n,\epsilon)/2\), while \(r(n,\epsilon)\ge n^{-1/4}\); hence
\[
10V(h,\delta)>5\cdot 2^{20}n^{-1/2}.
\]
For \(2\le n\le 2^{40}\), its square root exceeds the diameter of \([-1,1]\), and the clipped centered interval is therefore exactly \([-1,1]\), with the theorem's uniform coverage guarantee. The fallback branch has the same output. The practical interpretation of these constants is discussed under Limitations and future work.

The interval also has a candidate-value interpretation. For each candidate effect \leanref{sym:vartheta}{\ensuremath{\vartheta}}, inversion compares a clipped numerator with the candidate multiplied by the stabilized denominator. The resulting interval \leanref{sym:Iopt}{\ensuremath{I^*_{n,\epsilon}}} uses the same joint release as the point estimate. To express the connected output of the acceptance rule, we first define the interval hull.
\begin{definitionv}{synth_17}[Interval hull]
For a nonempty subset \(S\subseteq\mathbb R\), define its interval hull as the smallest connected interval containing \(S\):
\[
\operatorname{hull}(S)
=
\{x\in\mathbb R:\text{there exist }a,b\in S\text{ with }a\le x\le b\}.
\]
\end{definitionv}

With this operation, the inversion procedure expresses the interval directly in terms of the released coordinates and public error envelope. In the inversion algorithm, the symbol \(b\) denotes \(b_{\mathrm{env}}\).
\begin{algorithmv}{def:interval-handle}[Interval inversion $I^*_{n,\epsilon}$]
The inputs are the sample size \(n\), privacy parameter \(\epsilon\), benchmark \(r(n,\epsilon)\), public tuning parameters \(h,k,\delta\), denominator floor \(d_0\), and mean-squared-error envelope \(V(h,\delta)\). The envelope incorporates the bias envelope \(b\), occupancy scale \(\mathcal A\), and Laplace second moment.
\begin{enumerate}
\item If \(r(n,\epsilon)\ge 1/8\), accept every candidate \(\vartheta\in[-1,1]\) and output
\[
I^*_{n,\epsilon}=\widehat I^*=[-1,1].
\]
Otherwise, proceed to the following steps.
\item Compute the single joint noisy release \(\widetilde S\), with numerator coordinate \(\widetilde S_N\) and denominator coordinate \(\widetilde S_D\), using the public tuning parameters. Define
\[
B_D=\max\{d_0,\widetilde S_D\},\qquad
B_N=\operatorname{clip}_{[-B_D,B_D]}(\widetilde S_N),\qquad
\rho=\sqrt{10V(h,\delta)}.
\]
\item Accept a candidate \(\vartheta\in[-1,1]\) precisely when
\[
|B_N-\vartheta B_D|\le \rho B_D.
\]
Output the interval hull of the accepted set, with \([-1,1]\) as the output if the accepted set is empty:
\[
I^*_{n,\epsilon}=\widehat I^*=
\begin{cases}
\operatorname{hull}\{\vartheta\in[-1,1]:|B_N-\vartheta B_D|\le\rho B_D\},
& \{\vartheta\in[-1,1]:|B_N-\vartheta B_D|\le\rho B_D\}\ne\varnothing,\\
[-1,1],
& \{\vartheta\in[-1,1]:|B_N-\vartheta B_D|\le\rho B_D\}=\varnothing.
\end{cases}
\]
\end{enumerate}
\end{algorithmv}

Because \(B_D\) is positive, dividing the acceptance inequality by \(B_D\) gives an interval centered at \(B_N/B_D=\widehat T_{h,k}\), with radius \(\rho\), intersected with \([-1,1]\). The acceptance set contains its center and is connected. Hence the inversion reproduces the interval in \cref{obj:def:private-ratio,obj:def:public-tuning}. Computing the point estimate and this interval together uses one draw of the joint noise.

The next result establishes privacy for every dataset and uniform statistical guarantees over the causal class. It gives both the local squared-error envelope and the absolute-error and expected-length bounds obtained by public tuning.
\begin{theoremv}{thm:uniform-private-upper}[Uniform private upper bounds]*
Let the public parameters satisfy
\begin{itemize}
\item \textbf{(Sample size.)} \(n\in\mathbb N\) and \(n\ge2\).
\item \textbf{(Cell count.)} \(k\in\mathbb N\) and \(k\ge2\).
\item \textbf{(Privacy budget.)} \(0<\epsilon\le1\).
\item \textbf{(Localization radius.)} \(0<h\le1/4\).
\end{itemize}
Define the cell width \(\delta\), public occupancy scale \(\mathcal A\), and public squared-error envelope \(V(h,\delta)\) by
\[
\delta=\frac{2h}{k},
\qquad
\mathcal A=nh\min\{1,n\delta\},
\qquad
V(h,\delta)=2^{20}\left(
h^2+\delta^{2/5}+\mathcal A^{-1}+(\epsilon\mathcal A)^{-2}
\right).
\]
Define the numerical benchmark for public tuning by
\[
r(n,\epsilon)=
\min\left\{1,\,
\max\left\{
n^{-1/4},\,
(n^2\epsilon)^{-1/7},\,
(n\epsilon)^{-1/2}
\right\}\right\}.
\]
The joint release in \cref{obj:def:private-ratio} is
\[
\widetilde S=
\bigl(S_N(D)+\xi_N,\ S_D(D)+\xi_D\bigr),
\qquad
\xi_N,\xi_D\overset{\mathrm{ind}}{\sim}
\operatorname{Laplace}(0,12/\epsilon).
\]
This release and its scalar and interval postprocessings
\(\widehat T_{h,k}\) and \(\widehat I_{h,k}\) in
\cref{obj:def:private-ratio}, together with the publicly tuned outputs
\(\widehat T^*\) and \(\widehat I^*\) in \cref{obj:def:public-tuning},
are everywhere-defined Markov kernels satisfying pure
\(\epsilon\)-replacement privacy on all of \(\mathcal O^n\) in the sense of
\cref{obj:def:private-kernels}.

For every \(P\in\mathcal P\), the causal-law class in
\cref{obj:def:complete-model}, let \(\theta(P)=\tau_P(x_0)\) be the point
target at \(x_0=1/2\). Under the observed sample law
\((P^{\mathrm{obs}})^{\otimes n}\), the local outputs satisfy
\[
E_{(P^{\mathrm{obs}})^{\otimes n},\widehat T_{h,k}}
\left[(\widehat T_{h,k}(D)-\theta(P))^2\right]
\le V(h,\delta),
\]
\[
\Pr_{(P^{\mathrm{obs}})^{\otimes n},\widehat I_{h,k}}
\left\{\theta(P)\in\widehat I_{h,k}(D)\right\}
\ge\frac9{10},
\qquad
E_{(P^{\mathrm{obs}})^{\otimes n},\widehat I_{h,k}}
\operatorname{Leb}(\widehat I_{h,k}(D))
\le\min\left\{2,\ 2\sqrt{10V(h,\delta)}\right\}.
\]
The publicly tuned outputs satisfy
\[
\sup_{P\in\mathcal P}
E_{(P^{\mathrm{obs}})^{\otimes n},\widehat T^*}
\left|\widehat T^*(D)-\theta(P)\right|
\le2^{18}r(n,\epsilon),
\]
\[
\sup_{P\in\mathcal P}
E_{(P^{\mathrm{obs}})^{\otimes n},\widehat I^*}
\operatorname{Leb}(\widehat I^*(D))
\le2^{22}r(n,\epsilon),
\]
and, for every \(P\in\mathcal P\),
\[
\Pr_{(P^{\mathrm{obs}})^{\otimes n},\widehat I^*}
\left\{\theta(P)\in\widehat I^*(D)\right\}
\ge\frac9{10}.
\]
\end{theoremv}
% CAUSALSMITH-CITED-SCOPE-BEGIN thm:uniform-private-upper
\verificationfootnotetext{\textbf{Formalization scope.} This result uses the cited conclusion \citep{BoucheronLugosiMassart2013} (Theorem 3.1, replacement-copy form, printed page 54). The cited source proof is not formalized here: Lean verifies its use and all remaining steps. If that cited conclusion is false, the portions of this result that depend on it are not certified.}
% CAUSALSMITH-CITED-SCOPE-END thm:uniform-private-upper

The four terms in the error envelope separate the statistical mechanisms. Localization of the Lipschitz contrast contributes \(h^2\), and within-cell approximation of the two rough nuisance functions contributes \(\delta^{2/5}\). Pair-statistic variation contributes at scale \(\mathcal A^{-1}\). The replacement-copy variance inequality of \citet[Theorem 3.1, replacement-copy form, p.~54]{BoucheronLugosiMassart2013} controls this variation through changes induced by replacing individual independent records. Denominator stabilization keeps the noisy ratio defined across all occupancies, while the Laplace second moments contribute the term \((\epsilon\mathcal A)^{-2}\). The occupancy and replacement-sensitivity details appear in \cref{obj:lem:occupancy-cancellation,obj:lem:all-count-stability}.

Under the public choice \(\delta=h^5\), the two squared-bias terms have the same order. The benchmark balances this bias against pair variation and privacy noise in the sparse and dense occupancy regimes. Finally, the squared-error bound supplies the \(90\%\) coverage guarantee through the radius \(\sqrt{10V(h,\delta)}\). Intersecting the interval with the target range preserves coverage and bounds its length by two. Thus \cref{obj:thm:uniform-private-upper} supplies a common private release with uniform estimation and finite-sample interval guarantees throughout the causal class in \cref{obj:def:complete-model}.


\section{Optimal estimation and honest confidence intervals}

The uniform guarantees in \cref{obj:thm:uniform-private-upper} lead to two decision-theoretic conclusions. The first characterizes the smallest maximal expected absolute error among all pure-private scalar procedures. The second characterizes the smallest maximal expected length among pure-private connected intervals with class-wide \(90\%\) coverage. Both criteria are governed by the same numerical benchmark, under the causal-law restrictions in \cref{obj:def:complete-model}.

For scalar estimation, the comparison includes every everywhere-defined private randomized procedure admitted by \cref{obj:def:private-kernels}. The publicly tuned estimator of \cref{obj:def:public-tuning} attains the optimal order uniformly over the complete model class, including its admissible measurable covariate densities.

\begin{theoremv}{thm:matched-risk-frontier}[Matched private risk frontier]*
For every integer \(n\ge2\) and privacy parameter \(0<\epsilon\le1\), define the numerical benchmark
\[
r(n,\epsilon)
=
\min\left\{1,\,
\max\left\{n^{-1/4},\,(n^2\epsilon)^{-1/7},\,(n\epsilon)^{-1/2}\right\}\right\}.
\]
The private minimax absolute risk \(R_{n,\epsilon}\) of \cref{obj:def:risk-criterion}, with the infimum taken over all everywhere-defined randomized pure-private scalar Markov kernels as in \cref{obj:def:private-kernels}, satisfies
\[
2^{-20}r(n,\epsilon)
\le R_{n,\epsilon}
\le 2^{18}r(n,\epsilon).
\]
Moreover, the publicly tuned scalar estimator \(\widehat T^*\), which releases zero when \(r(n,\epsilon)\ge1/8\) and otherwise uses the specified publicly tuned dyadic pair-ratio release, is an everywhere-defined Markov kernel satisfying pure replacement \(\epsilon\)-differential privacy. Its maximal expected absolute error satisfies
\[
\sup_{P\in\mathcal P}
E_{(P^{\mathrm{obs}})^{\otimes n},\widehat T^*}
\left|\widehat T^*(D)-\theta(P)\right|
\le 2^{18}r(n,\epsilon),
\]
where \(\mathcal P\) is the complete causal-law class and \(\theta(P)\) is its pointwise treatment-effect target, with expectation integrating the observed sample and release randomization as in \cref{obj:def:risk-criterion}. The supremum includes every admissible measurable covariate density \(f_P\) in that class.
\end{theoremv}
% CAUSALSMITH-CITED-SCOPE-BEGIN thm:matched-risk-frontier
\verificationfootnotetext{\textbf{Formalization scope.} This result uses the cited conclusion \citep{BoucheronLugosiMassart2013} (Theorem 3.1, replacement-copy form, printed page 54), \citep{BoucheronLugosiMassart2013} (Theorem 6.2, displayed log-mgf conclusion in its proof, printed page 167) and \citep{KellerTrotterAppliedCombinatorics} (Section 5.6, Theorem 5.40). The cited source proof is not formalized here: Lean verifies its use and all remaining steps. If that cited conclusion is false, the portions of this result that depend on it are not certified.}
% CAUSALSMITH-CITED-SCOPE-END thm:matched-risk-frontier

The benchmark expresses three competing accuracy scales. The term \(n^{-1/4}\) reflects the statistical balance between localization bias and variation of the within-cell pairs. The two budget-dependent terms reflect the interaction of privacy noise with sparse and dense cell occupancy. Public tuning balances these contributions, and the outer truncation accommodates the bounded target range. The lower bound establishes that their maximum is an unavoidable order of maximal absolute risk over the stated class. The constants are explicit finite-sample bounds; optimality here concerns the common order in sample size and privacy budget.

The lower-bound intuition is observational indistinguishability. Separated treatment-effect targets can be embedded in causal laws whose observed samples are sufficiently difficult to distinguish. Any scalar release that estimates both targets accurately would also distinguish the associated experiments. Privacy strengthens this comparison by controlling how much the distribution of a released statistic can change when records are replaced. The testing priors and the private release comparison are presented in \cref{obj:lem:frontier-testing-priors,obj:lem:private-kernel-comparison}; together they connect target separation to a lower bound that applies to every admissible scalar procedure.

The construction of the sign-indexed alternatives and their component couplings appears in \cref{obj:def:cosine-family,obj:def:common-mass-coupling,obj:lem:measurable-component-contraction}, while \cref{obj:def:direct-family} supplies the direct comparison. These constructions organize the observational and privacy comparisons within the same causal model class. The quantitative arguments use the replacement-copy variance inequality of \citet[Theorem 3.1, p.~54]{BoucheronLugosiMassart2013}, the bounded-differences exponential-moment bound in \citet[Theorem 6.2, displayed log-mgf conclusion in the proof, p.~167]{BoucheronLugosiMassart2013}, and the labelled-tree enumeration formula of \citet[Section 5.6, Theorem 5.40]{KellerTrotterAppliedCombinatorics}. The appendix places these tools alongside the experiments to which they apply.

Honest inference imposes a distinct decision criterion: coverage must hold for every law in the class, while expected length is assessed in the worst case, as specified in \cref{obj:def:interval-criterion}. Define the interval benchmark \leanref{sym:ell}{\ensuremath{\ell(n,\epsilon)}} by
\[
\ell(n,\epsilon)=r(n,\epsilon),
\]
and set the lower and upper numerical constants \leanref{sym:cI}{\ensuremath{c_I}} and \leanref{sym:CI}{\ensuremath{C_I}} to
\[
c_I=2^{-20},
\qquad
C_I=2^{22}.
\]
With these conventions, the candidate-value inversion of \cref{obj:def:interval-handle} attains the same order for honest expected length.

The theorem below has two branches: the full-range interval applies when \(r(n,\epsilon)\ge1/8\), while every subsequent display involving \(h\), \(k\), \(\delta\), \(\mathcal A\), \(V\), and the centered pair-ratio interval applies only on the local branch \(r(n,\epsilon)<1/8\).

\begin{theoremv}{thm:sharp-interval-frontier}[Sharp honest interval frontier]*
For every integer \(n\ge2\) and privacy budget \(0<\epsilon\le1\), define the numerical benchmark
\[
r(n,\epsilon)
=
\min\left\{1,\,
\max\left\{
n^{-1/4},
(n^2\epsilon)^{-1/7},
(n\epsilon)^{-1/2}
\right\}\right\}.
\]
The honest expected-length benchmark and its constants satisfy
\[
\ell(n,\epsilon)=r(n,\epsilon)>0,
\qquad
c_I=2^{-20}>0,
\qquad
C_I=2^{22}>0.
\]

The interval kernel \(I^*_{n,\epsilon}\) defined by candidate-value inversion in \cref{obj:def:interval-handle} equals the total publicly tuned connected-interval kernel of \cref{obj:def:public-tuning}:
\[
I^*_{n,\epsilon}=\widehat I^*.
\]
For \(r(n,\epsilon)\ge1/8\), this kernel returns \([-1,1]\). For \(r(n,\epsilon)<1/8\), its public tuning is
\[
h=2^{-\lceil\log_2(1/r(n,\epsilon))\rceil},
\qquad
k=2h^{-4},
\qquad
\delta=\frac{2h}{k}=h^5,
\]
where \(0<h\le1/4\) and \(k\) is an integer with \(k\ge2\). Define the public occupancy scale and mean-squared-error envelope by
\[
\mathcal A=nh\min\{1,n\delta\},
\qquad
V(h,\delta)
=
2^{20}\left[
h^2+\delta^{2/5}+\mathcal A^{-1}
+(\epsilon\mathcal A)^{-2}
\right].
\]
Using the clipped pair-ratio estimator from the single joint release in \cref{obj:def:private-ratio}, the tuned interval is
\[
\widehat I^*(D)
=
[-1,1]\cap
\left[
\widehat T_{h,k}(D)-\sqrt{10V(h,\delta)},
\widehat T_{h,k}(D)+\sqrt{10V(h,\delta)}
\right].
\]

The kernel \(I^*_{n,\epsilon}\) is a Markov kernel satisfying pure replacement \(\epsilon\)-differential privacy in the sense of \cref{obj:def:private-kernels}. For the causal-law class \(\mathcal P\) in \cref{obj:def:complete-model},
\[
\Pr_{(P^{\mathrm{obs}})^{\otimes n},I^*_{n,\epsilon}}
\left\{\theta(P)\in I^*_{n,\epsilon}(D)\right\}
\ge\frac9{10}
\qquad\text{for every }P\in\mathcal P.
\]
The honest private expected-length frontier \(H_{n,\epsilon}\) of \cref{obj:def:interval-criterion}, whose infimum ranges over all total private Markov connected-interval kernels with this class-wide coverage, satisfies
\[
2^{-20}r(n,\epsilon)
\le H_{n,\epsilon}
\le
\sup_{P\in\mathcal P}
E_{(P^{\mathrm{obs}})^{\otimes n},I^*_{n,\epsilon}}
\operatorname{Leb}\!\left(I^*_{n,\epsilon}(D)\right)
\le2^{22}r(n,\epsilon).
\]
For the private minimax absolute-error risk \(R_{n,\epsilon}\) of \cref{obj:def:risk-criterion}, the extended nonnegative real frontiers also satisfy
\[
2^{-38}R_{n,\epsilon}
\le H_{n,\epsilon}
\le2^{42}R_{n,\epsilon}.
\]
\end{theoremv}
% CAUSALSMITH-CITED-SCOPE-BEGIN thm:sharp-interval-frontier
\verificationfootnotetext{\textbf{Formalization scope.} This result uses the cited conclusion \citep{BoucheronLugosiMassart2013} (Theorem 3.1, replacement-copy form, printed page 54), \citep{BoucheronLugosiMassart2013} (Theorem 6.2, displayed log-mgf conclusion in its proof, printed page 167) and \citep{KellerTrotterAppliedCombinatorics} (Section 5.6, Theorem 5.40). The cited source proof is not formalized here: Lean verifies its use and all remaining steps. If that cited conclusion is false, the portions of this result that depend on it are not certified.}
% CAUSALSMITH-CITED-SCOPE-END thm:sharp-interval-frontier

The upper bound turns the uniform mean-squared-error envelope into a public interval radius. The factor \(\sqrt{10}\) yields the stated coverage level, and intersection with the target range preserves containment of the true effect. Both endpoints are functions of the single joint private release and public quantities. Thus the estimator and interval share a release, while the interval guarantee controls coverage separately at every member of the causal class. The expected-length bound then places this particular procedure within a fixed numerical factor of the optimal honest length.

The interval converse has its own geometric intuition. Honesty requires containment of each target under its corresponding experiment. Observational indistinguishability and privacy contraction transfer these containment probabilities to a common output distribution. With positive probability, the released interval must consequently contain both separated targets. Connectedness forces its length on that event to be at least their separation. This simultaneous-containment argument links the experiments in \cref{obj:lem:frontier-testing-priors} directly to the expected-length criterion. It explains why honest connected intervals inherit the same numerical order as scalar estimation, with the explicit comparison between the two criteria given in \cref{obj:thm:sharp-interval-frontier}.

The benchmark also provides a direct guide to sample-size and budget choices. Its transition points identify which accuracy scale is largest, and its behavior along budget sequences determines uniform consistency.

\begin{propositionv}{prop:regime-and-sanity}[Benchmark regimes and consistency]*
For every integer \(n\ge2\) and every privacy budget \(0<\epsilon\le1\), define the numerical benchmark
\[
r(n,\epsilon)
=
\min\left\{1,\,
\max\left\{
n^{-1/4},\,
(n^2\epsilon)^{-1/7},\,
(n\epsilon)^{-1/2}
\right\}\right\}.
\]
Its regime identities are
\[
\begin{aligned}
n^{-1/4}\le\epsilon
&\ \Longrightarrow\ r(n,\epsilon)=n^{-1/4},\\
n^{-3/5}\le\epsilon\le n^{-1/4}
&\ \Longrightarrow\ r(n,\epsilon)=(n^2\epsilon)^{-1/7},\\
n^{-1}\le\epsilon\le n^{-3/5}
&\ \Longrightarrow\ r(n,\epsilon)=(n\epsilon)^{-1/2},\\
\epsilon\le n^{-1}
&\ \Longrightarrow\ r(n,\epsilon)=1.
\end{aligned}
\]
In particular,
\[
r(n,1)=n^{-1/4}.
\]

For every budget sequence \((\epsilon_n)_{n\in\mathbb N}\) satisfying \(0<\epsilon_n\le1\) for every \(n\ge2\), the private minimax absolute risk \(R_{n,\epsilon_n}\) defined in \cref{obj:def:risk-criterion}, with values in \([0,\infty]\), satisfies
\[
R_{n,\epsilon_n}\longrightarrow0
\quad\Longleftrightarrow\quad
n\epsilon_n\longrightarrow+\infty
\qquad(n\to\infty).
\]
\end{propositionv}
% CAUSALSMITH-CITED-SCOPE-BEGIN prop:regime-and-sanity
\verificationfootnotetext{\textbf{Formalization scope.} This result uses the cited conclusion \citep{BoucheronLugosiMassart2013} (Theorem 3.1, replacement-copy form, printed page 54), \citep{BoucheronLugosiMassart2013} (Theorem 6.2, displayed log-mgf conclusion in its proof, printed page 167) and \citep{KellerTrotterAppliedCombinatorics} (Section 5.6, Theorem 5.40). The cited source proof is not formalized here: Lean verifies its use and all remaining steps. If that cited conclusion is false, the portions of this result that depend on it are not certified.}
% CAUSALSMITH-CITED-SCOPE-END prop:regime-and-sanity

For each fixed sample size, \cref{obj:prop:regime-and-sanity} gives the following budget regimes. The expressions agree at every shared endpoint.
\[
\begin{array}{lll}
\hline
\text{Privacy budget} & \text{Benchmark} & \text{Dominant scale}\\
\hline
n^{-1/4}\le\epsilon\le1
& n^{-1/4}
& \text{Statistical variation}\\
n^{-3/5}\le\epsilon\le n^{-1/4}
& (n^2\epsilon)^{-1/7}
& \text{Sparse-pair privacy noise}\\
n^{-1}\le\epsilon\le n^{-3/5}
& (n\epsilon)^{-1/2}
& \text{Dense-occupancy privacy noise}\\
0<\epsilon\le n^{-1}
& 1
& \text{Bounded-range saturation}\\
\hline
\end{array}
\]

The first regime includes the budget \(\epsilon=1\), at which the benchmark is \(n^{-1/4}\). As the budget decreases, the two privacy terms successively determine the optimal order. In the saturation regime, the lower bound in \cref{obj:thm:matched-risk-frontier} keeps maximal absolute risk bounded away from zero, and the interval lower bound in \cref{obj:thm:sharp-interval-frontier} gives the corresponding positive expected-length bound.

For a sequence of admissible budgets, the decisive consistency condition is \(n\epsilon_n\to+\infty\). By \cref{obj:prop:regime-and-sanity}, this condition is necessary and sufficient for the private minimax absolute risk to vanish. The comparison in \cref{obj:thm:sharp-interval-frontier} yields the same criterion for vanishing optimal honest expected length. Consequently, the two principal results characterize both uniform estimation accuracy and the attainable length of honest connected intervals through a common sample-size and privacy-budget scale.


\section{Discussion and limitations}

For the binary observational class in \cref{obj:def:complete-model}, \cref{obj:thm:matched-risk-frontier,obj:thm:sharp-interval-frontier} give the common order \(r(n,\epsilon)\) for private point estimation and honest expected interval length. The lower bounds range over every procedure admitted by the respective decision criteria, while the pair-ratio release supplies attainment.

The shared occupancy weighting is the main reusable feature of the upper-bound construction. It gives the population numerator and denominator common density-dependent occupancy weights, while the occupancy scale \(\mathcal A=nh\min\{1,n\delta\}\) measures the availability of usable pairs in localized cells. Nuisance roughness determines the required cell resolution, and privacy noise is then amplified by the same occupancy shortage. This construction shows how density robustness, rough nuisances, and privacy interact through the occupancy scale.

The two private scales have different origins. The term \((n^2\epsilon)^{-1/7}\) reflects privacy in the sparse-pair regime, while \((n\epsilon)^{-1/2}\) reflects privacy once cells are densely occupied. Their transition to saturation yields the consistency condition \(n\epsilon_n\to+\infty\), and the regime comparison characterizes optimal dependence on sample size and privacy budget within the specified model.

Honest uncertainty quantification shares the estimation benchmark through complementary upper and lower arguments. The upper bound converts the uniform squared-error envelope into an interval radius. For the converse, transferred coverage forces a connected interval to contain two separated targets with positive probability under a common release distribution, and connectedness converts that separation into expected length. This simultaneous-containment argument establishes the interval frontier directly. Privacy is uniform over adjacent datasets, while error and coverage are uniform over laws in the causal class and include release randomization.

\subsection{Limitations and future work}

The certified interval constants are extremely conservative. As shown after \cref{obj:def:public-tuning}, the local-branch radius already exceeds two for every \(2\le n\le2^{40}\), so the reported interval is deterministically \([-1,1]\) throughout that range; the fallback branch is also full-range. The interval frontier should therefore be read as an asymptotic order result and a finite-sample coverage certificate, rather than as a claim of useful finite-sample width at these constants.

The regime identities characterize rates and do not by themselves prescribe finite-sample sample-size or privacy-budget choices.

The release is specified using real-valued arithmetic and ideal Laplace random variables. A represented-real implementation requires explicit treatment of numerical rounding, noise generation, denominator stabilization, and interval endpoints. Software privacy certification is a separate task: it must establish the privacy guarantee for the implemented output distribution and account for the numerical operations used by the program.

The minimax comparisons range over pure-private procedures with budgets in \((0,1]\). Establishing the unrestricted nonprivate minimax risk is a separate comparison. In particular, the benchmark at \(\epsilon=1\) describes the stated private decision problem; identifying its relation to an unrestricted nonprivate optimum requires an appropriate converse and matching construction for that larger procedure class.

The model uses binary treatment and outcomes, a scalar covariate, independent records, and fixed density, overlap, and smoothness restrictions. Broader outcome distributions, higher-dimensional covariates, and dependent or alternative sampling designs call for corresponding identification conditions, sensitivity bounds, and lower-bound experiments. Allowing overlap to vary with sample size or changing the nuisance and contrast smoothness classes would likewise require a new analysis of the bias, occupancy, and privacy balance. Adaptation to unknown smoothness introduces the additional question of selecting resolution while preserving privacy and honest coverage.

Finally, the results evaluate worst-case decision criteria and supply a concrete procedure attaining their optimal orders. Evaluation of particular causal-effect estimators remains a useful direction, including comparisons of finite-sample error, coverage, interval length, and computational cost under specified data-generating laws. Such evaluations would clarify how the explicit construction and its conservative error envelope perform at sample sizes and privacy budgets relevant to empirical applications.


\appendix

\section*{Appendices}

\section{Identification and upper-bound proofs}

The auxiliary results establish the pointwise interpretation of the target, the population bias of the pair ratio, and the stability of its two coordinates. Their statistical roles are complementary: identification fixes the estimand, common occupancy weights preserve the cellwise approximation, and replacement bounds support both privacy and variance control. Together they supply the ingredients used in \cref{obj:thm:uniform-private-upper}.

\subsection{Point identification}

The selected contrast satisfies the Lipschitz restriction in \cref{obj:ass:contrast-lipschitz}. Its continuity, together with the positive density in \cref{obj:ass:density}, determines a unique continuous version throughout the covariate domain. The following result connects this version to the potential-outcome contrast and establishes that its value at the evaluation point is determined by the observed-record law.

\begin{lemmav}{lem:point-version}[Continuous contrast and point identification]
Let \(P\in\mathcal P\), the class of causal laws specified in \cref{obj:def:complete-model}. Write \(P_X\) for the law of the covariate \(X\), and \(P^{\mathrm{obs}}\) for the law of the observed record \((X,A,Y)\). The selected observational contrast \(\tau_P\) is the difference of the selected treated and control outcome regressions:
\[
\tau_P(x)
=
E_P[Y\mid X=x,A=1]
-
E_P[Y\mid X=x,A=0],
\qquad x\in[0,1].
\]
Here the conditional means denote the selected regression versions carried by \(P\). Define the point target by
\[
x_0=\tfrac12,
\qquad
\theta(P)=\tau_P(x_0).
\]

Then \(\tau_P\) is continuous on \([0,1]\). For every continuous function \(v:[0,1]\to\mathbb R\),
\[
v=\tau_P\quad P_X\text{-almost everywhere}
\quad\Longrightarrow\quad
v(x)=\tau_P(x)\quad\text{for every }x\in[0,1].
\]
Moreover, the point target and the potential-outcome contrast satisfy
\[
\theta(P)\in[-1,1],
\qquad
\tau_P(X)=E_P[Y(1)-Y(0)\mid X]
\quad P\text{-almost surely},
\]
where \(Y(0)\) and \(Y(1)\) are the binary control and treated potential outcomes. Finally, for every \(P'\in\mathcal P\),
\[
P'^{\mathrm{obs}}=P^{\mathrm{obs}}
\quad\Longrightarrow\quad
\theta(P')=\theta(P).
\]
\end{lemmav}

\begin{proof}[Proof of \cref{obj:lem:point-version}]
\begin{enumerate}
\item For the selected arm regressions, write
\[
\mu_{a,P}(x)=E_P[Y\mid X=x,A=a],
\qquad a\in\{0,1\},\quad x\in[0,1].
\]
By \cref{obj:ass:contrast-lipschitz}, with \(L\) the common smoothness radius,
\[
|\tau_P(x)-\tau_P(x')|\le L|x-x'|,
\qquad x,x'\in[0,1].
\]
Thus \(\tau_P\) is Lipschitz and hence continuous. % lean: CausalSmith.Stat.PrivateCateRoughdesign.continuous_tau

\item For a Borel set \(B_{\mathrm{cov}}\subseteq[0,1]\), the density representation and \cref{obj:ass:density} give
\[
P_X(B_{\mathrm{cov}})
=\int_{B_{\mathrm{cov}}}f_P(x)\,dx,
\qquad
f_P(x)\ge\tfrac12
\quad\text{for Lebesgue-almost every }x.
\]
The density is therefore positive almost everywhere, so
\[
P_X(B_{\mathrm{cov}})=0
\quad\Longrightarrow\quad
\int_{B_{\mathrm{cov}}}dx=0.
\]
Consequently, equality \(P_X\)-almost everywhere implies equality Lebesgue-almost everywhere.

We record the continuous-function argument used here and below. Define the continuous projection
\[
c_{\mathrm{proj}}:\mathbb R\longrightarrow[0,1],
\qquad
c_{\mathrm{proj}}(u)=\max\{0,\min\{1,u\}\}.
\]
Let
\[
v_1,v_2:[0,1]\longrightarrow\mathbb R
\]
be continuous functions agreeing Lebesgue-almost everywhere. Their continuous extensions \(v_1\circ c_{\mathrm{proj}}\) and \(v_2\circ c_{\mathrm{proj}}\) agree almost everywhere on \([0,1]\), because
\[
c_{\mathrm{proj}}(x)=x,\qquad x\in[0,1].
\]
If their values differed at a point of \([0,1]\), continuity would give a relative neighborhood of that point on which they differed. Such a neighborhood has positive Lebesgue measure, including when the point is an endpoint. Hence
\[
v_1(c_{\mathrm{proj}}(x))
=v_2(c_{\mathrm{proj}}(x)),
\qquad x\in[0,1],
\]
and therefore \(v_1(x)=v_2(x)\) everywhere.

Now let \(v:[0,1]\to\mathbb R\) be continuous and satisfy \(v=\tau_P\) \(P_X\)-almost everywhere. The density argument transfers this equality to Lebesgue-almost everywhere, and the continuous-function argument yields
\[
v(x)=\tau_P(x),\qquad x\in[0,1].
\] % lean: CausalSmith.Stat.PrivateCateRoughdesign.continuous_tau_unique

\item We next establish the range of the selected arm means. By \cref{obj:ass:control-holder},
\[
|\mu_{0,P}(x)-\mu_{0,P}(x')|
\le L|x-x'|^{1/10},
\qquad x,x'\in[0,1],
\]
so \(\mu_{0,P}\) is continuous. The identity
\[
\mu_{1,P}(x)=\mu_{0,P}(x)+\tau_P(x),
\qquad x\in[0,1],
\]
then makes \(\mu_{1,P}\) continuous as well.

For \(a\in\{0,1\}\), define the arm propensity and the arm-specific covariate measure by
\[
e_{a,P}(x)=
\begin{cases}
1-e_P(x),&a=0,\\
e_P(x),&a=1,
\end{cases}
\qquad x\in[0,1],
\]
\[
\Lambda_{a,P}(B_{\mathrm{cov}})
=P\{X\in B_{\mathrm{cov}},A=a\},
\qquad B_{\mathrm{cov}}\subseteq[0,1]\ \text{Borel}.
\]
The selected propensity identity gives
\[
\int_{B_{\mathrm{cov}}}e_{1,P}(x)\,dP_X(x)
=\Lambda_{1,P}(B_{\mathrm{cov}}).
\]
For the control arm, partitioning \(\{X\in B_{\mathrm{cov}}\}\) by treatment gives
\[
P_X(B_{\mathrm{cov}})
=\Lambda_{0,P}(B_{\mathrm{cov}})
+\Lambda_{1,P}(B_{\mathrm{cov}}),
\]
and hence
\[
\int_{B_{\mathrm{cov}}}e_{0,P}(x)\,dP_X(x)
=P_X(B_{\mathrm{cov}})
-\int_{B_{\mathrm{cov}}}e_P(x)\,dP_X(x)
=\Lambda_{0,P}(B_{\mathrm{cov}}).
\]
By \cref{obj:ass:overlap}, both arm propensities satisfy
\[
\tfrac14\le e_{a,P}(x)\le1,
\qquad a\in\{0,1\},\quad x\in[0,1].
\]
Thus
\[
\tfrac14 P_X(B_{\mathrm{cov}})
\le
\int_{B_{\mathrm{cov}}}e_{a,P}(x)\,dP_X(x)
=\Lambda_{a,P}(B_{\mathrm{cov}}).
\]
In particular, a \(\Lambda_{a,P}\)-null Borel set is \(P_X\)-null and, by the preceding density argument, Lebesgue-null.

The selected regression identities and the binary range of \(Y\) give, for every Borel \(B_{\mathrm{cov}}\),
\[
0\le
\int_{B_{\mathrm{cov}}}\mu_{a,P}(x)\,d\Lambda_{a,P}(x)
=
\int_{\{X\in B_{\mathrm{cov}},\,A=a\}}Y\,dP
\le
\Lambda_{a,P}(B_{\mathrm{cov}}).
\]
The selected arm regressions are integrable under their respective arm measures. The set-integral order criterion therefore yields
\[
0\le\mu_{a,P}(x)\le1
\quad\text{for }\Lambda_{a,P}\text{-almost every }x.
\]
Indeed, the displayed inequalities applied to the measurable sets where the regression is negative or exceeds one exclude either violation on a set of positive arm measure. Transferring null sets as above gives the same bound Lebesgue-almost everywhere. Consequently,
\[
\mu_{a,P}(x)=c_{\mathrm{proj}}(\mu_{a,P}(x))
\quad\text{for Lebesgue-almost every }x.
\]
Both sides are continuous, so the continuous-function argument makes this equality pointwise. Since the projection takes values in \([0,1]\),
\[
0\le\mu_{a,P}(x)\le1,
\qquad a\in\{0,1\},\quad x\in[0,1].
\]
Applying these bounds at \(x_0=\tfrac12\) gives
\[
-1\le
\theta(P)=\mu_{1,P}(x_0)-\mu_{0,P}(x_0)
\le1.
\] % lean: CausalSmith.Stat.PrivateCateRoughdesign.theta_mem_Icc

\item We establish the conditional causal identity separately for each arm. All variables used below are measurable and bounded, so their conditional expectations are well defined. The arm-probability identities imply, for every Borel \(B_{\mathrm{cov}}\),
\[
\int_{\{X\in B_{\mathrm{cov}}\}}e_{a,P}(X)\,dP
=
\int_{B_{\mathrm{cov}}}e_{a,P}(x)\,dP_X(x)
=
\int_{\{X\in B_{\mathrm{cov}}\}}\mathbf1\{A=a\}\,dP.
\]
The characterization of conditional expectation by integrals over covariate events therefore gives
\[
E_P[\mathbf1\{A=a\}\mid X]=e_{a,P}(X)
\quad P\text{-almost surely}.
\]

Equality on all Borel covariate sets also identifies the arm measure as
\[
d\Lambda_{a,P}(x)=e_{a,P}(x)\,dP_X(x).
\]
Using this representation and the selected arm-regression identity,
\[
\begin{aligned}
\int_{\{X\in B_{\mathrm{cov}}\}}
e_{a,P}(X)\mu_{a,P}(X)\,dP
&=\int_{B_{\mathrm{cov}}}
e_{a,P}(x)\mu_{a,P}(x)\,dP_X(x)\\
&=\int_{B_{\mathrm{cov}}}\mu_{a,P}(x)\,d\Lambda_{a,P}(x)\\
&=\int_{\{X\in B_{\mathrm{cov}},\,A=a\}}Y\,dP\\
&=\int_{\{X\in B_{\mathrm{cov}}\}}\mathbf1\{A=a\}Y\,dP.
\end{aligned}
\]
Thus
\[
E_P[\mathbf1\{A=a\}Y\mid X]
=e_{a,P}(X)\mu_{a,P}(X)
\quad P\text{-almost surely}.
\]

By \cref{obj:ass:exchangeability}, projecting the potential-outcome pair onto either coordinate gives
\[
Y(a)\perp\!\!\!\perp A\mid X\quad[P],
\qquad a\in\{0,1\}.
\]
On \(\{A=a\}\), \cref{obj:ass:consistency} gives \(Y=Y(a)\) almost surely; on its complement the treatment indicator vanishes. Hence
\[
\mathbf1\{A=a\}Y
=\mathbf1\{A=a\}Y(a)
\quad P\text{-almost surely}.
\]
Conditional independence then factors the conditional expectation:
\[
\begin{aligned}
E_P[\mathbf1\{A=a\}Y\mid X]
&=E_P[\mathbf1\{A=a\}Y(a)\mid X]\\
&=E_P[\mathbf1\{A=a\}\mid X]\,
  E_P[Y(a)\mid X].
\end{aligned}
\]
Combining the preceding identities yields
\[
e_{a,P}(X)\mu_{a,P}(X)
=e_{a,P}(X)E_P[Y(a)\mid X]
\quad P\text{-almost surely}.
\]
Since \(e_{a,P}(X)\ge\tfrac14\), cancellation gives
\[
E_P[Y(a)\mid X]=\mu_{a,P}(X)
\quad P\text{-almost surely}.
\]
Applying this first to the treated arm and then to the control arm, and using conditional-expectation linearity for the bounded potential outcomes, gives
\[
\begin{aligned}
E_P[Y(1)-Y(0)\mid X]
&=E_P[Y(1)\mid X]-E_P[Y(0)\mid X]\\
&=\mu_{1,P}(X)-\mu_{0,P}(X)
=\tau_P(X)
\quad P\text{-almost surely}.
\end{aligned}
\] % lean: CausalSmith.Stat.PrivateCateRoughdesign.tau_conditional_causal_contrast

\item Finally, let \(P'\in\mathcal P\) satisfy \(P'^{\mathrm{obs}}=P^{\mathrm{obs}}\). Use the same arm-measure definition with \(P'\) in place of \(P\). Each arm measure is determined by the observed marginal, so
\[
\Lambda_{a,P'}(B_{\mathrm{cov}})
=P'^{\mathrm{obs}}\{X\in B_{\mathrm{cov}},A=a\}
=P^{\mathrm{obs}}\{X\in B_{\mathrm{cov}},A=a\}
=\Lambda_{a,P}(B_{\mathrm{cov}}).
\]
Thus \(\Lambda_{a,P'}=\Lambda_{a,P}\). The outcome integrals in the selected regression identities are likewise integrals of observed coordinates. For every Borel \(B_{\mathrm{cov}}\),
\[
\begin{aligned}
\int_{B_{\mathrm{cov}}}\mu_{a,P'}(x)\,d\Lambda_{a,P}(x)
&=\int_{\{X\in B_{\mathrm{cov}},\,A=a\}}Y\,dP'^{\mathrm{obs}}\\
&=\int_{\{X\in B_{\mathrm{cov}},\,A=a\}}Y\,dP^{\mathrm{obs}}\\
&=\int_{B_{\mathrm{cov}}}\mu_{a,P}(x)\,d\Lambda_{a,P}(x).
\end{aligned}
\]
Both regressions are integrable under this common arm measure. Uniqueness of integrable functions from their integrals over all measurable sets yields
\[
\mu_{a,P'}(x)=\mu_{a,P}(x)
\quad\text{for }\Lambda_{a,P}\text{-almost every }x.
\]
The null-set transfers established above make this equality Lebesgue-almost everywhere. Both regressions are continuous by their respective model conditions, so the continuous-function argument gives
\[
\mu_{a,P'}(x)=\mu_{a,P}(x),
\qquad a\in\{0,1\},\quad x\in[0,1].
\]
Evaluating the control and treated equalities at \(x_0\) and subtracting gives
\[
\theta(P')
=\mu_{1,P'}(x_0)-\mu_{0,P'}(x_0)
=\mu_{1,P}(x_0)-\mu_{0,P}(x_0)
=\theta(P).
\] % lean: CausalSmith.Stat.PrivateCateRoughdesign.theta_eq_of_observed_eq
\end{enumerate}
\end{proof}

Consistency, conditional exchangeability, and overlap give the observational contrast its causal interpretation under \cref{obj:ass:consistency,obj:ass:exchangeability,obj:ass:overlap}. Positive design density and contrast continuity then give that interpretation a unique continuous representative. Thus \cref{obj:lem:point-version} supplies the pointwise target used in the estimation and coverage criteria, while its range statement justifies clipping estimates and truncating intervals to \([-1,1]\).

\subsection{Population pair moments and occupancy}

Fix the public cells \(C_j\) of \cref{obj:synth_11}. Write the population cell mass as \leanref{sym:p}{\ensuremath{p_j}}, with \(p_j=P_X(C_j)\). For two observed records \(z=(x,a,y)\) and \(z'=(x',a',y')\), the numerator pair contribution \leanref{sym:KN}{\ensuremath{K_N(z,z')}} and denominator pair contribution \leanref{sym:KD}{\ensuremath{K_D(z,z')}} are
\[
K_N(z,z')=(a-a')(y-y'),
\qquad
K_D(z,z')=(a-a')^2.
\]
The moments \(\nu_j\) and \(d_j\) are their expectations under two independent draws from \(P^{\mathrm{obs}}\) conditional on membership in \(C_j\). Equivalently,
\[
\nu_j=E\{K_N(Z,Z')\mid X,X'\in C_j\},
\qquad
d_j=E\{K_D(Z,Z')\mid X,X'\in C_j\}.
\]
The density restriction ensures that each cell has positive probability, so these conditional laws are well defined.

The count weighting in \cref{obj:def:private-ratio} assigns a cell with \(s\ge2\) records \(s\) times its average unordered-pair contribution. Its population occupancy weight and the \leanref{sym:Wocc}{\ensuremath{W(P)}}, total population occupancy weight, are therefore specified by
\[
w_j
=
E_Q\{N_j(D)\mathbf1\{N_j(D)\ge2\}\}
=
np_j\bigl[1-(1-p_j)^{n-1}\bigr],
\qquad
W(P)=\sum_{j=1}^k w_j,
\]
where \(Q=(P^{\mathrm{obs}})^{\otimes n}\) is the dataset law under \cref{obj:ass:iid}. The population means are
\[
\overline N=\int S_N(D)\,dQ(D),
\qquad
\overline D=\int S_D(D)\,dQ(D).
\]
These quantities retain the within-cell conditional distribution and the random occupancy separately.

For the cellwise certificate below, write \(b_{\mathrm{cell}}=3h+24\delta^{1/5}\). This bound describes the population approximation; the public mean-squared-error envelope in \cref{obj:thm:uniform-private-upper} uses the distinct conservative quantity \(b_{\mathrm{env}}=2^{10}\sqrt{h^2+\delta^{2/5}}\). In the occupancy lemma, its symbol \(b\) denotes \(b_{\mathrm{cell}}\). The next result establishes the shared occupancy representation and its denominator and bias bounds.

\begin{lemmav}{lem:occupancy-cancellation}[Occupancy cancellation and ratio bias]
Suppose the following conditions hold:
\begin{itemize}
\item \textbf{(Sample and cells.)} The sample size \(n\) and public cell count \(k\) are integers satisfying \(n\ge 2\) and \(k\ge 2\).
\item \textbf{(Localization.)} The public localization radius \(h\) satisfies \(0<h\le 1/4\).
\item \textbf{(Model.)} The causal law \(P\) belongs to the class \(\mathcal P\) of \cref{obj:def:complete-model}.
\item \textbf{(Sampling.)} The dataset law \(Q\) satisfies \cref{obj:ass:iid} for sample size \(n\) and causal law \(P\).
\end{itemize}
Write \(\delta=2h/k\) for the cell width, and define the public bias envelope \(b\) and occupancy scale \(\mathcal A\) by
\[
b=3h+24\delta^{1/5},
\qquad
\mathcal A=nh\min\{1,n\delta\}.
\]
Let \(\overline N\) and \(\overline D\) be the population means of the numerator and denominator statistics:
\[
\overline N=\int S_N(D)\,dQ(D),
\qquad
\overline D=\int S_D(D)\,dQ(D).
\]
For each public cell \(j=1,\ldots,k\), let \(\nu_j\) and \(d_j\) be the expectations of the numerator and denominator pair contributions, respectively, under two independent observations drawn from the conditional observed law in that cell. Define its occupancy weight by
\[
w_j
=
nP_X(\text{cell }j)
\left[1-\left(1-P_X(\text{cell }j)\right)^{n-1}\right].
\]
With \(W(P)\) denoting the total population occupancy weight, \(\theta(P)=\tau_P(1/2)\) the point target, and \(d_0\) the public denominator floor, the population moments satisfy
\[
\overline N=\sum_{j=1}^{k}w_j\nu_j,
\qquad
\overline D=\sum_{j=1}^{k}w_jd_j.
\]
For every \(j=1,\ldots,k\),
\[
d_j\ge\frac38,
\qquad
\left|\nu_j-\theta(P)d_j\right|\le b\,d_j.
\]
Moreover,
\[
\frac{\mathcal A}{8}\le W(P)\le\frac{9\mathcal A}{2},
\qquad
\overline D\ge d_0,
\qquad
\left|\frac{\overline N}{\overline D}-\theta(P)\right|\le b.
\]
\end{lemmav}

\begin{proof}[Proof of \cref{obj:lem:occupancy-cancellation}]
By \cref{obj:ass:iid}, \(Q=(P^{\mathrm{obs}})^{\otimes n}\). All conditioning below uses this product law.

\begin{enumerate}
\item \textbf{Conditional cell contributions.}
Use the cells \(C_j\) of \cref{obj:synth_11}, and define
\[
p_j=P_X(C_j),\qquad
\mathsf P_j=\frac{1}{p_j}
 P^{\mathrm{obs}}\big|_{\{X\in C_j\}},
\qquad j=1,\ldots,k.
\]
Each cell has Lebesgue length \(\delta\), including the final cell with its right endpoint. Integrating the density bounds in \cref{obj:ass:density} gives
\[
\frac{\delta}{2}\le p_j\le\frac{3\delta}{2}.
\]
Moreover,
\[
0<\delta=\frac{2h}{k}\le h\le\frac14,
\qquad
0<p_j\le\frac38<1.
\]
Thus both cell membership and its complement have positive probability, and \(\mathsf P_j\) is a probability law.

Write the observed records as
\[
D=(O_i)_{i=1}^n,\qquad O_i=(X_i,A_i,Y_i),
\]
and define the pair kernels on two observed records by
\[
\begin{aligned}
K_N\bigl((X,A,Y),(X',A',Y')\bigr)
  &=(A-A')(Y-Y'),\\
K_D\bigl((X,A,Y),(X',A',Y')\bigr)
  &=(A-A')^2.
\end{aligned}
\]
Let \(S_{N,j}(D)\) and \(S_{D,j}(D)\) denote the respective cell summands in \cref{obj:def:private-ratio}. For every label set
\(\mathcal J\subseteq\{1,\ldots,n\}\), define
\[
\mathcal E_{j,\mathcal J}
=
\left\{D:\{i:X_i\in C_j\}=\mathcal J\right\},
\qquad
\mathsf Q_{j,\mathcal J}
=
\frac{Q|_{\mathcal E_{j,\mathcal J}}}
     {Q(\mathcal E_{j,\mathcal J})}.
\]
The product law gives
\[
Q(\mathcal E_{j,\mathcal J})
=p_j^{|\mathcal J|}(1-p_j)^{n-|\mathcal J|}>0.
\]
Indeed, restricting a product measure to this membership rectangle and dividing by its mass gives
\[
\mathsf Q_{j,\mathcal J}
=
\bigotimes_{i=1}^n
\begin{cases}
\mathsf P_j,&i\in\mathcal J,\\[2pt]
\displaystyle
\frac{P^{\mathrm{obs}}|_{\{X\notin C_j\}}}{1-p_j},
&i\notin\mathcal J.
\end{cases}
\]
Consequently, any two distinct coordinates with labels in \(\mathcal J\) have joint law
\(\mathsf P_j\otimes\mathsf P_j\).

The pair representations of the cell summands follow from
\[
\begin{aligned}
\sum_{\substack{i,i'\in\mathcal J\\i<i'}}
 (A_i-A_{i'})(Y_i-Y_{i'})
&=
|\mathcal J|\sum_{i\in\mathcal J}A_iY_i
-\left(\sum_{i\in\mathcal J}A_i\right)
 \left(\sum_{i\in\mathcal J}Y_i\right),\\
\sum_{\substack{i,i'\in\mathcal J\\i<i'}}
 (A_i-A_{i'})^2
&=
|\mathcal J|\sum_{i\in\mathcal J}A_i
-\left(\sum_{i\in\mathcal J}A_i\right)^2.
\end{aligned}
\]
Here the second identity uses \(A_i^2=A_i\). On
\(\mathcal E_{j,\mathcal J}\), therefore,
\[
S_{\bullet,j}(D)=
\begin{cases}
0,&|\mathcal J|<2,\\[3pt]
\displaystyle
\frac{|\mathcal J|}{\binom{|\mathcal J|}{2}}
\sum_{\substack{i,i'\in\mathcal J\\i<i'}}
K_\bullet(O_i,O_{i'}),&|\mathcal J|\ge2,
\end{cases}
\qquad \bullet\in\{N,D\}.
\]
Define the eligible cell count by
\[
s_j(D)=\sum_{i=1}^n\mathbf1\{X_i\in C_j\},
\qquad
s_{j,\ge2}(D)=
\begin{cases}
0,&s_j(D)<2,\\
s_j(D),&s_j(D)\ge2.
\end{cases}
\]
For \(|\mathcal J|\ge2\), the sum has
\(\binom{|\mathcal J|}{2}\) pairs, each with the corresponding moment
\(\nu_j\) or \(d_j\). For the empty and singleton label sets, both contributions and the eligible count are zero. Hence, in every case,
\[
\begin{aligned}
\int S_{N,j}(D)\,d\mathsf Q_{j,\mathcal J}(D)
&=
\left(\int s_{j,\ge2}(D)\,
 d\mathsf Q_{j,\mathcal J}(D)\right)\nu_j,\\
\int S_{D,j}(D)\,d\mathsf Q_{j,\mathcal J}(D)
&=
\left(\int s_{j,\ge2}(D)\,
 d\mathsf Q_{j,\mathcal J}(D)\right)d_j.
\end{aligned}
\]
The membership events form a finite measurable partition. Multiplying these identities by their event probabilities and summing gives
\[
\begin{aligned}
\int S_{N,j}(D)\,dQ(D)
&=\left(\int s_{j,\ge2}(D)\,dQ(D)\right)\nu_j,\\
\int S_{D,j}(D)\,dQ(D)
&=\left(\int s_{j,\ge2}(D)\,dQ(D)\right)d_j.
\end{aligned}
\]
The bounded kernels and the finite sample size ensure integrability throughout.
% lean: CausalSmith.Stat.PrivateCateRoughdesign.conditional_cell_contribution_expectation
% lean: CausalSmith.Stat.PrivateCateRoughdesign.cell_membership_setIntegral_normalize
% lean: CausalSmith.Stat.PrivateCateRoughdesign.integral_cell_membership_partition

\item \textbf{The occupancy weight.}
Pointwise, the eligible cell count counts precisely those cell members with a distinct partner:
\[
s_{j,\ge2}(D)
=
\sum_{i=1}^n
\mathbf1\left\{
X_i\in C_j,\ 
\exists i'\ne i:\ X_{i'}\in C_j
\right\}.
\]
For a fixed \(i\), the event that it belongs to the cell and has a partner is the cell-membership event minus the event that it is the sole member. Independence gives
\[
\begin{aligned}
Q\left\{X_i\in C_j,\ 
 \exists i'\ne i:\ X_{i'}\in C_j\right\}
&=
p_j-p_j(1-p_j)^{n-1}\\
&=p_j\left[1-(1-p_j)^{n-1}\right].
\end{aligned}
\]
Integrating the finite indicator sum therefore yields
\[
\int s_{j,\ge2}(D)\,dQ(D)
=
np_j\left[1-(1-p_j)^{n-1}\right]
=w_j.
\]
% lean: CausalSmith.Stat.PrivateCateRoughdesign.iid_cell_partner_count_expectation

\item \textbf{Population moment identities.}
By the definitions of the statistics,
\[
S_N(D)=\sum_{j=1}^k S_{N,j}(D),
\qquad
S_D(D)=\sum_{j=1}^k S_{D,j}(D).
\]
Linearity of integration, followed by the preceding two steps, gives
\[
\overline N
=\sum_{j=1}^k\int S_{N,j}(D)\,dQ(D)
=\sum_{j=1}^k w_j\nu_j,
\qquad
\overline D
=\sum_{j=1}^k\int S_{D,j}(D)\,dQ(D)
=\sum_{j=1}^k w_jd_j.
\]
% lean: CausalSmith.Stat.PrivateCateRoughdesign.population_cell_sum
% lean: CausalSmith.Stat.PrivateCateRoughdesign.occupancy_cancellation

\item \textbf{The conditional denominator floor.}
For each cell, define its conditional covariate law and mean treatment probability by
\[
\mathsf G_j=\frac{P_X|_{C_j}}{p_j},
\qquad
\overline e_j
=\int A\,d\mathsf P_j
=\frac{1}{p_j}\int_{C_j}e_P(x)\,dP_X(x)
=\int e_P(x)\,d\mathsf G_j(x).
\]
The middle equality is the defining conditional-probability identity for the selected propensity, restricted to the cell and normalized by \(p_j\). Averaging \cref{obj:ass:overlap} gives
\[
\frac14\le\overline e_j\le\frac34.
\]
For two independent binary treatments,
\[
(A-A')^2=A+A'-2AA'.
\]
Integration under \(\mathsf P_j\otimes\mathsf P_j\) thus gives
\[
d_j=2\overline e_j(1-\overline e_j).
\]
Finally,
\[
d_j-\frac38
=
2\left(\overline e_j-\frac14\right)
 \left(\frac34-\overline e_j\right)\ge0.
\]
Therefore \(d_j\ge3/8\).
% lean: CausalSmith.Stat.PrivateCateRoughdesign.conditional_cell_treatment_overlap
% lean: CausalSmith.Stat.PrivateCateRoughdesign.independent_binary_denominator_moment
% lean: CausalSmith.Stat.PrivateCateRoughdesign.conditional_cell_denominator_lower

\item \textbf{The conditional bias certificate.}
First express the observed pair moments through the selected regressions. Restriction to the cell and the defining arm-mean identities give
\[
\begin{aligned}
\int AY\,d\mathsf P_j
&=\int e_P(x)\mu_{1,P}(x)\,d\mathsf G_j(x),\\
\int (1-A)Y\,d\mathsf P_j
&=\int (1-e_P(x))\mu_{0,P}(x)\,d\mathsf G_j(x),\\
\int Y\,d\mathsf P_j
&=\int\left[(1-e_P(x))\mu_{0,P}(x)
             +e_P(x)\mu_{1,P}(x)\right]\,d\mathsf G_j(x).
\end{aligned}
\]
For example, the treated covariate submeasure has density \(e_P\) with respect to \(P_X\); integrating its selected outcome regression over \(C_j\) gives the first identity after division by \(p_j\). The control submeasure has density \(1-e_P\), giving the second, and their sum gives the third.

Expanding the numerator kernel and factoring the independent-record integrals yields
\[
\nu_j
=
2\left[
\int AY\,d\mathsf P_j
-\left(\int A\,d\mathsf P_j\right)
 \left(\int Y\,d\mathsf P_j\right)
\right].
\]
Define, for \(x,x'\in[0,1]\),
\[
\begin{aligned}
\mathcal K_{N,P}(x,x')
&=
e_P(x)(1-e_P(x'))
 \bigl(\mu_{1,P}(x)-\mu_{0,P}(x')\bigr)\\
&\quad+
e_P(x')(1-e_P(x))
 \bigl(\mu_{1,P}(x')-\mu_{0,P}(x)\bigr),\\
\mathcal K_{D,P}(x,x')
&=
e_P(x)(1-e_P(x'))
+e_P(x')(1-e_P(x)).
\end{aligned}
\]
Factoring the four terms in the integral of \(\mathcal K_{N,P}\) gives
\[
\begin{aligned}
&\int\mathcal K_{N,P}\,d(\mathsf G_j\otimes\mathsf G_j)\\
&\quad=
2\left[
\int e_P\mu_{1,P}\,d\mathsf G_j
-\overline e_j
 \left\{
 \int(1-e_P)\mu_{0,P}\,d\mathsf G_j
 +\int e_P\mu_{1,P}\,d\mathsf G_j
 \right\}
\right]
=\nu_j.
\end{aligned}
\]
Likewise, factoring the two denominator terms gives
\[
\int\mathcal K_{D,P}\,d(\mathsf G_j\otimes\mathsf G_j)
=2\overline e_j(1-\overline e_j)=d_j.
\]
% lean: CausalSmith.Stat.PrivateCateRoughdesign.conditional_cell_numerator_design_pair
% lean: CausalSmith.Stat.PrivateCateRoughdesign.conditional_cell_denominator_design_pair

For \(x,x'\in C_j\), the cell endpoints give
\[
|x-x_0|\le h,\qquad
|x'-x_0|\le h,\qquad
|x-x'|\le\delta.
\]
These inequalities include the final cell's right endpoint. The public radius is \(L=3\) by \cref{obj:synth_4}. Hence
\cref{obj:ass:contrast-lipschitz} gives
\[
|\tau_P(x)-\theta(P)|
\le3|x-x_0|\le3h,
\qquad
|\tau_P(x')-\theta(P)|\le3h.
\]
Similarly, \cref{obj:ass:propensity-holder,obj:ass:control-holder} give
\[
\begin{aligned}
&\left|
(e_P(x)-e_P(x'))
(\mu_{0,P}(x)-\mu_{0,P}(x'))
\right|\\
&\quad=
|e_P(x)-e_P(x')|\,
|\mu_{0,P}(x)-\mu_{0,P}(x')|\\
&\quad\le
(3\delta^{1/10})(3\delta^{1/10})
=9\delta^{1/5}.
\end{aligned}
\]
The overlap bounds also imply
\[
\begin{aligned}
\mathcal K_{D,P}(x,x')-\frac38
&=
\left(e_P(x)-\frac14\right)
\left(\frac34-e_P(x')\right)\\
&\quad+
\left(\frac34-e_P(x)\right)
\left(e_P(x')-\frac14\right)
\ge0.
\end{aligned}
\]
Both coefficients \(e_P(x)(1-e_P(x'))\) and
\(e_P(x')(1-e_P(x))\) are nonnegative.

Using \(\mu_{1,P}=\mu_{0,P}+\tau_P\), direct expansion gives
\[
\begin{aligned}
\mathcal K_{N,P}(x,x')
&=
(e_P(x)-e_P(x'))
(\mu_{0,P}(x)-\mu_{0,P}(x'))\\
&\quad+
e_P(x)(1-e_P(x'))\tau_P(x)
+e_P(x')(1-e_P(x))\tau_P(x').
\end{aligned}
\]
Subtracting \(\theta(P)\mathcal K_{D,P}(x,x')\), applying the triangle inequality, and using the preceding estimates yields
\[
\begin{aligned}
&|\mathcal K_{N,P}(x,x')
      -\theta(P)\mathcal K_{D,P}(x,x')|\\
&\quad\le
9\delta^{1/5}
+3h\,e_P(x)(1-e_P(x'))
+3h\,e_P(x')(1-e_P(x))\\
&\quad=
9\delta^{1/5}+3h\,\mathcal K_{D,P}(x,x')\\
&\quad\le
(24\delta^{1/5}+3h)\mathcal K_{D,P}(x,x')
=b\,\mathcal K_{D,P}(x,x').
\end{aligned}
\]
The last inequality uses
\[
9\delta^{1/5}
\le24\delta^{1/5}\mathcal K_{D,P}(x,x'),
\]
which follows from \(\mathcal K_{D,P}(x,x')\ge3/8\).

The probability law \(\mathsf G_j\otimes\mathsf G_j\) is supported on
\(C_j\times C_j\). Therefore,
\[
\begin{aligned}
|\nu_j-\theta(P)d_j|
&=
\left|\int
 \bigl(\mathcal K_{N,P}-\theta(P)\mathcal K_{D,P}\bigr)
 \,d(\mathsf G_j\otimes\mathsf G_j)\right|\\
&\le
\int
 \left|\mathcal K_{N,P}-\theta(P)\mathcal K_{D,P}\right|
 \,d(\mathsf G_j\otimes\mathsf G_j)\\
&\le
b\int\mathcal K_{D,P}\,d(\mathsf G_j\otimes\mathsf G_j)
=bd_j.
\end{aligned}
\]
% lean: CausalSmith.Stat.PrivateCateRoughdesign.cell_cross_treatment_certificate
% lean: CausalSmith.Stat.PrivateCateRoughdesign.conditional_design_pair_bias
% lean: CausalSmith.Stat.PrivateCateRoughdesign.conditional_cell_bias_bound

\item \textbf{Lower partner-probability bound.}
Since \(n\ge2\) and \(p_j\ge\delta/2\),
\[
n-1\ge\frac n2,
\qquad
(n-1)p_j\ge\frac{n\delta}{4}.
\]
The inequality \(1-p_j\le e^{-p_j}\), with \(1-p_j\ge0\), gives
\[
(1-p_j)^{n-1}
\le e^{-(n-1)p_j}
\le e^{-n\delta/4}.
\]
To compare the rational envelope with the truncated occupancy scale, split according to \(n\delta\). If \(n\delta\le1\), then
\[
\frac{\min\{1,n\delta\}}8(4+n\delta)
=\frac{n\delta(4+n\delta)}8
\le n\delta.
\]
If \(n\delta>1\), then
\[
\frac{\min\{1,n\delta\}}8(4+n\delta)
=\frac{4+n\delta}{8}
\le n\delta.
\]
Since \(4+n\delta>0\), both cases yield
\[
\frac{\min\{1,n\delta\}}8
\le\frac{n\delta}{4+n\delta}.
\]
Furthermore,
\[
e^{n\delta/4}\ge1+\frac{n\delta}{4}
\quad\Longrightarrow\quad
e^{-n\delta/4}\le
\frac{1}{1+n\delta/4}.
\]
Combining these estimates in order gives, for every cell,
\[
\begin{aligned}
\frac{\min\{1,n\delta\}}8
&\le\frac{n\delta}{4+n\delta}
=\frac{n\delta/4}{1+n\delta/4}\\
&\le1-e^{-n\delta/4}
\le1-(1-p_j)^{n-1}.
\end{aligned}
\]
% lean: CausalSmith.Stat.PrivateCateRoughdesign.occupancy_partner_lower

\item \textbf{Upper partner-probability bound.}
Bernoulli's inequality and nonnegativity give
\[
(1-p_j)^{n-1}\ge1-(n-1)p_j,
\qquad
(1-p_j)^{n-1}\ge0.
\]
Thus
\[
1-(1-p_j)^{n-1}
\le\min\{1,(n-1)p_j\},
\qquad
(n-1)p_j\le\frac32n\delta.
\]
If \(n\delta\le1\), these inequalities imply
\[
1-(1-p_j)^{n-1}
\le\frac32n\delta
=\frac32\min\{1,n\delta\}.
\]
If \(n\delta>1\), they imply
\[
1-(1-p_j)^{n-1}
\le1\le\frac32
=\frac32\min\{1,n\delta\}.
\]
Hence the same upper envelope applies to every cell.
% lean: CausalSmith.Stat.PrivateCateRoughdesign.occupancy_partner_upper

\item \textbf{Total occupancy bounds.}
Summing the cell-mass bounds and using \(k\delta=2h\) gives
\[
h=\frac{k\delta}{2}
\le\sum_{j=1}^k p_j
\le\frac{3k\delta}{2}=3h.
\]
Define the total occupancy weight by
\[
W(P)=\sum_{j=1}^k w_j.
\]
Multiplication of the lower partner envelope by \(np_j\ge0\), followed by summation, yields
\[
\begin{aligned}
W(P)
&\ge
\sum_{j=1}^k np_j\,
 \frac{\min\{1,n\delta\}}8\\
&=
\frac{n\min\{1,n\delta\}}8\sum_{j=1}^k p_j
\ge\frac{nh\min\{1,n\delta\}}8
=\frac{\mathcal A}{8}.
\end{aligned}
\]
The upper partner envelope similarly yields
\[
\begin{aligned}
W(P)
&\le
\sum_{j=1}^k np_j\,
 \frac32\min\{1,n\delta\}\\
&=
\frac32n\min\{1,n\delta\}\sum_{j=1}^k p_j
\le\frac92nh\min\{1,n\delta\}
=\frac{9\mathcal A}{2}.
\end{aligned}
\]
% lean: CausalSmith.Stat.PrivateCateRoughdesign.occupancy_comparisons

\item \textbf{Denominator floor and ratio bias.}
Since \(0\le1-p_j\le1\), every weight satisfies
\[
w_j=np_j\left[1-(1-p_j)^{n-1}\right]\ge0.
\]
The population denominator identity and the cellwise lower bound therefore imply
\[
\overline D
=\sum_{j=1}^k w_jd_j
\ge\frac38\sum_{j=1}^k w_j
=\frac38W(P)
\ge\frac{3\mathcal A}{64}
=d_0,
\]
where the formula for \(d_0\) is supplied by
\cref{obj:def:private-ratio}. Because \(n>0\), \(h>0\), and \(\delta>0\),
\[
\mathcal A=nh\min\{1,n\delta\}>0,
\qquad
d_0=\frac{3\mathcal A}{64}>0.
\]
In particular, \(\overline D>0\). The common nonnegative weights also preserve the bias certificates:
\[
\begin{aligned}
|\overline N-\theta(P)\overline D|
&=
\left|\sum_{j=1}^k
w_j\bigl(\nu_j-\theta(P)d_j\bigr)\right|\\
&\le
\sum_{j=1}^k w_j
 \left|\nu_j-\theta(P)d_j\right|\\
&\le
\sum_{j=1}^k w_jbd_j
=b\overline D.
\end{aligned}
\]
Finally, division by the positive denominator gives
\[
\left|\frac{\overline N}{\overline D}-\theta(P)\right|
=
\frac{|\overline N-\theta(P)\overline D|}{\overline D}
\le b.
\]
% lean: CausalSmith.Stat.PrivateCateRoughdesign.occupancy_ratio_assembly
\end{enumerate}
\end{proof}

The same nonnegative occupancy weights enter both population moments. Consequently, variation in cell probabilities changes their common weighting while preserving the cellwise approximation to the target. The bound \(d_j\ge3/8\) expresses the contribution of overlap to denominator strength, and the comparison of \(W(P)\) with \(\mathcal A\) translates random occupancy into a public scale. The bias bound separates localization of the Lipschitz contrast from within-cell approximation of the two Hölder nuisance regressions.

\subsection{Stability across all cell counts}

Privacy compares arbitrary replacement-adjacent datasets in \(\mathcal O^n\), whereas variance concerns independent records sampled from the model. The next result provides the respective bounds for the same statistics. Its replacement statement includes changes in cell membership and occupancies zero, one, and two.

\begin{lemmav}{lem:all-count-stability}[All-count stability and variance]*
Let \(S_N(D)\) and \(S_D(D)\) be the sums over the \(k\) public cells of the count-weighted within-cell treatment–outcome pair contributions and treatment pair contributions, respectively. Let \(W(P)\), the total population occupancy weight, be
\[
W(P)=\sum_{j=1}^{k}w_j,
\]
where \(w_j\) is the population occupancy weight of cell \(j\). Suppose:
\begin{itemize}
\item \textbf{(Parameters.)} The integers \(n,k\) satisfy \(n\ge2\) and \(k\ge2\), and the public window radius \(h\) satisfies \(0<h\le1/4\).
\item \textbf{(Model.)} The causal law \(P\) belongs to \(\mathcal P\) as specified in \cref{obj:def:complete-model}.
\item \textbf{(Sampling.)} The dataset law satisfies \cref{obj:ass:iid}.
\end{itemize}
For every pair of ordered datasets \(D,D^{\prime}\in\mathcal O^n\) whose replacement Hamming distance, the number of differing record positions, is \(d_{\mathrm H}(D,D^{\prime})=1\),
\[
\left|S_N(D)-S_N(D^{\prime})\right|
+\left|S_D(D)-S_D(D^{\prime})\right|
\le12.
\]
Under the stated dataset law,
\[
\max\left\{
\operatorname{Var}(S_N(D)),
\operatorname{Var}(S_D(D))
\right\}
\le18W(P).
\]
The replacement bound applies to every input dataset, including cells with occupancies zero, one, and two.
\end{lemmav}
% CAUSALSMITH-CITED-SCOPE-BEGIN lem:all-count-stability
\verificationfootnotetext{\textbf{Formalization scope.} This result uses the cited conclusion \citep{BoucheronLugosiMassart2013} (Theorem 3.1, replacement-copy form, printed page 54). The cited source proof is not formalized here: Lean verifies its use and all remaining steps. If that cited conclusion is false, the portions of this result that depend on it are not certified.}
% CAUSALSMITH-CITED-SCOPE-END lem:all-count-stability

\begin{proof}[Proof of \cref{obj:lem:all-count-stability}]
\begin{enumerate}
\item We first establish the insertion bounds for the two scalar pair contributions. For observed records \(o=(x,a,y)\) and \(o'=(x',a',y')\) in \(\mathcal O\), define
\[
X(o)=x,\qquad A(o)=a,\qquad Y(o)=y,
\]
\[
K_N(o,o')=(a-a')(y-y'),
\qquad
K_D(o,o')=(a-a')^2.
\]
Both kernels are symmetric. Since the treatment and outcome coordinates are binary,
\[
|K_N(o,o')|\le K_D(o,o')\le1,
\qquad
K_D(o,o')\ge0.
\]
Indeed, equal treatment coordinates make both kernels zero; unequal treatment coordinates give \(K_D=1\) and \(|K_N|=|y-y'|\le1\).

For \(\kappa\in\{K_N,K_D\}\), a nonnegative integer \(s\), and an ordered tuple \(\mathscr R=(o_1,\ldots,o_s)\in\mathcal O^s\), define
\[
F_\kappa(\mathscr R)=
\begin{cases}
0,&s<2,\\[2pt]
\displaystyle\frac{s}{\binom{s}{2}}
\sum_{1\le i<i'\le s}\kappa(o_i,o_{i'}),&s\ge2.
\end{cases}
\]
In particular, \(F_\kappa\) is defined on the empty tuple.

These pair contributions agree with the closed cell formulas in \cref{obj:def:private-ratio}. To see this, set
\[
a_i=A(o_i),\qquad y_i=Y(o_i)\qquad(1\le i\le s).
\]
Symmetry and the vanishing diagonal give
\[
\begin{aligned}
\sum_{i<i'}K_N(o_i,o_{i'})
&=\frac12\sum_{i=1}^s\sum_{i'=1}^s
(a_i-a_{i'})(y_i-y_{i'})\\
&=s\sum_{i=1}^s a_i y_i
-\left(\sum_{i=1}^s a_i\right)
 \left(\sum_{i=1}^s y_i\right),\\
\sum_{i<i'}K_D(o_i,o_{i'})
&=\frac12\sum_{i=1}^s\sum_{i'=1}^s(a_i-a_{i'})^2\\
&=s\sum_{i=1}^s a_i-\left(\sum_{i=1}^s a_i\right)^2,
\end{aligned}
\]
where the last equality uses \(a_i^2=a_i\). Multiplication by
\(s/\binom{s}{2}=2/(s-1)\) for \(s\ge2\) produces the stated cell formulas; for \(s<2\), both constructions give zero.
% lean: CausalSmith.Stat.PrivateCateRoughdesign.pair_moment_formulas

The unit bounds on the kernels imply
\[
\left|\sum_{i<i'}\kappa(o_i,o_{i'})\right|
\le\sum_{i<i'}|\kappa(o_i,o_{i'})|
\le\binom{s}{2},
\qquad
|F_\kappa(\mathscr R)|\le s.
\]
For \(s\ge2\), separating the pairs involving an appended record \(o\in\mathcal O\) gives the exact identity
\[
F_\kappa(o_1,\ldots,o_s,o)-F_\kappa(o_1,\ldots,o_s)
=
\frac2s\sum_{i=1}^s\kappa(o_i,o)
-\frac{F_\kappa(o_1,\ldots,o_s)}s.
\]
Consequently,
\[
\begin{aligned}
\left|F_\kappa(o_1,\ldots,o_s,o)-F_\kappa(o_1,\ldots,o_s)\right|
&\le\frac2s\left|\sum_{i=1}^s\kappa(o_i,o)\right|
+\frac{|F_\kappa(o_1,\ldots,o_s)|}{s}\\
&\le\frac2s\,s+\frac{s}{s}=3.
\end{aligned}
\]
For the remaining occupancies, the defining formulas give
\[
F_\kappa(o)-F_\kappa(\varnothing)=0,
\qquad
F_\kappa(o_1,o)-F_\kappa(o_1)=2\kappa(o_1,o),
\]
so the respective absolute increments are zero and at most \(2\le3\). Thus, at every occupancy, each scalar insertion increment is at most \(3\), and the sum of the two absolute insertion increments is at most \(6\).
% lean: CausalSmith.Stat.PrivateCateRoughdesign.pairContribution_insertion_bound

\item We now prove the replacement bound. Use the public cells \(C_j\) from \cref{obj:synth_11}. For every dataset \(D\in\mathcal O^n\), cell \(j\), and record position \(i_\star\in\{1,\ldots,n\}\), define
\[
\begin{aligned}
\mathscr R_j(D)&=(D_i)_{\,i:\,X(D_i)\in C_j},
&
s_j(D)&=|\mathscr R_j(D)|,\\
\mathscr R_{j,-i_\star}(D)
&=(D_i)_{\,i:\,i\ne i_\star,\ X(D_i)\in C_j},
&
F_{\kappa,j}(D)&=F_\kappa(\mathscr R_j(D)),\\
F_{\kappa,j}^{(-i_\star)}(D)
&=F_\kappa(\mathscr R_{j,-i_\star}(D)),
&
S_\kappa(D)&=\sum_{j=1}^kF_{\kappa,j}(D).
\end{aligned}
\]
The tuples are listed in increasing record position. Symmetry makes their pair contributions invariant under reordering, so the preceding calculation yields
\[
S_{K_N}(D)=S_N(D),\qquad S_{K_D}(D)=S_D(D).
\]

If \(D\) and \(D'\) have replacement Hamming distance one, their set of differing positions is a singleton. Write its position as \(i_\star\); then
\[
D_i=D_i'\qquad(i\ne i_\star).
\]
The erased cell tuples are therefore identical, and hence
\[
F_{\kappa,j}^{(-i_\star)}(D)
=F_{\kappa,j}^{(-i_\star)}(D').
\]
If the erased position belongs to the cell, reversing its insertion and applying the joint insertion bound gives
\[
\sum_{\kappa\in\{K_N,K_D\}}
\left|F_{\kappa,j}(D)-F_{\kappa,j}^{(-i_\star)}(D)\right|
\le6\,\mathbf1\{X(D_{i_\star})\in C_j\}.
\]
If it belongs elsewhere, the erased and original tuples coincide and the left side is zero. Applying this bound to both datasets and using their common erased contribution gives
\[
\sum_{\kappa\in\{K_N,K_D\}}
|F_{\kappa,j}(D)-F_{\kappa,j}(D')|
\le
6\,\mathbf1\{X(D_{i_\star})\in C_j\}
+6\,\mathbf1\{X(D'_{i_\star})\in C_j\}.
\]
The cells form a disjoint partition of the public window, so every covariate belongs to at most one cell. Thus
\[
\begin{aligned}
|S_N(D)-S_N(D')|+|S_D(D)-S_D(D')|
&\le\sum_{j=1}^k\sum_{\kappa\in\{K_N,K_D\}}
|F_{\kappa,j}(D)-F_{\kappa,j}(D')|\\
&\le6\sum_{j=1}^k\mathbf1\{X(D_{i_\star})\in C_j\}
+6\sum_{j=1}^k\mathbf1\{X(D'_{i_\star})\in C_j\}\\
&\le12.
\end{aligned}
\]
This argument includes empty cells, singleton cells, and the transition to occupancy two.
% lean: CausalSmith.Stat.PrivateCateRoughdesign.statistic_hamming_sensitivity

\item For the variance calculation, take a dataset and an independent replacement record with joint law
\[
\mathcal L(D,o^\star)
=(P^{\mathrm{obs}})^{\otimes n}\otimes P^{\mathrm{obs}}.
\]
All probabilities and expectations in this step and the next are under this law. For every \(j\) and \(i_\star\), define
\[
p_j=P_X(C_j)
=P^{\mathrm{obs}}\{o:X(o)\in C_j\},
\qquad
B_{j,i_\star}(D)=|\mathscr R_{j,-i_\star}(D)|.
\]
The occupancy weights and their sum are
\[
w_j=np_j\bigl[1-(1-p_j)^{n-1}\bigr],
\qquad
W(P)=\sum_{j=1}^k w_j.
\]
For every proposed record \(o\in\mathcal O\), define its partner indicator by
\[
J_{i_\star}(D,o)
=
\mathbf1\{\exists j:\ X(o)\in C_j,\ B_{j,i_\star}(D)\ge1\},
\]
and define the old and new indicators by
\[
J_{i_\star}^{\mathrm{old}}=J_{i_\star}(D,D_{i_\star}),
\qquad
J_{i_\star}^{\mathrm{new}}=J_{i_\star}(D,o^\star).
\]
Disjointness of the cells gives the identity
\[
J_{i_\star}(D,o)
=\sum_{j=1}^k
\mathbf1\{X(o)\in C_j,\ B_{j,i_\star}(D)\ge1\}.
\]

The product law evaluates the event that all unreplaced records avoid cell \(j\), and the event that the old record is isolated in that cell, as
\[
\begin{aligned}
\Pr\{B_{j,i_\star}(D)=0\}
&=\prod_{i\ne i_\star}(1-p_j)
=(1-p_j)^{n-1},\\
\Pr\{X(D_{i_\star})\in C_j,\ B_{j,i_\star}(D)=0\}
&=p_j\prod_{i\ne i_\star}(1-p_j)
=p_j(1-p_j)^{n-1}.
\end{aligned}
\]
Subtracting the isolated-record event from the old record's cell-membership event yields
\[
\begin{aligned}
\Pr\{X(D_{i_\star})\in C_j,\ B_{j,i_\star}(D)\ge1\}
&=p_j-p_j(1-p_j)^{n-1}\\
&=p_j\bigl[1-(1-p_j)^{n-1}\bigr].
\end{aligned}
\]
For the independent replacement, the corresponding event is a product of its cell-membership event and the event that the unreplaced records have a partner. Hence
\[
\Pr\{X(o^\star)\in C_j,\ B_{j,i_\star}(D)\ge1\}
=p_j\bigl[1-(1-p_j)^{n-1}\bigr].
\]
These formulas include \(p_j=0\) and \(p_j=1\), since \(n-1\ge1\). Summing the cell indicators and using \(n>0\), we obtain
\[
\begin{aligned}
\mathbb E J_{i_\star}^{\mathrm{old}}
=\mathbb E J_{i_\star}^{\mathrm{new}}
&=\sum_{j=1}^k p_j\bigl[1-(1-p_j)^{n-1}\bigr]
=\frac{W(P)}n,\\
\mathbb E\bigl[J_{i_\star}^{\mathrm{old}}
+J_{i_\star}^{\mathrm{new}}\bigr]
&=\frac{2W(P)}n.
\end{aligned}
\]
% lean: CausalSmith.Stat.PrivateCateRoughdesign.iid_partner_sum_expectation

\item Define the replacement dataset, for every \(D\in\mathcal O^n\) and \(o^\star\in\mathcal O\), by
\[
D_i^{(i_\star)}
=
\begin{cases}
o^\star,&i=i_\star,\\
D_i,&i\ne i_\star.
\end{cases}
\]
The scalar insertion bound sharpens when a cell has no remaining partner. For each kernel,
\[
\left|F_{\kappa,j}(D)-F_{\kappa,j}^{(-i_\star)}(D)\right|
\le
3\,\mathbf1\{X(D_{i_\star})\in C_j,\ B_{j,i_\star}(D)\ge1\}.
\]
Indeed, if the record belongs to the cell and a partner remains, its erasure reverses a scalar insertion of size at most \(3\). If it belongs to the cell and no partner remains, the original tuple is a singleton and the erased tuple is empty, so both contributions are zero. If it belongs elsewhere, erasure leaves the cell tuple unchanged.

The original and replacement datasets have the same erased cell tuples. Applying the preceding estimate to both gives
\[
\begin{aligned}
|F_{\kappa,j}(D)-F_{\kappa,j}(D^{(i_\star)})|
\le{}&
3\,\mathbf1\{X(D_{i_\star})\in C_j,\ B_{j,i_\star}(D)\ge1\}\\
&+3\,\mathbf1\{X(o^\star)\in C_j,\ B_{j,i_\star}(D)\ge1\}.
\end{aligned}
\]
Summing over cells and using the partner-indicator identity therefore yields
\[
|S_\kappa(D)-S_\kappa(D^{(i_\star)})|
\le3\bigl(J_{i_\star}^{\mathrm{old}}+J_{i_\star}^{\mathrm{new}}\bigr).
\]
Both indicators take values in \(\{0,1\}\). Squaring this bound and checking their four possible pairs gives
\[
\begin{aligned}
\bigl(S_\kappa(D)-S_\kappa(D^{(i_\star)})\bigr)^2
&\le9\bigl(J_{i_\star}^{\mathrm{old}}+J_{i_\star}^{\mathrm{new}}\bigr)^2\\
&\le18\bigl(J_{i_\star}^{\mathrm{old}}+J_{i_\star}^{\mathrm{new}}\bigr).
\end{aligned}
\]
For the pairs \((0,0),(1,0),(0,1),(1,1)\), the middle expression is respectively \(0,9,9,36\), while the final expression is \(0,18,18,36\).
% lean: CausalSmith.Stat.PrivateCateRoughdesign.statistic_partner_square_bound

The bounded measurable statistics and indicators make these expressions integrable. Taking expectations and applying the partner expectation from the preceding step gives, for every record position,
\[
\begin{aligned}
\mathbb E\bigl(S_N(D)-S_N(D^{(i_\star)})\bigr)^2
&\le18\,\mathbb E\bigl[J_{i_\star}^{\mathrm{old}}
+J_{i_\star}^{\mathrm{new}}\bigr]
=\frac{36W(P)}n,\\
\mathbb E\bigl(S_D(D)-S_D(D^{(i_\star)})\bigr)^2
&\le18\,\mathbb E\bigl[J_{i_\star}^{\mathrm{old}}
+J_{i_\star}^{\mathrm{new}}\bigr]
=\frac{36W(P)}n.
\end{aligned}
\]
% lean: CausalSmith.Stat.PrivateCateRoughdesign.statistic_replacement_moments_of_partners

\item To apply the replacement Efron--Stein inequality, observe that disjoint cell membership and the scalar contribution bound imply, for every dataset,
\[
|S_\kappa(D)|
\le\sum_{j=1}^k|F_{\kappa,j}(D)|
\le\sum_{j=1}^k s_j(D)
=\sum_{i=1}^n\sum_{j=1}^k\mathbf1\{X(D_i)\in C_j\}
\le n.
\]
The cell formulas are measurable, and this uniform bound gives square integrability.

The replacement-copy form of the Efron--Stein inequality
\citep[Theorem 3.1, p.~54]{BoucheronLugosiMassart2013}
now gives, for each \(\kappa\in\{K_N,K_D\}\),
\[
\begin{aligned}
\operatorname{Var}(S_\kappa(D))
&\le\frac12\sum_{i_\star=1}^n
\mathbb E\bigl(S_\kappa(D)-S_\kappa(D^{(i_\star)})\bigr)^2\\
&\le\frac12\sum_{i_\star=1}^n\frac{36W(P)}n
=18W(P).
\end{aligned}
\]
Consequently,
\[
\max\{\operatorname{Var}(S_N(D)),\operatorname{Var}(S_D(D))\}
\le18W(P).
\]
Finally, \cref{obj:ass:iid} identifies the stated dataset law with the product law used in this calculation, proving the asserted variance bound.
% lean: CausalSmith.Stat.PrivateCateRoughdesign.all_count_stability
\end{enumerate}
\end{proof}

Count weighting allows a record to contribute to several pairs while retaining a uniform bound on the change in the joint statistic. The all-count statement covers transitions between the zero contribution for occupancies below two and the pair contribution for larger occupancies. Its constant \(12\) is the sensitivity used to calibrate the joint Laplace release in \cref{obj:def:private-ratio}.

The variance argument uses the replacement-copy inequality of \citet[Theorem 3.1, replacement-copy form, p.~54]{BoucheronLugosiMassart2013}. For a square-integrable measurable statistic \(Z=\Psi(D)\) of independent records, let \(D^{(i)}\) replace record \(i\) by an independent copy. The cited inequality states
\[
\operatorname{Var}(Z)
\le
\frac12\sum_{i=1}^n
E\bigl\{\bigl(\Psi(D)-\Psi(D^{(i)})\bigr)^2\bigr\}.
\]
In \cref{obj:lem:all-count-stability}, this comparison yields variance control in terms of \(W(P)\). That occupancy dependence connects sampling variation to the same scale that controls the population denominator.

\subsection{Role in the uniform upper bound}

The public envelope in \cref{obj:thm:uniform-private-upper} combines the preceding ingredients through
\[
V(h,\delta)
=
2^{20}\left(
h^2+\delta^{2/5}
+\mathcal A^{-1}
+(\epsilon\mathcal A)^{-2}
\right).
\]
The first two terms accommodate the squared population bias from \cref{obj:lem:occupancy-cancellation}. The variance bound in \cref{obj:lem:all-count-stability}, together with the population denominator and occupancy comparisons, supplies the sampling term \(\mathcal A^{-1}\). Independent Laplace perturbations of scale \(12/\epsilon\) have second moments \(288/\epsilon^2\); denominator stabilization places their contribution on the scale \((\epsilon\mathcal A)^{-2}\). The public floor keeps the ratio defined throughout the dataset space, and clipping respects the target range established in \cref{obj:lem:point-version}.

The joint sensitivity bound supports pure replacement privacy for the release, and deterministic postprocessing preserves that guarantee \citep[Theorem 3.6 and Proposition 2.1]{Dwork2014}. Thus the estimator, centered interval, and candidate-value inversion share the privacy guarantee of the same released coordinates. The squared-error envelope supports the radius \(\sqrt{10V(h,\delta)}\) used for uniform \(90\%\) coverage in \cref{obj:thm:uniform-private-upper}. Truncation to the target range also supplies the stated upper bound of two on interval length.

Public tuning in \cref{obj:def:public-tuning} sets \(\delta=h^5\) in the local branch, aligning the localization and within-cell squared-bias scales. The benchmark accounts for sampling variation and privacy noise across sparse and dense occupancy regimes. For the complementary public branch, the constant estimate and full target interval supply the outputs specified there. These choices are the procedures covered by the uniform absolute-error and expected-length bounds in \cref{obj:thm:uniform-private-upper}.

Finally, \cref{obj:def:interval-handle} expresses interval construction through candidate values using the floored denominator \(B_D\), clipped numerator \(B_N\), and public radius \(\rho\) defined there. Its acceptance rule describes the same target-range-truncated interval centered on the private ratio. Inversion therefore uses the numerator and denominator already released for estimation, with the public envelope controlling the acceptance radius.


\section{Lower-bound experiments and component coupling}

The lower-bound experiments localize a treatment-effect separation while varying the nuisance functions within the model class of \cref{obj:def:complete-model}. Their internal coordinates serve two comparisons: a chi-square bound for the observed-sample mixture and, under sparse occupancy, a coupling bound for private output laws. These constructions support the testing arguments underlying \cref{obj:thm:matched-risk-frontier,obj:thm:sharp-interval-frontier}.

The geometry uses the target location \(x_0=1/2\), a macro localization radius \leanref{sym:hL}{\ensuremath{h_{\mathrm L}}} in \((0,1/4]\), and a micro localization radius \leanref{sym:deltaL}{\ensuremath{\delta_{\mathrm L}}} calibrated as \(\delta_{\mathrm L}=h_{\mathrm L}^{5}\). The sign vector \leanref{sym:lambda}{\ensuremath{\lambda}} indexes perturbations whose micro supports intersect the macro window. Binary coordinates encode the signed coefficients, and the target-separation amplitude \leanref{sym:t}{\ensuremath{t}} scales with the macro radius. The following definitions specify the indexing and amplitude.

\begin{definitionv}{synth_6}[Binary sign vector encoding]
The sign vector \(\lambda\) is a binary encoding of signs in \(\{-1,1\}\), with one coordinate for each active integer index determined by the location \(x_0\), macro radius \(h_{\mathrm L}\), and micro radius \(\delta_{\mathrm L}\). Specifically, \(\lambda\) is any function
\[
\lambda:
\left\{
j\in\mathbb Z:
\begin{gathered}
\left\lfloor\frac{x_0-h_{\mathrm L}}{\delta_{\mathrm L}}\right\rfloor-1
\le j\le
\left\lceil\frac{x_0+h_{\mathrm L}}{\delta_{\mathrm L}}\right\rceil+1,\\
j\delta_{\mathrm L}-\delta_{\mathrm L}<x_0+h_{\mathrm L},
\qquad
x_0-h_{\mathrm L}<j\delta_{\mathrm L}+\delta_{\mathrm L}
\end{gathered}
\right\}
\longrightarrow \{0,1\}.
\]
\end{definitionv}

The active indices in \cref{obj:synth_6} retain every micro support that meets the interior of the localization window. With this finite indexing fixed, the separation is specified by the macro scale.

\begin{definitionv}{synth_5}[Target separation amplitude]
The target-separation amplitude \(t\) is defined by
\[
t := 2^{-12}h_{\mathrm L},
\]
where \(h_{\mathrm L}\) is the macro localization radius for the lower-bound experiments.
\end{definitionv}

The localized contrast profile \leanref{sym:q}{\ensuremath{q(x)}} is a fourth-power cosine on the macro window. The nuisance perturbation \leanref{sym:S}{\ensuremath{S_\lambda(x)}} combines a squared-cosine envelope with overlapping micro cosine profiles carrying the coefficients \(2\lambda_j-1\). Both profiles vanish outside the macro window. Uniform averaging over the sign vectors assigns independent fair binary coordinates to the active indices; a single draw of this vector is shared across the records in each alternative sample.

For comparison, the causal null \leanref{sym:P0}{\ensuremath{P_0}} has a uniform covariate and, conditional on that covariate, mutually independent fair binary treatment and potential outcomes, with \(Y=Y(A)\). Its observed treatment--outcome cells each have conditional probability \(1/4\). The treatment mark \leanref{sym:U}{\ensuremath{U}} and outcome mark \(V_Y\) encode the observed binary values as
\[
U=2A-1,\qquad V_Y=2Y-1.
\]
These marks express the conditional likelihood relative to the fair null. The profiles and sign distribution now determine the causal alternative \leanref{sym:Plambda}{\ensuremath{P_\lambda}} and the observed-sample mixture \leanref{sym:Qn}{\ensuremath{Q_n}} in the following construction.

\begin{definitionv}{def:cosine-family}[Shared-sign causal alternative family]
The causal alternatives \(P_\lambda\) and their observed-sample mixture \(Q_n\) are defined with localization center \(x_0=1/2\), macro localization radius \(0<h_{\mathrm L}\le 1/4\), micro localization radius \(\delta_{\mathrm L}\), and target-separation amplitude \(t\), where
\[
\delta_{\mathrm L}=h_{\mathrm L}^{5},
\qquad
t=2^{-12}h_{\mathrm L}.
\]
Each sign vector \(\lambda\) assigns a binary value \(\lambda_j\in\{0,1\}\) to every active integer index in
\[
\left\{
j\in\mathbb Z:
\left\lfloor\frac{x_0-h_{\mathrm L}}{\delta_{\mathrm L}}\right\rfloor-1
\le j\le
\left\lceil\frac{x_0+h_{\mathrm L}}{\delta_{\mathrm L}}\right\rceil+1,\quad
j\delta_{\mathrm L}-\delta_{\mathrm L}<x_0+h_{\mathrm L},\quad
x_0-h_{\mathrm L}<j\delta_{\mathrm L}+\delta_{\mathrm L}
\right\}.
\]
The corresponding signed coefficient is \(2\lambda_j-1\in\{-1,1\}\). The localized contrast profile \(q(x)\) and shared-sign perturbation \(S_\lambda(x)\) are
\[
q(x)=
\begin{cases}
\displaystyle
\cos^{4}\!\left(\frac{\pi(x-x_0)}{2h_{\mathrm L}}\right),
& |x-x_0|\le h_{\mathrm L},\\
0,& |x-x_0|>h_{\mathrm L},
\end{cases}
\]
and
\[
S_\lambda(x)=
\begin{cases}
\displaystyle
\cos^{2}\!\left(\frac{\pi(x-x_0)}{2h_{\mathrm L}}\right)
\sum_j(2\lambda_j-1)
\begin{cases}
\displaystyle
\cos\!\left(\frac{\pi(x-j\delta_{\mathrm L})}{2\delta_{\mathrm L}}\right),
& |x-j\delta_{\mathrm L}|\le\delta_{\mathrm L},\\
0,& |x-j\delta_{\mathrm L}|>\delta_{\mathrm L},
\end{cases}
& |x-x_0|\le h_{\mathrm L},\\
0,& |x-x_0|>h_{\mathrm L},
\end{cases}
\]
where the sum runs over the active indices above.

For each \(\lambda\), the causal law \(P_\lambda\) has covariate \(X\sim\operatorname{Unif}[0,1]\). Conditional on \(X=x\), the binary treatment indicator \(A\) and binary potential outcomes \(Y(0),Y(1)\) are mutually independent, with
\[
\begin{aligned}
A\mid X=x
&\sim\operatorname{Bernoulli}
\left(\frac{1+\sqrt t\,S_\lambda(x)}{2}\right),\\
Y(0)\mid X=x
&\sim\operatorname{Bernoulli}
\left(\frac{1-\sqrt t\,S_\lambda(x)-tq(x)}{2}\right),\\
Y(1)\mid X=x
&\sim\operatorname{Bernoulli}
\left(\frac{1-\sqrt t\,S_\lambda(x)+tq(x)}{2}\right).
\end{aligned}
\]
The observed outcome is \(Y=Y(A)\). The selected covariate density is \(f_{P_\lambda}(x)=1\), and the displayed Bernoulli parameters specify the selected propensity and arm-mean versions.

For every nonnegative integer \(n\), let \(P_\lambda^{\mathrm{obs}}\) be the marginal law of \((X,A,Y)\). Define
\[
Q_n=E_\lambda\!\left[(P_\lambda^{\mathrm{obs}})^{\otimes n}\right],
\]
where \(E_\lambda\) is uniform averaging over all binary sign vectors on the active index set, equivalently assigning independent fair binary values to its indices.
\end{definitionv}

The common nuisance perturbation in \cref{obj:def:cosine-family} enters both the propensity and the potential-outcome means, while their difference equals \(tq(x)\). Consequently, the sign vector changes the nuisance functions while preserving the target separation. The conditional likelihood \leanref{sym:likelihood}{\ensuremath{L_\lambda(x,U,V_Y)}} is four times the alternative conditional cell probability, using the fair-null cell probability as its denominator. The next result certifies model membership, the common target value, and the mixture comparison for this explicit family.

\begin{lemmav}{lem:positive-family-certificate}[Positive family and mixture certificate]*
For every \(h_{\mathrm L}\in(0,1/4]\), the macro localization radius, and every integer \(n\ge0\), set the target-separation amplitude \(t\) and micro localization radius \(\delta_{\mathrm L}\) to
\[
t=2^{-12}h_{\mathrm L},
\qquad
\delta_{\mathrm L}=h_{\mathrm L}^{5}.
\]
Let \(P_0\) be the causal null with a uniform covariate and conditionally independent fair binary treatment and potential outcomes, with the observed outcome determined by consistency. Use the alternatives \(P_\lambda\), localized profile \(q(x)\), nuisance perturbation \(S_\lambda(x)\), uniform sign average \(E_\lambda\), and mixture \(Q_n\) from \cref{obj:def:cosine-family}. Throughout, \(\lambda\) is the binary encoding on the active frame indices:
\[
\lambda_j\in\{0,1\},
\qquad
2\lambda_j-1\in\{-1,1\}.
\]
Thus the coefficients of the cosine-frame sum defining \(S_\lambda(x)\) are \(2\lambda_j-1\).

For the model class \(\mathcal P\) in \cref{obj:def:complete-model} and the point target \(\theta(P)=\tau_P(1/2)\),
\[
P_0\in\mathcal P,\qquad
\theta(P_0)=0,\qquad
P_\lambda\in\mathcal P,\qquad
\theta(P_\lambda)=t
\quad\text{for every }\lambda.
\]
For every sign vector \(\lambda\), covariate value \(x\), and binary treatment and outcome values \(A,Y\), define the likelihood marks
\[
U=2A-1,\qquad V_Y=2Y-1.
\]
The conditional likelihood \(L_\lambda(x,U,V_Y)\), equal to four times the alternative conditional probability of the cell \((A,Y)\) given \(x\), satisfies
\[
\begin{aligned}
L_\lambda(x,U,V_Y)
={}&1+U\sqrt t\,S_\lambda(x)
 +V_Y\sqrt t\,\bigl(tq(x)-1\bigr)S_\lambda(x)\\
&+UV_Yt\bigl(q(x)-S_\lambda(x)^2\bigr).
\end{aligned}
\]
The one-record mixture equals the observed null law:
\[
E_\lambda P_\lambda^{\mathrm{obs}}=P_0^{\mathrm{obs}}.
\]
Moreover,
\[
\chi^2\!\left(Q_n,(P_0^{\mathrm{obs}})^{\otimes n}\right)
\le
\exp\!\left(160n^2t^2h_{\mathrm L}\delta_{\mathrm L}\right)-1,
\qquad
Q_n\ll(P_0^{\mathrm{obs}})^{\otimes n},
\]
and the centered likelihood ratio is square-integrable:
\[
\int_{\mathcal O^n}
\left(
\frac{dQ_n}{d(P_0^{\mathrm{obs}})^{\otimes n}}-1
\right)^2
\,d(P_0^{\mathrm{obs}})^{\otimes n}
<\infty.
\]
\end{lemmav}
% CAUSALSMITH-CITED-SCOPE-BEGIN lem:positive-family-certificate
\verificationfootnotetext{\textbf{Formalization scope.} This result uses the cited conclusion \citep{BoucheronLugosiMassart2013} (Theorem 6.2, displayed log-mgf conclusion in its proof, printed page 167). The cited source proof is not formalized here: Lean verifies its use and all remaining steps. If that cited conclusion is false, the portions of this result that depend on it are not certified.}
% CAUSALSMITH-CITED-SCOPE-END lem:positive-family-certificate

\begin{proof}[Proof of \cref{obj:lem:positive-family-certificate}]
\begin{enumerate}
\item \textbf{Frame identities.}
Write the active index set from \cref{obj:def:cosine-family} as
\[
\begin{split}
\mathcal J=\biggl\{j\in\mathbb Z:\;&
\left\lfloor\frac{x_0-h_{\mathrm L}}{\delta_{\mathrm L}}\right\rfloor-1
\le j\le
\left\lceil\frac{x_0+h_{\mathrm L}}{\delta_{\mathrm L}}\right\rceil+1,\\
&j\delta_{\mathrm L}-\delta_{\mathrm L}<x_0+h_{\mathrm L},
\quad
x_0-h_{\mathrm L}<j\delta_{\mathrm L}+\delta_{\mathrm L}
\biggr\}.
\end{split}
\]
For \(x\in[0,1]\), \(j\in\mathbb Z\), and
\(\lambda\in\{0,1\}^{\mathcal J}\), define
\[
g(x)=
\begin{cases}
\cos^2\!\left(\dfrac{\pi(x-x_0)}{2h_{\mathrm L}}\right),
 & |x-x_0|\le h_{\mathrm L},\\
0,& |x-x_0|>h_{\mathrm L},
\end{cases}
\qquad
\phi_j(x)=
\begin{cases}
\cos\!\left(\dfrac{\pi(x-j\delta_{\mathrm L})}{2\delta_{\mathrm L}}\right),
 & |x-j\delta_{\mathrm L}|\le\delta_{\mathrm L},\\
0,& |x-j\delta_{\mathrm L}|>\delta_{\mathrm L},
\end{cases}
\]
and
\[
\sigma_{\lambda,j}=2\lambda_j-1\quad(j\in\mathcal J),
\qquad
F_\lambda(x)=\sum_{j\in\mathcal J}\sigma_{\lambda,j}\phi_j(x).
\]
Thus
\[
q(x)=g(x)^2,\qquad S_\lambda(x)=g(x)F_\lambda(x),
\qquad
0\le g(x)\le1,\quad 0\le q(x)\le1,\quad |\phi_j(x)|\le1.
\]
The frame functions vanish at their support endpoints.

To establish the partition identity, put
\[
j_x=\left\lfloor\frac{x}{\delta_{\mathrm L}}\right\rfloor
\quad(x\in[0,1]).
\]
Then \(j_x\delta_{\mathrm L}\le x<(j_x+1)\delta_{\mathrm L}\).
For \(j<j_x\) or \(j>j_x+1\), the distance from \(x\) to
\(j\delta_{\mathrm L}\) is at least \(\delta_{\mathrm L}\), so
\(\phi_j(x)=0\). The two remaining terms satisfy
\[
\begin{aligned}
\phi_{j_x}(x)
 &=\cos\!\left(\frac{\pi(x-j_x\delta_{\mathrm L})}
                         {2\delta_{\mathrm L}}\right),\\
\phi_{j_x+1}(x)
 &=\sin\!\left(\frac{\pi(x-j_x\delta_{\mathrm L})}
                         {2\delta_{\mathrm L}}\right),\\
\phi_{j_x}(x)^2+\phi_{j_x+1}(x)^2&=1.
\end{aligned}
\]
If \(|x-x_0|\ge h_{\mathrm L}\), then \(q(x)=0\).
If \(|x-x_0|<h_{\mathrm L}\) and \(j\notin\mathcal J\), the
open support of \(\phi_j\) fails to intersect the open macro interval;
hence \(|x-j\delta_{\mathrm L}|\ge\delta_{\mathrm L}\) and
\(\phi_j(x)=0\). Consequently, including the macro and lattice
boundaries,
\[
q(x)\sum_{j\in\mathcal J}\phi_j(x)^2=q(x).
\]

For later support counts, define
\[
\mathcal J_x=\{j\in\mathcal J:|x-j\delta_{\mathrm L}|
                                      \le\delta_{\mathrm L}\},
\qquad
\mathcal J_x^\circ=\{j\in\mathcal J:\phi_j(x)\ne0\}.
\]
The closed-support condition places \(j\) between \(j_x-1\) and
\(j_x+1\), while the preceding two-frame argument gives
\[
|\mathcal J_x|\le3,\qquad |\mathcal J_x^\circ|\le2,
\qquad
|F_\lambda(x)|\le\sum_{j\in\mathcal J_x}|\phi_j(x)|\le3,
\qquad |S_\lambda(x)|\le3.
\]

Uniform binary signs are independent, and pairing encodings that differ
in one coordinate gives
\[
E_\lambda\sigma_{\lambda,j}=0,
\qquad
E_\lambda(\sigma_{\lambda,j}\sigma_{\lambda,j'})
=
\begin{cases}
1,&j=j',\\
0,&j\ne j'.
\end{cases}
\]
Expanding the finite frame sum and its square therefore yields
\[
E_\lambda S_\lambda(x)=0,
\qquad
E_\lambda S_\lambda(x)^2
=q(x)\sum_{j\in\mathcal J}\phi_j(x)^2=q(x).
\]
% lean: CausalSmith.Stat.PrivateCateRoughdesign.perturbation_sign_sum
% lean: CausalSmith.Stat.PrivateCateRoughdesign.perturbation_sign_square_sum_exact
% lean: CausalSmith.Stat.PrivateCateRoughdesign.weighted_frame_square_partition

\item \textbf{The centered overlap.}
The prescribed amplitude satisfies
\[
0<t\le\frac1{16384},
\qquad
0\le\sqrt t\le\frac1{128},
\qquad
|\sqrt t\,S_\lambda(x)|\le\frac3{128},
\qquad
0\le tq(x)\le\frac1{16384}.
\]
In particular,
\[
\left|\sqrt t\,S_\lambda(x)\right|<1,
\qquad
\left|-\sqrt t\,S_\lambda(x)\pm tq(x)\right|
\le\frac3{128}+\frac1{16384}<1.
\]
Thus all three Bernoulli parameters in
\cref{obj:def:cosine-family} belong to \([0,1]\).

For each of the four pairs
\((U,V_Y)\in\{-1,1\}^2\), with \(U=2A-1\) and \(V_Y=2Y-1\),
the conditional cell probability divided by its null probability
\(1/4\) is
\[
L_\lambda(x,U,V_Y)
=(1+U\sqrt t\,S_\lambda(x))
 \bigl(1+V_Y(-\sqrt t\,S_\lambda(x)+Utq(x))\bigr).
\]
Expanding in these four cases and using \((\sqrt t)^2=t\) gives
\[
\begin{aligned}
L_\lambda(x,U,V_Y)
={}&1+U\sqrt t\,S_\lambda(x)
 +V_Y\sqrt t\,(tq(x)-1)S_\lambda(x)\\
&+UV_Yt\bigl(q(x)-S_\lambda(x)^2\bigr).
\end{aligned}
\]
Since \(L_\lambda\) is four times a conditional cell probability,
\[
0\le L_\lambda(x,U,V_Y)\le4.
\]
% lean: CausalSmith.Stat.PrivateCateRoughdesign.conditionalLikelihood_fourier

For two sign vectors, define the overlap integrand and its integral by
\[
\begin{aligned}
\mathcal H_{\lambda,\lambda'}(x)
={}&t\bigl(1+(tq(x)-1)^2\bigr)S_\lambda(x)S_{\lambda'}(x)\\
&+t^2\bigl(q(x)-S_\lambda(x)^2\bigr)
       \bigl(q(x)-S_{\lambda'}(x)^2\bigr),
\qquad x\in[0,1],\\
H_{\lambda,\lambda'}&=\int_0^1
                         \mathcal H_{\lambda,\lambda'}(x)\,dx.
\end{aligned}
\]
Expanding the four fair mark cells cancels every product of distinct
members of \(1,U,V_Y,UV_Y\), and gives
\begin{equation}
\label{eq:positive-family-fair-mark-identity}
\frac14\sum_{U\in\{-1,1\}}\sum_{V_Y\in\{-1,1\}}
 L_\lambda(x,U,V_Y)L_{\lambda'}(x,U,V_Y)
=1+\mathcal H_{\lambda,\lambda'}(x).
\end{equation}
Each summand lies in \([0,16]\), so
\[
-1\le\mathcal H_{\lambda,\lambda'}(x)\le15,
\qquad
-1\le H_{\lambda,\lambda'}\le15,
\qquad
|H_{\lambda,\lambda'}|\le15.
\]
All the displayed functions are measurable and bounded.

For fixed \(\lambda'\), the sign moments established above imply
\[
\begin{aligned}
E_\lambda\mathcal H_{\lambda,\lambda'}(x)
={}&t\bigl(1+(tq(x)-1)^2\bigr)
       S_{\lambda'}(x)E_\lambda S_\lambda(x)\\
&+t^2\bigl(q(x)-E_\lambda S_\lambda(x)^2\bigr)
       \bigl(q(x)-S_{\lambda'}(x)^2\bigr)
=0.
\end{aligned}
\]
Commuting the finite sign average with the integral gives
\[
E_\lambda H_{\lambda,\lambda'}=0.
\]
% lean: CausalSmith.Stat.PrivateCateRoughdesign.signOverlap_abs_le_fifteen
% lean: CausalSmith.Stat.PrivateCateRoughdesign.signOverlap_sign_sum

\item \textbf{Coordinate changes and their squared budget.}
Cauchy--Schwarz on the nonzero frames gives
\[
\begin{aligned}
S_\lambda(x)^2
&=q(x)\left(
  \sum_{j\in\mathcal J_x^\circ}\sigma_{\lambda,j}\phi_j(x)
       \right)^2\\
&\le q(x)|\mathcal J_x^\circ|
               \sum_{j\in\mathcal J_x^\circ}\phi_j(x)^2\\
&\le2q(x)\sum_{j\in\mathcal J}\phi_j(x)^2
=2q(x).
\end{aligned}
\]
The same calculation covers an empty nonzero-frame set. Hence
\[
0\le S_\lambda(x)^2\le2q(x),
\qquad
|q(x)-S_\lambda(x)^2|\le q(x),
\qquad
|S_\lambda(x)|\le\frac32.
\]

For \(j\in\mathcal J\) and \(\beta\in\{0,1\}\), define the updated
encoding by
\[
\lambda^{(j,\beta)}_{j'}=
\begin{cases}
\beta,&j'=j,\\
\lambda_{j'},&j'\ne j,
\end{cases}
\qquad j'\in\mathcal J.
\]
Subtracting the two frame sums gives
\[
S_\lambda(x)-S_{\lambda^{(j,\beta)}}(x)
=2g(x)(\lambda_j-\beta)\phi_j(x),
\]
and therefore
\[
|S_\lambda(x)-S_{\lambda^{(j,\beta)}}(x)|
\le2g(x)|\phi_j(x)|\le2.
\]
Both squared perturbations belong to \([0,2q(x)]\), so
\[
\left|S_{\lambda^{(j,\beta)}}(x)^2-S_\lambda(x)^2\right|
\le2q(x)\le2.
\]
Also, \(0\le tq(x)\le1\) implies
\[
0\le1+(tq(x)-1)^2\le2.
\]
The exact difference of the overlap integrands is
\[
\begin{aligned}
\mathcal H_{\lambda,\lambda'}(x)
-\mathcal H_{\lambda^{(j,\beta)},\lambda'}(x)
={}&t\bigl(1+(tq(x)-1)^2\bigr)
   \bigl(S_\lambda(x)-S_{\lambda^{(j,\beta)}}(x)\bigr)
   S_{\lambda'}(x)\\
&+t^2\bigl(S_{\lambda^{(j,\beta)}}(x)^2-S_\lambda(x)^2\bigr)
       \bigl(q(x)-S_{\lambda'}(x)^2\bigr).
\end{aligned}
\]
The absolute values of the two displayed terms are bounded, respectively,
by
\[
t\cdot2\cdot2\cdot\frac32=6t,
\qquad
t^2\cdot2\cdot1=2t^2.
\]
Since \(0\le t\le1\), this proves
\[
\left|\mathcal H_{\lambda,\lambda'}(x)
-\mathcal H_{\lambda^{(j,\beta)},\lambda'}(x)\right|
\le6t+2t^2\le8t.
\]

Define the open frame support within the covariate interval by
\[
\mathcal S_j=\{x\in[0,1]:|x-j\delta_{\mathrm L}|
                                      <\delta_{\mathrm L}\}.
\]
Outside \(\mathcal S_j\), including its endpoints,
\(\phi_j(x)=0\), so the updated and original integrands agree.
Moreover,
\[
\mathcal S_j\subset
(j\delta_{\mathrm L}-\delta_{\mathrm L},
 j\delta_{\mathrm L}+\delta_{\mathrm L}),
\qquad
\operatorname{Leb}(\mathcal S_j)\le2\delta_{\mathrm L}.
\]
It follows that
\[
\begin{aligned}
|H_{\lambda,\lambda'}-H_{\lambda^{(j,\beta)},\lambda'}|
&\le\int_0^1
 \left|\mathcal H_{\lambda,\lambda'}(x)
       -\mathcal H_{\lambda^{(j,\beta)},\lambda'}(x)\right|\,dx\\
&\le8t\,\operatorname{Leb}(\mathcal S_j)
\le16t\delta_{\mathrm L}.
\end{aligned}
\]
% lean: CausalSmith.Stat.PrivateCateRoughdesign.signOverlap_update_abs_le

To count the coordinates, the strict intersection conditions defining
\(\mathcal J\) imply
\[
\frac{x_0-h_{\mathrm L}}{\delta_{\mathrm L}}<j+1,
\qquad
j-1<\frac{x_0+h_{\mathrm L}}{\delta_{\mathrm L}}.
\]
Taking integer bounds yields
\[
\mathcal J\subset
\left[
\left\lfloor\frac{x_0-h_{\mathrm L}}{\delta_{\mathrm L}}\right\rfloor,
\left\lceil\frac{x_0+h_{\mathrm L}}{\delta_{\mathrm L}}\right\rceil
\right]\cap\mathbb Z.
\]
Counting this integer interval and using the floor and ceiling bounds,
\[
\begin{aligned}
|\mathcal J|
&\le
\left\lceil\frac{x_0+h_{\mathrm L}}{\delta_{\mathrm L}}\right\rceil
-\left\lfloor\frac{x_0-h_{\mathrm L}}{\delta_{\mathrm L}}\right\rfloor+1\\
&\le\frac{2h_{\mathrm L}}{\delta_{\mathrm L}}+3
\le\frac{5h_{\mathrm L}}{\delta_{\mathrm L}}.
\end{aligned}
\]
The last inequality uses
\(0<\delta_{\mathrm L}=h_{\mathrm L}^5\le h_{\mathrm L}\).
Thus the squared coordinate-change budget satisfies
\[
\begin{aligned}
\sum_{j\in\mathcal J}(16t\delta_{\mathrm L})^2
&=|\mathcal J|(16t\delta_{\mathrm L})^2\\
&\le\frac{5h_{\mathrm L}}{\delta_{\mathrm L}}
                  (16t\delta_{\mathrm L})^2
=1280t^2h_{\mathrm L}\delta_{\mathrm L}.
\end{aligned}
\]
% lean: CausalSmith.Stat.PrivateCateRoughdesign.sign_coordinate_squared_budget

\item \textbf{The joint exponential average.}
For fixed \(\lambda'\), the function
\(\lambda\mapsto H_{\lambda,\lambda'}\) is a bounded measurable
function of independent fair binary coordinates, has mean zero,
and changes by at most \(16t\delta_{\mathrm L}\) in each coordinate.
Apply the bounded-differences exponential-moment conclusion of
\citep[Theorem~6.2, displayed log-mgf bound in its proof,
p.~167]{BoucheronLugosiMassart2013}, with exponential parameter
\(n\ge0\). It gives
\[
\begin{aligned}
E_\lambda\exp(nH_{\lambda,\lambda'})
&\le
\exp\!\left\{\frac{n^2}{8}
       \sum_{j\in\mathcal J}(16t\delta_{\mathrm L})^2\right\}\\
&\le\exp\!\left\{160n^2t^2h_{\mathrm L}\delta_{\mathrm L}\right\}.
\end{aligned}
\]
Here uniform averaging is the finite product expectation:
\[
E_\lambda\exp(nH_{\lambda,\lambda'})
=\frac1{2^{|\mathcal J|}}
  \sum_{\lambda\in\{0,1\}^{\mathcal J}}\exp(nH_{\lambda,\lambda'}).
\]
Its denominator is positive, including for an empty coordinate set,
when the product law is concentrated on the unique empty encoding.

Let \(\lambda'\) be an independent uniform encoding, and define
\(E_{\lambda,\lambda'}=E_{\lambda'}E_\lambda\).
Averaging the fixed-\(\lambda'\) inequality gives
\[
\begin{aligned}
E_{\lambda,\lambda'}\exp(nH_{\lambda,\lambda'})
&=\frac1{2^{2|\mathcal J|}}
 \sum_{\lambda,\lambda'\in\{0,1\}^{\mathcal J}}
                      \exp(nH_{\lambda,\lambda'})\\
&\le\exp\!\left\{160n^2t^2h_{\mathrm L}\delta_{\mathrm L}\right\}.
\end{aligned}
\]
% lean: CausalSmith.Stat.PrivateCateRoughdesign.signOverlap_exp_average_le
% lean: CausalSmith.Stat.PrivateCateRoughdesign.signOverlap_joint_exp_average_le

\item \textbf{From overlap to chi-square divergence.}
Each conditional mass of \((A,Y(0),Y(1))\) under an alternative is
a product of three Bernoulli masses and is at most \(1\).
The corresponding null mass is \(1/8\).
The covariate law is uniform in both experiments, and the consistency
map is shared. Integrating these conditional inequalities and then
taking the observed marginal gives
\[
P_\lambda\le8P_0,
\qquad
P_\lambda^{\mathrm{obs}}\le8P_0^{\mathrm{obs}},
\qquad
P_\lambda^{\mathrm{obs}}\ll P_0^{\mathrm{obs}}.
\]

Use the observed-record space
\(\mathcal O=[0,1]\times\{0,1\}^2\), and choose the one-record density
version
\[
\mathscr L_\lambda(x,A,Y)
=L_\lambda(x,2A-1,2Y-1)
=\frac{dP_\lambda^{\mathrm{obs}}}{dP_0^{\mathrm{obs}}}(x,A,Y),
\qquad (x,A,Y)\in\mathcal O.
\]
For an ordered dataset
\(D=((X_i,A_i,Y_i))_{i=1}^n\in\mathcal O^n\), define
\[
\mathscr D_n(D)
=E_\lambda\prod_{i=1}^n\mathscr L_\lambda(X_i,A_i,Y_i)
=\frac{dQ_n}{d(P_0^{\mathrm{obs}})^{\otimes n}}(D).
\]
These versions are bounded; in particular,
\(0\le\mathscr D_n\le8^n\).
Finite averaging and product integration therefore apply to all their
first and second moments.

Integrating \cref{eq:positive-family-fair-mark-identity} gives
\[
\int_{\mathcal O}\mathscr L_\lambda\mathscr L_{\lambda'}
                      \,dP_0^{\mathrm{obs}}
=1+H_{\lambda,\lambda'}\ge0.
\]
Since \(\int\mathscr D_n\,d(P_0^{\mathrm{obs}})^{\otimes n}=1\),
expanding the centered square, the finite mixture, and the product
integral gives
\[
\begin{aligned}
1+\chi^2\!\left(Q_n,(P_0^{\mathrm{obs}})^{\otimes n}\right)
&=\int_{\mathcal O^n}\mathscr D_n(D)^2\,
                         d(P_0^{\mathrm{obs}})^{\otimes n}(D)\\
&=E_{\lambda,\lambda'}
 \left(\int_{\mathcal O}
       \mathscr L_\lambda\mathscr L_{\lambda'}\,dP_0^{\mathrm{obs}}
 \right)^n\\
&=E_{\lambda,\lambda'}(1+H_{\lambda,\lambda'})^n.
\end{aligned}
\]
For each pair of encodings, the inequality \(1+H_{\lambda,\lambda'}
\le\exp(H_{\lambda,\lambda'})\), together with
\(1+H_{\lambda,\lambda'}\ge0\), implies
\[
(1+H_{\lambda,\lambda'})^n
\le\exp(nH_{\lambda,\lambda'}).
\]
The empty product and every zeroth power equal \(1\), so these
identities and inequalities include \(n=0\). Combining with the joint
exponential bound and subtracting one proves
\[
\chi^2\!\left(Q_n,(P_0^{\mathrm{obs}})^{\otimes n}\right)
\le
\exp\!\left\{160n^2t^2h_{\mathrm L}\delta_{\mathrm L}\right\}-1.
\]
% lean: CausalSmith.Stat.PrivateCateRoughdesign.signMixture_chiSq_overlap_identity
% lean: CausalSmith.Stat.PrivateCateRoughdesign.signDensityOverlap_eq_one_add_signOverlap
% lean: CausalSmith.Stat.PrivateCateRoughdesign.signMixture_chiSq_le_fourier_exp_average

\item \textbf{Regularity and alternative membership.}
We next verify the smoothness conditions using the two scales of the
construction. For \(\varrho>0\), \(\zeta\in\mathbb R\), and
\(x\in[0,1]\), define the localized cosine
\[
\mathcal C_{\varrho,\zeta}(x)=
\begin{cases}
\cos\!\left(\dfrac{\pi(x-\zeta)}{2\varrho}\right),
   &|x-\zeta|\le\varrho,\\
0,&|x-\zeta|>\varrho.
\end{cases}
\]
It has the clamped representation
\[
\mathcal C_{\varrho,\zeta}(x)
=\cos\!\left(
 \frac{\pi}{2\varrho}
 \max\{-\varrho,\min\{\varrho,x-\zeta\}\}\right).
\]
Indeed, beyond either endpoint the clamped cosine equals
\(\cos(\pm\pi/2)=0\). Clamping is \(1\)-Lipschitz and cosine is
\(1\)-Lipschitz, so
\[
|\mathcal C_{\varrho,\zeta}(x)
 -\mathcal C_{\varrho,\zeta}(x')|
\le\frac{\pi}{2\varrho}|x-x'|.
\]
Also, for two values of absolute value at most \(1\), factoring the
difference of their squares bounds that difference by twice their
distance. Since
\[
\phi_j=\mathcal C_{\delta_{\mathrm L},j\delta_{\mathrm L}},
\qquad
g=\mathcal C_{h_{\mathrm L},x_0}^{\,2},
\qquad q=g^2,
\]
we obtain
\[
\begin{aligned}
|\phi_j(x)-\phi_j(x')|
 &\le\frac{\pi}{2\delta_{\mathrm L}}|x-x'|,\\
|g(x)-g(x')|
 &\le\frac{\pi}{h_{\mathrm L}}|x-x'|,\\
|q(x)-q(x')|
 &\le\frac{2\pi}{h_{\mathrm L}}|x-x'|.
\end{aligned}
\]
% lean: CausalSmith.Stat.PrivateCateRoughdesign.localized_cosine_difference
% lean: CausalSmith.Stat.PrivateCateRoughdesign.bump_difference_bound

Only indices in \(\mathcal J_x\cup\mathcal J_{x'}\) contribute to
\(F_\lambda(x)-F_\lambda(x')\), and that union has at most six
elements. Hence
\[
\begin{aligned}
|F_\lambda(x)-F_\lambda(x')|
&\le\sum_{j\in\mathcal J_x\cup\mathcal J_{x'}}
                       |\phi_j(x)-\phi_j(x')|\\
&\le6\frac{\pi}{2\delta_{\mathrm L}}|x-x'|
=\frac{3\pi}{\delta_{\mathrm L}}|x-x'|.
\end{aligned}
\]
Using \(|F_\lambda(x')|\le3\) and the product identity
\[
S_\lambda(x)-S_\lambda(x')
=g(x)\bigl(F_\lambda(x)-F_\lambda(x')\bigr)
 +\bigl(g(x)-g(x')\bigr)F_\lambda(x'),
\]
we obtain, since \(\delta_{\mathrm L}\le h_{\mathrm L}\) and
\(\pi<4\),
\[
\begin{aligned}
|S_\lambda(x)-S_\lambda(x')|
&\le\left(\frac{3\pi}{\delta_{\mathrm L}}
                  +\frac{3\pi}{h_{\mathrm L}}\right)|x-x'|\\
&\le\frac{6\pi}{\delta_{\mathrm L}}|x-x'|
\le\frac{24}{\delta_{\mathrm L}}|x-x'|.
\end{aligned}
\]
Similarly,
\[
|q(x)-q(x')|
\le\frac{2\pi}{h_{\mathrm L}}|x-x'|
\le\frac{24}{\delta_{\mathrm L}}|x-x'|.
\]
The uniform bounds give
\[
|S_\lambda(x)-S_\lambda(x')|\le6\le24,
\qquad
|q(x)-q(x')|\le1\le24.
\]

If \(x=x'\), both differences vanish. If
\(0<|x-x'|/\delta_{\mathrm L}\le1\), use
\[
\frac{|x-x'|}{\delta_{\mathrm L}}
\le
\left(\frac{|x-x'|}{\delta_{\mathrm L}}\right)^{1/10}
\]
in the Lipschitz bounds. If
\(|x-x'|/\delta_{\mathrm L}>1\), use
\[
1\le
\left(\frac{|x-x'|}{\delta_{\mathrm L}}\right)^{1/10}
\]
in the uniform bounds. Thus, for every \(x,x'\in[0,1]\),
\[
\begin{aligned}
|S_\lambda(x)-S_\lambda(x')|
&\le24\frac{|x-x'|^{1/10}}{\delta_{\mathrm L}^{1/10}},\\
|q(x)-q(x')|
&\le24\frac{|x-x'|^{1/10}}{\delta_{\mathrm L}^{1/10}}.
\end{aligned}
\]
% lean: CausalSmith.Stat.PrivateCateRoughdesign.perturbation_holder_bound
% lean: CausalSmith.Stat.PrivateCateRoughdesign.bump_holder_bound

The scale calibration is
\[
\delta_{\mathrm L}^{1/10}=h_{\mathrm L}^{1/2},
\qquad
\sqrt t=2^{-6}h_{\mathrm L}^{1/2},
\qquad
\frac{\sqrt t}{\delta_{\mathrm L}^{1/10}}=\frac1{64},
\qquad
0\le t\le\sqrt t.
\]
The selected propensity and control mean in
\cref{obj:def:cosine-family} therefore satisfy
\[
\begin{aligned}
|e_{P_\lambda}(x)-e_{P_\lambda}(x')|
&=\frac{\sqrt t}{2}|S_\lambda(x)-S_\lambda(x')|\\
&\le\frac3{16}|x-x'|^{1/10}
\le L|x-x'|^{1/10},
\end{aligned}
\]
and
\[
\begin{aligned}
|\mu_{0,P_\lambda}(x)-\mu_{0,P_\lambda}(x')|
&\le\frac12\left(
 \sqrt t\,|S_\lambda(x)-S_\lambda(x')|
 +t\,|q(x)-q(x')|\right)\\
&\le\frac38|x-x'|^{1/10}
\le L|x-x'|^{1/10},
\end{aligned}
\]
where \cref{obj:synth_4} supplies \(L=3\).
Subtracting the two arm means cancels the shared perturbation:
\[
\tau_{P_\lambda}(x)=tq(x).
\]
Consequently,
\[
\begin{aligned}
|\tau_{P_\lambda}(x)-\tau_{P_\lambda}(x')|
&\le t\frac{2\pi}{h_{\mathrm L}}|x-x'|\\
&=2\pi\,2^{-12}|x-x'|
\le L|x-x'|.
\end{aligned}
\]
These establish
\cref{obj:ass:propensity-holder,obj:ass:control-holder,obj:ass:contrast-lipschitz}.
The amplitude bound also gives
\[
\frac{1-3/128}{2}\le e_{P_\lambda}(x)
\le\frac{1+3/128}{2},
\qquad
\frac14\le e_{P_\lambda}(x)\le\frac34,
\]
as required by \cref{obj:ass:overlap}.
The construction in \cref{obj:def:cosine-family} supplies consistency,
conditional exchangeability, and density \(f_{P_\lambda}=1\), which
satisfy
\cref{obj:ass:consistency,obj:ass:exchangeability,obj:ass:density}.
Every defining condition in \cref{obj:def:complete-model} is therefore
satisfied:
\[
P_\lambda\in\mathcal P\qquad
(\lambda\in\{0,1\}^{\mathcal J}).
\]
% lean: CausalSmith.Stat.PrivateCateRoughdesign.cosineFamily_completeModel

\item \textbf{The null, targets, and one-record cancellation.}
Under \(P_0\), consistency holds by construction, treatment is
conditionally independent of the potential outcomes, and the selected
functions satisfy
\[
f_{P_0}(x)=1,\qquad
e_{P_0}(x)=\mu_{0,P_0}(x)=\frac12,\qquad
\tau_{P_0}(x)=0
\quad(x\in[0,1]).
\]
The density and overlap bounds follow immediately; each smoothness
difference is zero. Thus \cref{obj:def:complete-model} gives
\[
P_0\in\mathcal P,\qquad \theta(P_0)=0.
\]
At the localization center,
\[
g(x_0)=\cos^2(0)=1,\qquad q(x_0)=1,
\qquad
\theta(P_\lambda)=\tau_{P_\lambda}(x_0)=t.
\]

The likelihood expansion established above is the asserted formula.
Its sign average is, pointwise in every covariate and mark pair,
\[
\begin{aligned}
E_\lambda L_\lambda(x,U,V_Y)
={}&1+U\sqrt t\,E_\lambda S_\lambda(x)
 +V_Y\sqrt t\,(tq(x)-1)E_\lambda S_\lambda(x)\\
&+UV_Yt\bigl(q(x)-E_\lambda S_\lambda(x)^2\bigr)
=1.
\end{aligned}
\]
Hence the uniformly averaged conditional probability of each observed
cell is \(1/4\). Integrating against the common uniform covariate law
identifies the measures:
\[
E_\lambda P_\lambda^{\mathrm{obs}}=P_0^{\mathrm{obs}}.
\]
The pointwise frame identities ensure this cancellation also at all
macro and frame boundaries.
% lean: CausalSmith.Stat.PrivateCateRoughdesign.fairNull_completeModel
% lean: CausalSmith.Stat.PrivateCateRoughdesign.fairNull_target
% lean: CausalSmith.Stat.PrivateCateRoughdesign.cosineFamily_target
% lean: CausalSmith.Stat.PrivateCateRoughdesign.oneRecordMixture_eq_fairNull

\item \textbf{Absolute continuity and square integrability.}
Taking finite products of the observed domination inequality and
then averaging gives
\[
(P_\lambda^{\mathrm{obs}})^{\otimes n}
\le8^n(P_0^{\mathrm{obs}})^{\otimes n},
\qquad
Q_n=E_\lambda(P_\lambda^{\mathrm{obs}})^{\otimes n}
\le8^n(P_0^{\mathrm{obs}})^{\otimes n}.
\]
Thus
\[
Q_n\ll(P_0^{\mathrm{obs}})^{\otimes n},
\qquad
0\le\frac{dQ_n}{d(P_0^{\mathrm{obs}})^{\otimes n}}
\le8^n
\quad (P_0^{\mathrm{obs}})^{\otimes n}\text{-almost everywhere}.
\]
The centered squared density is measurable and obeys
\[
0\le
\left(\frac{dQ_n}{d(P_0^{\mathrm{obs}})^{\otimes n}}-1\right)^2
\le(8^n+1)^2.
\]
Since the dominating law is a probability measure,
\[
\int_{\mathcal O^n}
\left(\frac{dQ_n}{d(P_0^{\mathrm{obs}})^{\otimes n}}-1\right)^2
\,d(P_0^{\mathrm{obs}})^{\otimes n}
\le(8^n+1)^2<\infty.
\]
For \(n=0\), both dataset laws are concentrated on the empty dataset
and the density equals \(1\), so the centered integral is zero.
This completes all the asserted certificates.
% lean: CausalSmith.Stat.PrivateCateRoughdesign.signMixture_absolutelyContinuous
% lean: CausalSmith.Stat.PrivateCateRoughdesign.signMixture_square_density_integrable
\end{enumerate}
\end{proof}

\Cref{obj:lem:positive-family-certificate} supplies separated targets within the same causal model and with a uniform covariate density. Its one-record identity explains the role of the shared signs: after averaging, each individual observed record has exactly the null distribution, whereas joint records can retain dependence through their common latent signs. The chi-square bound quantifies the resulting sample comparison through \(n^2t^2h_{\mathrm L}\delta_{\mathrm L}\).

The exponential comparison uses the bounded-differences log-moment inequality. For a bounded measurable function of independent random elements whose coordinate sensitivities are \(c_i\), this inequality bounds its centered exponential moment at \(u\ge0\) by \(\exp\{u^2\sum_i c_i^2/8\}\) \citep[Theorem 6.2, displayed log-mgf conclusion in its proof, p.~167]{BoucheronLugosiMassart2013}. This published concentration result supplies the exponential-moment control used in \cref{obj:lem:positive-family-certificate}.

To compare private output laws, we organize the same sign dependence conditional on the covariates. The macro envelope \leanref{sym:g}{\ensuremath{g(x)}}, micro frame functions \leanref{sym:phi}{\ensuremath{\phi_j(x)}}, and active-sign set \leanref{sym:Jsign}{\ensuremath{\mathcal J}} specify which records share a nonzero sign contribution. The shared-sign graph \leanref{sym:graph}{\ensuremath{G_X}} records these links.

\begin{definitionv}{synth_14}[Conditional shared-sign graph]
Let \(x_0=1/2\), \(0<h_{\mathrm L}\le1/4\), and \(\delta_{\mathrm L}=h_{\mathrm L}^5\). Define the macro envelope and cosine frame by
\[
g(x)=\cos^2\!\left\{\frac{\pi(x-x_0)}{2h_{\mathrm L}}\right\}
\mathbf1\{|x-x_0|\le h_{\mathrm L}\},
\qquad
\phi_j(x)=\cos\!\left\{\frac{\pi(x-j\delta_{\mathrm L})}{2\delta_{\mathrm L}}\right\}
\mathbf1\{|x-j\delta_{\mathrm L}|\le\delta_{\mathrm L}\},
\]
and define the finite active-sign index set
\[
\mathcal J=
\left\{j\in\mathbb Z:
(j\delta_{\mathrm L}-\delta_{\mathrm L},j\delta_{\mathrm L}+\delta_{\mathrm L})
\cap(x_0-h_{\mathrm L},x_0+h_{\mathrm L})\ne\varnothing
\right\}.
\]
For a covariate vector \(X=(X_1,\ldots,X_n)\), with \(n\ge0\), the conditional shared-sign graph \(G_X\) is the simple undirected graph with vertex set \(\{1,\ldots,n\}\). Distinct vertices \(i,l\) are adjacent precisely when
\[
\exists j\in\mathcal J:\quad
g(X_i)\phi_j(X_i)\ne0
\quad\text{and}\quad
g(X_l)\phi_j(X_l)\ne0.
\]
Its connected components are selected in lexicographic vertex order.
\end{definitionv}

A connected component \leanref{sym:component}{\ensuremath{C}} of the graph in \cref{obj:synth_14} groups records linked through shared signs; its size is \leanref{sym:m}{\ensuremath{m=|C|}}. Conditional on fixed covariates, let \leanref{sym:QC}{\ensuremath{Q_C}} denote the treatment--outcome mark law on that component under the mixture in \cref{obj:def:cosine-family}, and let \leanref{sym:PC}{\ensuremath{P_{0,C}}} denote the corresponding fair-null law. The latter is uniform over all component mark configurations. Their total common mass \leanref{sym:alpha}{\ensuremath{\alpha}} measures the probability mass that can be placed on identical configurations in both coordinates of a coupling.

\begin{definitionv}{synth_7}[Total common component mass]
The total common mass \(\alpha\), for fixed covariate values in \([0,1]\) and a finite vertex set \(C\), is defined by
\[
\alpha
=
\sum_{z:C\to\{0,1\}^2}
\min\bigl\{Q_C(\{z\}),\,P_{0,C}(\{z\})\bigr\},
\]
where \(z\) assigns a binary treatment–outcome pair to each vertex in \(C\), and \(Q_C\) and \(P_{0,C}\) are the alternative and null component mark laws conditional on the fixed covariate values.
\end{definitionv}

For the finite coupling formula, configurations are encoded by increasing vertex order. The alternative masses \leanref{sym:amass}{\ensuremath{a_v}}, null masses \leanref{sym:bmass}{\ensuremath{b_v}}, and common masses \leanref{sym:cmass}{\ensuremath{c_v}} then express the overlap in \cref{obj:synth_7} as arrays on a common finite space.

\begin{definitionv}{synth_15}[Component mark masses and common overlap]
Fix the covariates and a connected component \(C\) of \(G_X\), and write \(m=|C|\). Order its vertices increasingly and encode their treatment–outcome marks \((A_i,Y_i)_{i\in C}\) as elements of \(\{0,1\}^{2m}\). Let \(Q_C\) be their conditional law under the alternative dataset mixture
\[
Q_n=E_\lambda[(P_\lambda^{\mathrm{obs}})^{\otimes n}]
\]
from \(\cref{obj:def:cosine-family}\). Let \(P_{0,C}\) be their conditional law under \((P_0^{\mathrm{obs}})^{\otimes n}\), where the null observed law has \(X\sim\operatorname{Unif}[0,1]\) and, conditional on \(X\), independent fair Bernoulli treatment and outcome marks. Thus \(P_{0,C}\) is uniform on \(\{0,1\}^{2m}\). For every configuration \(v\in\{0,1\}^{2m}\), define
\[
a_v=Q_C(\{v\}),\qquad
b_v=P_{0,C}(\{v\}),\qquad
c_v=\min\{Q_C(\{v\}),P_{0,C}(\{v\})\}.
\]
In particular, \(a_z=Q_C(\{z\})\), \(b_w=P_{0,C}(\{w\})\), and \(c_z\) is the common diagonal mass at \(z\). Define the total common mass by
\[
\alpha=\sum_{v\in\{0,1\}^{2m}}c_v.
\]
\end{definitionv}

The component coupling \leanref{sym:Gamma}{\ensuremath{\Gamma_C}} assigns the common masses from \cref{obj:synth_15} to the diagonal and couples the remaining masses by a normalized product. Keeping the covariates identical in both datasets allows this construction to be combined across the ordered graph components.

\begin{definitionv}{def:common-mass-coupling}[Component coupling \(\Gamma_C\)]
For component mark configurations \(z,w\in\{0,1\}^{2m}\), define
\[
\Gamma_C(z,w)
=
\begin{cases}
\displaystyle
c_z\mathbf 1\{z=w\}
+\frac{(a_z-c_z)(b_w-c_w)}{1-\alpha},
&\alpha<1,\\[6pt]
c_z\mathbf 1\{z=w\},
&\alpha=1.
\end{cases}
\]
The dataset coupling uses the same covariates in both coordinates and the product of these kernels across the lexicographically selected graph components.
\end{definitionv}

The comparison below establishes that \cref{obj:def:common-mass-coupling} defines a measurable dataset coupling with the required marginals. It also bounds the component total variation distance \leanref{sym:TV}{\ensuremath{d_{\mathrm{TV}}(Q_C,P_{0,C})}}. Under the sparse-occupancy condition \(n\delta_{\mathrm L}\le1/128\), the resulting Hamming transportation cost \leanref{sym:WH}{\ensuremath{W_{\mathrm H}}}, the infimum of expected dataset Hamming distance over couplings, yields a comparison of output laws for every private release covered by \cref{obj:def:private-kernels}.

\begin{lemmav}{lem:measurable-component-contraction}[Measurable component coupling and contraction]*
Let \(n\) be a nonnegative integer. Impose the following conditions:
\begin{itemize}
\item \textbf{(Lower-family calibration.)} The macro localization radius \(h_{\mathrm L}\), micro localization radius \(\delta_{\mathrm L}\), and target-separation amplitude \(t\) satisfy
\[
0<h_{\mathrm L}\le\frac14,
\qquad
\delta_{\mathrm L}=h_{\mathrm L}^{5},
\qquad
t=2^{-12}h_{\mathrm L}.
\]
\item \textbf{(Output space.)} Let \(\mathcal B\) be a standard Borel output space.
\item \textbf{(Private release.)} Let \(0<\epsilon\le1\) and let \(M\in\mathfrak M_{n,\epsilon}(\mathcal B)\), as defined in \cref{obj:def:private-kernels}.
\end{itemize}
Let \(Q_n\) be the uniform sign mixture of the alternative observed product laws in \cref{obj:def:cosine-family}, and let \((P_0^{\mathrm{obs}})^{\otimes n}\) be the fair-null dataset law.

For every covariate vector \(X=x\), let \(G_X\) be the shared-sign graph, whose connected components group records linked by shared latent signs. For each connected component \(C\), set \(m=|C|\). Its conditional alternative and null mark laws \(Q_C\) and \(P_{0,C}\) have masses
\[
Q_C(\{z\})
 = E_\lambda\prod_{i\in C}
   \frac{L_\lambda(x_i,U_i,V_{Y,i})}{4},
\qquad
P_{0,C}(\{z\})=4^{-m},
\]
where \(z=((U_i,V_{Y,i}))_{i\in C}\) is a configuration of the binary treatment and outcome marks, \(L_\lambda\) is their conditional likelihood relative to the fair null, and \(E_\lambda\) averages uniformly over the latent sign vectors.

For every \(x\) and every full mark configuration \(z=((U_i,V_{Y,i}))_{i=1}^n\), the conditional masses factor over the components of \(G_X\):
\[
E_\lambda\prod_{i=1}^n
  \frac{L_\lambda(x_i,U_i,V_{Y,i})}{4}
 = \prod_{C\text{ component of }G_X} Q_C(\{z|_C\}),
\qquad
4^{-n}
 = \prod_{C\text{ component of }G_X} P_{0,C}(\{z|_C\}).
\]
For every such component,
\[
m=1\ \Longrightarrow\ Q_C=P_{0,C},
\qquad
d_{\mathrm{TV}}(Q_C,P_{0,C})\le4tm^2,
\]
where \(d_{\mathrm{TV}}\) denotes total variation distance. The common-mass measure \(\Gamma_C\) of \cref{obj:def:common-mass-coupling} is a coupling of \(Q_C\) and \(P_{0,C}\). The product of these component couplings, with components ordered by increasing least record index, defines a measurable conditional dataset coupling. Integrating it over the shared covariate law, the product of \(n\) uniform laws on \([0,1]\), gives a dataset coupling with marginals \(Q_n\) and \((P_0^{\mathrm{obs}})^{\otimes n}\).

If \(n\delta_{\mathrm L}\le1/128\), then
\[
W_{\mathrm H}\!\left(Q_n,(P_0^{\mathrm{obs}})^{\otimes n}\right)
 \le1024tn^2h_{\mathrm L}\delta_{\mathrm L},
\qquad
d_{\mathrm{TV}}\!\left(MQ_n,M(P_0^{\mathrm{obs}})^{\otimes n}\right)
 \le2048\epsilon tn^2h_{\mathrm L}\delta_{\mathrm L},
\]
where \(W_{\mathrm H}\) is the infimum of expected dataset Hamming distance over couplings, Hamming distance counts the record positions at which two datasets differ, and \(MQ\) denotes the output law obtained by applying \(M\) to a dataset with law \(Q\).
\end{lemmav}
% CAUSALSMITH-CITED-SCOPE-BEGIN lem:measurable-component-contraction
\verificationfootnotetext{\textbf{Formalization scope.} This result uses the cited conclusion \citep{KellerTrotterAppliedCombinatorics} (Section 5.6, Theorem 5.40) and \citep{BoucheronLugosiMassart2013} (Theorem 6.2, displayed log-mgf conclusion in its proof, printed page 167). The cited source proof is not formalized here: Lean verifies its use and all remaining steps. If that cited conclusion is false, the portions of this result that depend on it are not certified.}
% CAUSALSMITH-CITED-SCOPE-END lem:measurable-component-contraction

\begin{proof}[Proof of \cref{obj:lem:measurable-component-contraction}]
\begin{enumerate}
\item \textit{Conditional factorization.}
Fix a covariate vector \(x\in[0,1]^n\). To identify the signs used by each record, write the envelope and micro frames in \cref{obj:def:cosine-family} as
\[
g_{\mathrm L}(v)=
\begin{cases}
\cos^2\!\left(\dfrac{\pi(v-x_0)}{2h_{\mathrm L}}\right),
& |v-x_0|\le h_{\mathrm L},\\
0,& |v-x_0|>h_{\mathrm L},
\end{cases}
\qquad
\phi_{\mathrm L,\eta}(v)=
\begin{cases}
\cos\!\left(\dfrac{\pi(v-\eta\delta_{\mathrm L})}{2\delta_{\mathrm L}}\right),
& |v-\eta\delta_{\mathrm L}|\le\delta_{\mathrm L},\\
0,& |v-\eta\delta_{\mathrm L}|>\delta_{\mathrm L},
\end{cases}
\]
for \(v\in[0,1]\) and \(\eta\in\mathbb Z\), where \(x_0=1/2\). Thus
\[
q(v)=g_{\mathrm L}(v)^2,
\qquad
S_\lambda(v)=g_{\mathrm L}(v)
 \sum_{\eta\in\mathcal J_{\mathrm L}}(2\lambda_\eta-1)\phi_{\mathrm L,\eta}(v),
\]
with the finite active index set
\[
\mathcal J_{\mathrm L}
=
\left\{
\eta\in\mathbb Z:
\begin{array}{l}
\displaystyle
\left\lfloor\frac{x_0-h_{\mathrm L}}{\delta_{\mathrm L}}\right\rfloor-1
\le \eta\le
\left\lceil\frac{x_0+h_{\mathrm L}}{\delta_{\mathrm L}}\right\rceil+1,\\
\eta\delta_{\mathrm L}-\delta_{\mathrm L}<x_0+h_{\mathrm L},
\quad
x_0-h_{\mathrm L}<\eta\delta_{\mathrm L}+\delta_{\mathrm L}
\end{array}
\right\}.
\]
Define the record and component sign sets by
\[
J_i(x)=
 \{\eta\in\mathcal J_{\mathrm L}:
   g_{\mathrm L}(x_i)\phi_{\mathrm L,\eta}(x_i)\ne0\},
\qquad
J_C(x)=\bigcup_{i\in C}J_i(x).
\]
The graph and its ordered component list are therefore specified by
\[
i\sim_{G_x}i'
\quad\Longleftrightarrow\quad
i\ne i'\ \text{ and }\ J_i(x)\cap J_{i'}(x)\ne\varnothing,
\]
\[
\mathscr C(x)
=
\bigl(C:\ C\text{ is a connected component of }G_x\bigr),
\qquad
\text{listed by increasing }\min C.
\]
For \(n=0\), this list is empty.

We use binary mark configurations and their signed likelihood coordinates through the encoding
\[
z_i=(A_i,Y_i)\in\{0,1\}^2,
\qquad
U_i(z)=2A_i-1,
\qquad
V_{Y,i}(z)=2Y_i-1.
\]
This identifies binary configurations with the signed configurations in the statement. The likelihood expansion supplied by
\cref{obj:lem:positive-family-certificate} shows that the conditional likelihood at record \(i\) depends on \(\lambda\) only through \(J_i(x)\). Its quadratic term uses the same sign set.

Distinct components have disjoint sign sets:
\[
C,C'\in\mathscr C(x),\ C\ne C'
\quad\Longrightarrow\quad
J_C(x)\cap J_{C'}(x)=\varnothing.
\]
Indeed, a common sign would be used by a vertex in each component and would join them by an edge. Since the signs are independent and fair, summing separately over these disjoint sign blocks gives
\[
E_\lambda\prod_{i=1}^n
 \frac{L_\lambda(x_i,U_i(z),V_{Y,i}(z))}{4}
=
\prod_{C\in\mathscr C(x)}
 E_\lambda\prod_{i\in C}
 \frac{L_\lambda(x_i,U_i(z),V_{Y,i}(z))}{4}
=
\prod_{C\in\mathscr C(x)}Q_C(\{z|_C\}).
\]
Unused signs average to one. The component vertex sets partition the record indices, so
\[
4^{-n}
=
\prod_{C\in\mathscr C(x)}4^{-|C|}
=
\prod_{C\in\mathscr C(x)}P_{0,C}(\{z|_C\}).
\]
Both identities also hold for \(n=0\), with empty products equal to one.
% lean: CausalSmith.Stat.PrivateCateRoughdesign.fullAlternativeMass_component_factorization
% lean: CausalSmith.Stat.PrivateCateRoughdesign.nullMass_component_factorization

\item \textit{Singleton cancellation.}
We first establish the pointwise sign identities needed below. Define, for \(v\in[0,1]\),
\[
\eta_v=\left\lfloor\frac{v}{\delta_{\mathrm L}}\right\rfloor,
\qquad
\omega_{\mathrm L}(v)
=\frac{\pi(v-\eta_v\delta_{\mathrm L})}{2\delta_{\mathrm L}}.
\]
Then \(0\le\omega_{\mathrm L}(v)<\pi/2\), and the only possibly nonzero frames are the two neighboring ones:
\[
\phi_{\mathrm L,\eta_v}(v)=\cos\omega_{\mathrm L}(v),
\qquad
\phi_{\mathrm L,\eta_v+1}(v)=\sin\omega_{\mathrm L}(v),
\qquad
\phi_{\mathrm L,\eta}(v)=0
\quad(\eta\notin\{\eta_v,\eta_v+1\}).
\]
Consequently their squared values sum to one. When \(q(v)>0\), every nonzero frame at \(v\) belongs to the active set; when \(q(v)=0\), multiplying by \(q(v)\) annihilates every term. Hence, including support boundaries,
\[
q(v)\sum_{\eta\in\mathcal J_{\mathrm L}}
 \phi_{\mathrm L,\eta}(v)^2=q(v).
\]
Fair independent signs satisfy
\[
E_\lambda(2\lambda_\eta-1)=0,
\qquad
E_\lambda\bigl[(2\lambda_\eta-1)(2\lambda_{\eta'}-1)\bigr]
=\mathbf 1\{\eta=\eta'\}.
\]
Expanding the perturbation and its square therefore yields
\[
E_\lambda S_\lambda(v)=0,
\qquad
E_\lambda S_\lambda(v)^2
=
q(v)\sum_{\eta\in\mathcal J_{\mathrm L}}\phi_{\mathrm L,\eta}(v)^2
=q(v).
\]
Substitution into the likelihood expansion in
\cref{obj:lem:positive-family-certificate} gives the pointwise identity
\[
E_\lambda L_\lambda(v,U,V_Y)=1
\qquad
(v\in[0,1],\ U,V_Y\in\{-1,1\}).
\]
Thus, if \(C=\{i\}\), every binary mark pair has mass
\[
Q_C(\{z\})
=\frac14E_\lambda L_\lambda(x_i,U_i(z),V_{Y,i}(z))
=\frac14
=P_{0,C}(\{z\}),
\]
and \(Q_C=P_{0,C}\).
% lean: CausalSmith.Stat.PrivateCateRoughdesign.conditionalLikelihood_sign_sum_exact
% lean: CausalSmith.Stat.PrivateCateRoughdesign.alternativeComponentLaw_eq_null_of_card_one

\item \textit{Component curvature.}
Fix a component \(C\), and put \(m=|C|\). The two-frame property and Cauchy--Schwarz imply, for every \(\lambda\) and \(v\),
\[
S_\lambda(v)^2
\le
2q(v)\sum_{j\in\mathcal J_{\mathrm L}}\phi_{\mathrm L,j}(v)^2
=2q(v).
\]
Since \(0\le q(v)\le1\), this gives the coefficient bounds
\[
|S_\lambda(v)|\le\frac32,
\qquad
|q(v)-S_\lambda(v)^2|\le q(v)\le1.
\]

For \(i\in C\), \(\lambda\), \(z_i\in\{0,1\}^2\), and every real amplitude \(u\), define the polynomial mark mass
\[
\begin{aligned}
p_{i,\lambda}(z_i;u)
=\frac14\bigl[
1
&+U_i(z)uS_\lambda(x_i)\\
&+V_{Y,i}(z)u(u^2q(x_i)-1)S_\lambda(x_i)\\
&+U_i(z)V_{Y,i}(z)u^2
   \bigl(q(x_i)-S_\lambda(x_i)^2\bigr)
\bigr].
\end{aligned}
\]
Its factorization is
\[
p_{i,\lambda}(z_i;u)
=
\frac{1+U_i(z)uS_\lambda(x_i)}2
\frac{1+V_{Y,i}(z)
 \bigl(-uS_\lambda(x_i)+U_i(z)u^2q(x_i)\bigr)}2.
\]
For \(0\le u\le1/10\), the two perturbations satisfy
\[
|uS_\lambda(x_i)|\le\frac3{20},
\qquad
\bigl|-uS_\lambda(x_i)+U_i(z)u^2q(x_i)\bigr|
\le\frac3{20}+\frac1{100}
=\frac4{25}.
\]
The masses are therefore positive. Summing first over the outcome mark and then over the treatment mark gives
\[
\sum_{z_i\in\{0,1\}^2}p_{i,\lambda}(z_i;u)=1,
\qquad
\sum_{z_i\in\{0,1\}^2}|p_{i,\lambda}(z_i;u)|=1.
\]

Differentiation gives
\[
\begin{aligned}
\frac{\partial p_{i,\lambda}}{\partial u}(z_i;u)
=\frac14\bigl[
&U_i(z)S_\lambda(x_i)
+V_{Y,i}(z)(3u^2q(x_i)-1)S_\lambda(x_i)\\
&+2U_i(z)V_{Y,i}(z)u
 \bigl(q(x_i)-S_\lambda(x_i)^2\bigr)
\bigr],
\end{aligned}
\]
\[
\frac{\partial^2p_{i,\lambda}}{\partial u^2}(z_i;u)
=\frac14\bigl[
6V_{Y,i}(z)u q(x_i)S_\lambda(x_i)
+2U_i(z)V_{Y,i}(z)
 \bigl(q(x_i)-S_\lambda(x_i)^2\bigr)
\bigr].
\]
On \(0\le u\le1/10\), we have
\[
|3u^2q(x_i)-1|\le1.
\]
Taking absolute values of the displayed Fourier coefficients and summing over the four mark pairs consequently gives
\[
\begin{aligned}
\sum_{z_i}
 \left|\frac{\partial p_{i,\lambda}}{\partial u}(z_i;u)\right|
&\le
|S_\lambda(x_i)|
+|3u^2q(x_i)-1|\,|S_\lambda(x_i)|\\
&\qquad
+2u|q(x_i)-S_\lambda(x_i)^2|\\
&\le \frac32+\frac32+\frac15
\le4,
\end{aligned}
\]
and
\[
\begin{aligned}
\sum_{z_i}
 \left|\frac{\partial^2p_{i,\lambda}}{\partial u^2}(z_i;u)\right|
&\le
6u q(x_i)|S_\lambda(x_i)|
+2|q(x_i)-S_\lambda(x_i)^2|\\
&\le\frac9{10}+2
\le4.
\end{aligned}
\]

Define the averaged component array, for every \(u\in\mathbb R\), by
\[
F_{C,z}(u)=E_\lambda\prod_{i\in C}p_{i,\lambda}(z_i;u),
\qquad
z\in(\{0,1\}^2)^C.
\]
The second product derivative consists of one term for each twice-differentiated factor and one term for each ordered pair of distinct once-differentiated factors:
\[
\begin{aligned}
\frac{d^2}{du^2}\prod_{i\in C}p_{i,\lambda}(z_i;u)
={}&
\sum_{i\in C}
 \frac{\partial^2p_{i,\lambda}}{\partial u^2}(z_i;u)
 \prod_{i'\in C\setminus\{i\}}p_{i',\lambda}(z_{i'};u)\\
&+
\sum_{\substack{i,i'\in C\\i\ne i'}}
 \frac{\partial p_{i,\lambda}}{\partial u}(z_i;u)
 \frac{\partial p_{i',\lambda}}{\partial u}(z_{i'};u)
 \prod_{i''\in C\setminus\{i,i'\}}p_{i'',\lambda}(z_{i''};u).
\end{aligned}
\]
For finite arrays, summing the absolute value of a tensor product factors into the product of the absolute sums. Applying the preceding bounds and then averaging over signs yields
\[
\sum_z|F_{C,z}''(u)|
\le4m+16m(m-1)
\le16m^2
\qquad(0\le u\le1/10).
\]
At zero amplitude all factors equal \(1/4\), while
\[
\frac{\partial p_{i,\lambda}}{\partial u}(z_i;0)
=
\frac{U_i(z)-V_{Y,i}(z)}4S_\lambda(x_i).
\]
The pointwise sign cancellation from the preceding step thus gives
\[
F_{C,z}(0)=4^{-m},
\qquad
F_{C,z}'(0)=0.
\]
For an empty vertex set, the product is identically one and both derivatives vanish, so the same conclusions and curvature bound apply.
% lean: CausalSmith.Stat.PrivateCateRoughdesign.amplitudeMark_coefficients_valid
% lean: CausalSmith.Stat.PrivateCateRoughdesign.amplitudeMarkSlope_abs_sum_le_four
% lean: CausalSmith.Stat.PrivateCateRoughdesign.amplitudeMarkCurvature_abs_sum_le_four
% lean: CausalSmith.Stat.PrivateCateRoughdesign.componentAmplitudeCurvature_abs_sum_le
% lean: CausalSmith.Stat.PrivateCateRoughdesign.componentAmplitudeSlope_zero

\item \textit{Quadratic remainder and common overlap.}
The calibration implies
\[
0<t\le\frac1{16384},
\qquad
0\le\sqrt t\le\frac1{10}.
\]
At this amplitude, the likelihood expansion in
\cref{obj:lem:positive-family-certificate} identifies the component arrays as
\[
a_z:=Q_C(\{z\})=F_{C,z}(\sqrt t),
\qquad
b_z:=P_{0,C}(\{z\})=F_{C,z}(0)=4^{-m}.
\]
To control their entire absolute difference with one scalar remainder, define
\[
\sigma_C(z)=
\begin{cases}
1,&a_z-b_z\ge0,\\
-1,&a_z-b_z<0,
\end{cases}
\qquad
\mathcal F_C(u)=\sum_z\sigma_C(z)F_{C,z}(u)
\quad(u\in\mathbb R).
\]
Then
\[
\mathcal F_C'(0)=0,
\qquad
|\mathcal F_C''(u)|
\le\sum_z|F_{C,z}''(u)|
\le16m^2
\quad(0\le u\le\sqrt t).
\]
The mean-value inequality gives
\[
|\mathcal F_C'(u)|\le16m^2u
\qquad(0\le u\le\sqrt t).
\]
Integrating the first derivative, and using the defining signs, proves
\[
\begin{aligned}
\sum_z|a_z-b_z|
&=\mathcal F_C(\sqrt t)-\mathcal F_C(0)\\
&\le\int_0^{\sqrt t}|\mathcal F_C'(u)|\,du\\
&\le\int_0^{\sqrt t}16m^2u\,du
=8tm^2.
\end{aligned}
\]

Define the common entries and total overlap by
\[
c_z:=\min(a_z,b_z),
\qquad
\alpha:=\sum_zc_z.
\]
Both arrays are nonnegative and sum to one, by the mark-mass normalization above and the \(4^m\) null configurations. Therefore
\[
0\le\alpha\le1,
\qquad
\sum_z(a_z-c_z)=\sum_z(b_z-c_z)=1-\alpha,
\]
and
\[
\sum_z|a_z-b_z|
=\sum_z(a_z+b_z-2c_z)
=2(1-\alpha).
\]
For every mark event
\(\mathcal E_C\subseteq(\{0,1\}^2)^C\), each residual sum lies in \([0,1-\alpha]\), so
\[
\left|
 \sum_{z\in\mathcal E_C}a_z-\sum_{z\in\mathcal E_C}b_z
\right|
=
\left|
 \sum_{z\in\mathcal E_C}(a_z-c_z)
 -
 \sum_{z\in\mathcal E_C}(b_z-c_z)
\right|
\le1-\alpha.
\]
Taking the supremum over events gives
\[
d_{\mathrm{TV}}(Q_C,P_{0,C})
\le1-\alpha
=\frac12\sum_z|a_z-b_z|
\le4tm^2.
\]
In fact, the event
\[
\mathcal E_C^+:=\{z:b_z\le a_z\}
\]
has probability difference
\[
Q_C(\mathcal E_C^+)-P_{0,C}(\mathcal E_C^+)
=\sum_z(a_z-c_z)=1-\alpha.
\]
Thus the total variation distance equals \(1-\alpha\), as will also describe the coupling's failure mass.
% lean: CausalSmith.Stat.PrivateCateRoughdesign.finite_array_second_derivative_remainder_le
% lean: CausalSmith.Stat.PrivateCateRoughdesign.alternativeMass_abs_sub_nullMass_sum_le
% lean: CausalSmith.Stat.PrivateCateRoughdesign.overlapMass_deficit_le
% lean: CausalSmith.Stat.PrivateCateRoughdesign.component_TV_le_quadratic
% lean: CausalSmith.Stat.PrivateCateRoughdesign.component_TV_eq_overlap_deficit

\item \textit{Component coupling.}
Use the masses \(\Gamma_C(z,w)\) defined in
\cref{obj:def:common-mass-coupling}, with the arrays just specified. When \(\alpha<1\), summing its formula over the second coordinate gives
\[
\sum_w\Gamma_C(z,w)
=
c_z+\frac{(a_z-c_z)\sum_w(b_w-c_w)}{1-\alpha}
=a_z.
\]
Similarly,
\[
\sum_z\Gamma_C(z,w)
=
c_w+\frac{\sum_z(a_z-c_z)(b_w-c_w)}{1-\alpha}
=b_w.
\]
When \(\alpha=1\), all nonnegative residual entries have zero sum and hence vanish individually:
\[
a_z=c_z=b_z
\qquad\text{for every }z.
\]
The diagonal branch then has the same row and column sums. In both cases the masses are nonnegative and sum to one, establishing that \(\Gamma_C\) couples \(Q_C\) and \(P_{0,C}\).

At every configuration, at least one residual entry vanishes:
\[
(a_z-c_z)(b_z-c_z)=0.
\]
Consequently the diagonal and off-diagonal masses are
\[
\Gamma_C(z,z)=c_z,
\qquad
\sum_{z\ne w}\Gamma_C(z,w)=1-\alpha.
\]
For a singleton component, the cancellation already proved gives \(\alpha=1\).
% lean: CausalSmith.Stat.PrivateCateRoughdesign.componentCoupling_isCoupling
% lean: CausalSmith.Stat.PrivateCateRoughdesign.componentCouplingMass_diagonal
% lean: CausalSmith.Stat.PrivateCateRoughdesign.overlapMass_eq_one_of_card_one

\item \textit{Measurable dataset coupling and its marginals.}
For full binary mark vectors \(z,w\in(\{0,1\}^2)^n\), define
\[
p_{\mathrm{joint}}(x,z,w)
:=
\prod_{C\in\mathscr C(x)}
 \Gamma_C(z|_C,w|_C),
\qquad
D_{\mathrm{attach}}(x,z)
:=
\bigl((x_i,A_i,Y_i)\bigr)_{i=1}^n,
\]
where \(z_i=(A_i,Y_i)\). The conditional dataset measure is
\[
\gamma_x
:=
\sum_{z,w}p_{\mathrm{joint}}(x,z,w)\,
\delta_{\left(D_{\mathrm{attach}}(x,z),
               D_{\mathrm{attach}}(x,w)\right)}.
\]

Each envelope and frame is continuous, including at its support endpoints. The edge predicates defining \(G_x\) are finite unions of Borel nonzero predicates. For every fixed labeled graph \(\mathsf G\) on the record indices, the set
\[
\mathcal X_{\mathsf G}:=\{x\in[0,1]^n:G_x=\mathsf G\}
\]
is therefore Borel. There are finitely many such graphs. On each \(\mathcal X_{\mathsf G}\), the ordered component list is fixed, and the component masses are finite sums of products of Borel functions. Taking minima, finite sums, and the two overlap branches preserves measurability; the denominator in the branch \(\alpha<1\) is positive. It follows that \(x\mapsto p_{\mathrm{joint}}(x,z,w)\) is Borel. The attachment maps are Borel as well, so the finite atomic formula makes \(x\mapsto\gamma_x\) measurable.

The component marginal identities and factorization give, for every \(x,z,w\),
\[
\begin{aligned}
\sum_w p_{\mathrm{joint}}(x,z,w)
&=\prod_{C\in\mathscr C(x)}Q_C(\{z|_C\})\\
&=E_\lambda\prod_{i=1}^n
 \frac{L_\lambda(x_i,U_i(z),V_{Y,i}(z))}{4},
\end{aligned}
\qquad
\sum_zp_{\mathrm{joint}}(x,z,w)=4^{-n}.
\]
In particular, \(\gamma_x\) is a probability measure.

Define the shared covariate law and integrated coupling by
\[
\nu_X:=\bigotimes_{i=1}^n\operatorname{Unif}[0,1],
\qquad
\gamma(\mathcal E)
:=
\int_{[0,1]^n}\gamma_x(\mathcal E)\,\nu_X(dx)
\]
for every measurable dataset-pair event \(\mathcal E\).
The alternatives and null have this common covariate law. Integrating the first conditional marginal and interchanging the finite sign average with the integral yields the uniform mixture of alternative product laws; integrating the second yields the fair-null product law. Explicitly, with
\[
\operatorname{pr}_1(D,D')=D,
\qquad
\operatorname{pr}_2(D,D')=D',
\]
we obtain
\[
(\operatorname{pr}_1)_\#\gamma=Q_n,
\qquad
(\operatorname{pr}_2)_\#\gamma=(P_0^{\mathrm{obs}})^{\otimes n},
\qquad
\gamma(\mathcal O^n\times\mathcal O^n)=1.
\]
For \(n=0\), the mark space and dataset space each have one element; the empty component product gives the unit mass on the pair of empty datasets.
% lean: CausalSmith.Stat.PrivateCateRoughdesign.measurable_conditionalDatasetCoupling
% lean: CausalSmith.Stat.PrivateCateRoughdesign.commonMassCoupling_isCoupling

\item \textit{Component cost and tree probabilities.}
Assume now that \(n\delta_{\mathrm L}\le1/128\), and first consider \(n\ge1\). For component configurations define
\[
d_{\mathrm H,C}(z,w)
:=\sum_{i\in C}\mathbf 1\{z_i\ne w_i\},
\qquad
c_{\mathrm H,C}(x)
:=\sum_{z,w}d_{\mathrm H,C}(z,w)\Gamma_C(z,w).
\]
Because \(d_{\mathrm H,C}\) is zero on the diagonal and at most \(m\) elsewhere,
\[
c_{\mathrm H,C}(x)
\le m(1-\alpha)
\le4tm^3.
\]
For \(m=1\), the coupling is diagonal and its cost is zero. The shared covariates ensure that dataset Hamming distance is exactly the sum of component mark distances. Thus
\[
\begin{aligned}
\int d_{\mathrm H}(D,D')\,\gamma(dD,dD')
&=\int\sum_{C\in\mathscr C(x)}c_{\mathrm H,C}(x)\,\nu_X(dx)\\
&\le4t\sum_{m=2}^n m^3\int N_m(x)\,\nu_X(dx),
\end{aligned}
\]
where, for every integer \(m\ge0\), the component count is
\[
N_m(x)
:=\#\{C\in\mathscr C(x):|C|=m\}.
\]
In particular \(N_0(x)=0\), and \(N_m(x)=0\) for \(m>n\).

To bound these counts, a shared-sign edge has the geometric restrictions
\[
i\sim_{G_x}i'
\quad\Longrightarrow\quad
|x_i-x_0|<h_{\mathrm L},
\quad
|x_{i'}-x_0|<h_{\mathrm L},
\quad
|x_i-x_{i'}|<2\delta_{\mathrm L}.
\]
Indeed, the common nonzero weighted frame places both covariates in the open macro support and within \(\delta_{\mathrm L}\) of the same frame center. The last inequality follows by the triangle inequality.

Fix \(2\le m\le n\), an \(m\)-element record set \(C\), and a bijection
\[
e_C:\{1,\ldots,m\}\longrightarrow C.
\]
Let
\[
\mathscr T_m
:=\{\text{undirected trees on the labeled vertex set }\{1,\ldots,m\}\}.
\]
For \(\mathsf T\in\mathscr T_m\), define its shared-edge event by
\[
\mathcal E_{\mathsf T,C}
:=
\left\{
x:
e_C(i)\sim_{G_x}e_C(i')
\ \text{for every edge }\{i,i'\}\text{ of }\mathsf T
\right\}.
\]
A connected tree admits a parent-before-child order: start at a root and repeatedly adjoin a vertex adjacent to the vertices already selected. Connectivity permits this until all vertices are selected. Write this order and its parent positions as
\[
\rho_{\mathsf T}:\{1,\ldots,m\}\longrightarrow\{1,\ldots,m\}
\quad\text{a permutation},
\qquad
\pi_{\mathsf T}(i)\in\{1,\ldots,i-1\}
\quad(2\le i\le m),
\]
with
\[
\{\rho_{\mathsf T}(i),\rho_{\mathsf T}(\pi_{\mathsf T}(i))\}
\text{ an edge of }\mathsf T.
\]
Since \(m\ge2\), the first edge localizes the root. The shared-edge event is contained in the region
\[
\begin{split}
\mathcal R_{\mathsf T,C}:=\bigl\{x:\ &
|x_{e_C(\rho_{\mathsf T}(1))}-x_0|<h_{\mathrm L},\\
&
|x_{e_C(\rho_{\mathsf T}(i))}
 -x_{e_C(\rho_{\mathsf T}(\pi_{\mathsf T}(i)))}
 |<2\delta_{\mathrm L}
\quad(2\le i\le m)\bigr\}.
\end{split}
\]
The root interval has length at most \(2h_{\mathrm L}\), and, conditional on the preceding coordinates, each child interval has length at most \(4\delta_{\mathrm L}\). Integrating the last child first and repeating gives
\[
\nu_X(\mathcal E_{\mathsf T,C})
\le\nu_X(\mathcal R_{\mathsf T,C})
\le2h_{\mathrm L}(4\delta_{\mathrm L})^{m-1}.
\]
Coordinates outside \(C\) integrate to one.

If \(C\) is a graph component, its connected induced graph contains a spanning tree. Hence
\[
\{x:C\in\mathscr C(x)\}
\subseteq
\bigcup_{\mathsf T\in\mathscr T_m}\mathcal E_{\mathsf T,C}.
\]
The labeled-tree count is
\[
|\mathscr T_m|=m^{m-2},
\]
by \citep[Section 5.6, Theorem 5.40]{KellerTrotterAppliedCombinatorics}. A union bound, followed by summation over the \(\binom nm\) record subsets, proves
\[
\int N_m(x)\,\nu_X(dx)
\le
\binom nm\,2h_{\mathrm L}m^{m-2}
 (4\delta_{\mathrm L})^{m-1}.
\]
All geometric restrictions used here are strict nonzero-support restrictions, so the estimate includes every support-boundary configuration.
% lean: CausalSmith.Stat.PrivateCateRoughdesign.commonMassCoupling_hamming_cost_le_cubic_sum
% lean: CausalSmith.Stat.PrivateCateRoughdesign.uniform_labeled_tree_probability_bound
% lean: CausalSmith.Stat.PrivateCateRoughdesign.lintegral_sizeComponentCount_le_of_uniform_labeled_tree_bound

\item \textit{Summation of the transport envelope.}
Using \(\binom nm\le n^m/m!\) in the preceding cost bound gives
\[
\begin{aligned}
\int d_{\mathrm H}(D,D')\,\gamma(dD,dD')
&\le
\sum_{m=2}^n
4tm^3\binom nm\,2h_{\mathrm L}m^{m-2}
 (4\delta_{\mathrm L})^{m-1}\\
&\le
32tn^2h_{\mathrm L}\delta_{\mathrm L}
\sum_{m=2}^n
\frac{m^{m+1}}{m!}(4n\delta_{\mathrm L})^{m-2}.
\end{aligned}
\]
For \(n=1\), both sums are empty.

The nonnegative exponential series gives
\[
\frac{m^m}{m!}
\le
\sum_{j_{\mathrm{exp}}=0}^{\infty}
 \frac{m^{j_{\mathrm{exp}}}}{j_{\mathrm{exp}}!}
=\exp(m)
=\exp(1)^m.
\]
Define the series argument
\[
\zeta_{\mathrm{tree}}
:=4\exp(1)n\delta_{\mathrm L}.
\]
The sparsity cap and \(\exp(1)<3\) imply
\[
0\le\zeta_{\mathrm{tree}}
\le\frac3{32}
\le\frac18.
\]
Termwise,
\[
\frac{m^{m+1}}{m!}(4n\delta_{\mathrm L})^{m-2}
\le
\exp(1)^2\,m\,\zeta_{\mathrm{tree}}^{m-2}.
\]
The differentiated finite geometric sum yields, for \(n\ge1\),
\[
\begin{aligned}
\sum_{m=2}^n m\zeta_{\mathrm{tree}}^{m-2}
&=
\frac{
2-\zeta_{\mathrm{tree}}
-(n+1)\zeta_{\mathrm{tree}}^{n-1}
+n\zeta_{\mathrm{tree}}^n
}{(1-\zeta_{\mathrm{tree}})^2}\\
&\le
\frac{2-\zeta_{\mathrm{tree}}}
 {(1-\zeta_{\mathrm{tree}})^2}
\le3.
\end{aligned}
\]
For the first inequality, the removed tail equals
\(\zeta_{\mathrm{tree}}^{n-1}
 ((n+1)-n\zeta_{\mathrm{tree}})\ge0\).
For the second, \(1-\zeta_{\mathrm{tree}}\ge7/8\), and
\[
2-\zeta_{\mathrm{tree}}
\le2
\le3(7/8)^2
\le3(1-\zeta_{\mathrm{tree}})^2.
\]
Consequently,
\[
\sum_{m=2}^n
\frac{m^{m+1}}{m!}(4n\delta_{\mathrm L})^{m-2}
\le3\exp(1)^2
\le32.
\]
Substitution proves the required bound for the constructed coupling:
\[
\int d_{\mathrm H}(D,D')\,\gamma(dD,dD')
\le1024tn^2h_{\mathrm L}\delta_{\mathrm L}.
\]
% lean: CausalSmith.Stat.PrivateCateRoughdesign.tree_series_term_le
% lean: CausalSmith.Stat.PrivateCateRoughdesign.tree_series_sum_le_thirty_two
% lean: CausalSmith.Stat.PrivateCateRoughdesign.cubic_tree_envelope_sum_le
% lean: CausalSmith.Stat.PrivateCateRoughdesign.commonMassCoupling_hamming_cost_le_sparse

\item \textit{Transport and private contraction.}
For the remaining case \(n=0\), both datasets are empty, and
\[
d_{\mathrm H}(D,D')=0,
\qquad
\int d_{\mathrm H}(D,D')\,\gamma(dD,dD')
=0
=1024tn^2h_{\mathrm L}\delta_{\mathrm L}.
\]
Thus the coupling-cost bound holds over the full range of \(n\).

Since \(\gamma\) has the required marginals, the defining infimum of Hamming transport gives
\[
W_{\mathrm H}\!\left(Q_n,(P_0^{\mathrm{obs}})^{\otimes n}\right)
\le
\int d_{\mathrm H}(D,D')\,\gamma(dD,dD')
\le1024tn^2h_{\mathrm L}\delta_{\mathrm L}.
\]
Its marginals are probability laws, and the hypotheses on \(\epsilon\), the output space, and \(M\) permit application of
\cref{obj:lem:private-kernel-comparison}. That comparison, together with the defining infimum, supplies the coupling-cost inequality
\[
d_{\mathrm{TV}}\!\left(MQ_n,M(P_0^{\mathrm{obs}})^{\otimes n}\right)
\le
2\epsilon
\int d_{\mathrm H}(D,D')\,\gamma(dD,dD').
\]
Using the established cost bound and \(2\epsilon\ge0\), we conclude
\[
\begin{aligned}
d_{\mathrm{TV}}\!\left(MQ_n,M(P_0^{\mathrm{obs}})^{\otimes n}\right)
&\le
2\epsilon\bigl(1024tn^2h_{\mathrm L}\delta_{\mathrm L}\bigr)\\
&=
2048\epsilon tn^2h_{\mathrm L}\delta_{\mathrm L}.
\end{aligned}
\]
% lean: CausalSmith.Stat.PrivateCateRoughdesign.private_TV_le_coupling_cost
% lean: CausalSmith.Stat.PrivateCateRoughdesign.measurable_component_contraction
\end{enumerate}
\end{proof}

The factorization in \cref{obj:lem:measurable-component-contraction} makes the graph components the units of the conditional comparison. Singleton components have identical alternative and null mark laws, and the common-mass construction couples those marks identically. For larger components, the stated total variation bound controls their discrepancy. The measurable product construction turns these finite comparisons into a coupling of the full dataset laws.

The component-size control uses the labelled-tree enumeration: on a specified set of \(s\ge2\) vertex labels, there are exactly \(s^{s-2}\) undirected trees \citep[Section 5.6, Theorem 5.40]{KellerTrotterAppliedCombinatorics}. The bounded-differences exponential-moment inequality stated above also enters this comparison \citep[Theorem 6.2, displayed log-mgf conclusion in its proof, p.~167]{BoucheronLugosiMassart2013}. These published tools support the quantitative component analysis in \cref{obj:lem:measurable-component-contraction}.

Under \(n\delta_{\mathrm L}\le1/128\), and with the calibration and privacy conditions of \cref{obj:lem:measurable-component-contraction}, the private output comparison is controlled by \(\epsilon tn^2h_{\mathrm L}\delta_{\mathrm L}\). Its quantification over standard Borel output spaces accommodates both scalar releases and the interval-valued releases of \cref{obj:def:interval-criterion}. Thus the same sparse-occupancy comparison supplies a testing ingredient for estimation and honest interval length.

A complementary experiment for regime selection perturbs the treatment-arm mean directly while keeping the propensity and control regression constant. The direct localized alternative \leanref{sym:P1}{\ensuremath{P_1}} uses the same amplitude and contrast profile, as specified next.

\begin{definitionv}{def:direct-family}[Direct localized causal alternative]
The causal law \(P_1\), the direct localized alternative, is defined for the macro localization radius \(0<h_{\mathrm L}\le 1/4\), with target-separation amplitude and localized contrast profile
\[
t=2^{-12}h_{\mathrm L},
\qquad
q(x)=
\begin{cases}
\cos^4\!\left(\dfrac{\pi(x-x_0)}{2h_{\mathrm L}}\right),
& |x-x_0|\le h_{\mathrm L},\\
0,& |x-x_0|>h_{\mathrm L},
\end{cases}
\quad x\in[0,1],
\]
where \(x_0\) is the target covariate location. Under \(P_1\), draw the covariate \(X\sim\operatorname{Unif}[0,1]\). Conditional on \(X\), draw the treatment indicator \(A\) and the potential outcomes \(Y(0)\) and \(Y(1)\) independently according to
\[
A\sim\operatorname{Bernoulli}(1/2),
\qquad
Y(0)\sim\operatorname{Bernoulli}(1/2),
\qquad
Y(1)\sim\operatorname{Bernoulli}\bigl(1/2+tq(X)\bigr).
\]
Set the observed outcome \(Y\) to
\[
Y=Y(A).
\]
The selected covariate density, propensity, control regression, and treatment-effect contrast are
\[
f_{P_1}(x)=1,\qquad
e_{P_1}(x)=\frac12,\qquad
\mu_{0,P_1}(x)=\frac12,\qquad
\tau_{P_1}(x)=tq(x).
\]
\end{definitionv}

At the target location, the profile in \cref{obj:def:direct-family} equals one, so the direct alternative has target value \(t\). Its constant propensity and control regression isolate a localized change in the treatment-arm mean. Together with the shared-sign family in \cref{obj:def:cosine-family}, it supplies the complementary experiment used in selecting the testing comparison across regimes.


\section{Testing reductions and matching lower bounds}

The experiments in \cref{obj:def:cosine-family,obj:def:direct-family} provide separated point targets within the causal model. This section collects the probability comparisons that make those experiments useful for both absolute-error risk and honest connected-interval length. The common ingredient is a bound on the distinguishability of their released outputs, uniform over the private procedures admitted by \cref{obj:def:private-kernels}.

\subsection{Dataset laws, released outputs, and testing priors}

Let \(Q\) and \(Q'\) be probability laws on the ordered-dataset space \(\mathcal O^n\). Their total variation distance \(d_{\mathrm{TV}}(Q,Q')\) is the supremum of their probability differences over measurable events. Their Hamming transportation distance \(W_{\mathrm H}(Q,Q')\) is the infimum, over couplings of the two dataset laws, of the expected number of differing record positions. When \(Q\ll Q'\), the chi-square divergence \(\chi^2(Q,Q')\) is the squared deviation of the likelihood ratio from one, integrated under \(Q'\). For a randomized release \(M\), the induced output law \(MQ\) integrates the conditional release probabilities against the dataset law \(Q\).

A finite testing prior is a uniform distribution over a nonempty finite indexed family of causal laws. For \(a\in\{0,1\}\), write this family as \((P_{a,j})_{j=1}^{N_a}\), where \(N_a\) is a positive integer, and define its averaged observed-sample law by
\[
Q^{(a)}
=
\frac{1}{N_a}\sum_{j=1}^{N_a}
(P_{a,j}^{\mathrm{obs}})^{\otimes n}.
\]
The testing construction uses families satisfying
\[
\theta(P_{0,j})=0,
\qquad
\theta(P_{1,j})=\Delta,
\qquad
\Delta>0,
\]
where \(\Delta\) is their common target separation. Averaging takes place over product laws, so each member retains the independent-sampling interpretation of \cref{obj:ass:iid}.

The following comparison connects dataset discrepancies to the distribution of a released statistic. Its privacy bound uses Hamming transportation, while its ordinary probability bounds apply through total variation and chi-square divergence.

\begin{lemmav}{lem:private-kernel-comparison}[Private kernel comparison bounds]
Let \(n\) be a nonnegative integer and let \(Q,Q'\) be probability laws on \(\mathcal O^n\), the space of ordered datasets of \(n\) observed records. Suppose that:
\begin{itemize}
\item \textbf{(Privacy parameter.)} \(0<\epsilon\le 1\).
\item \textbf{(Output space.)} \(\mathcal B\) is a standard Borel output space.
\item \textbf{(Private release.)} \(M\in\mathfrak M_{n,\epsilon}(\mathcal B)\), the class of everywhere probability-valued private Markov kernels defined in \cref{obj:def:private-kernels}.
\end{itemize}
Write \(MQ\) for the induced output law, so that
\[
(MQ)(E)=\int M(D,E)\,Q(dD),
\qquad E\in\mathcal B(\mathcal B).
\]
For probability laws on a common measurable space, let \(d_{\mathrm{TV}}\) denote total variation distance:
\[
d_{\mathrm{TV}}(Q,Q')=\sup_{E\ \mathrm{measurable}}|Q(E)-Q'(E)|.
\]
Let \(W_{\mathrm H}\) denote the infimum of expected Hamming distance over all joint laws with the specified marginals:
\[
W_{\mathrm H}(Q,Q')
=
\inf_{\substack{\mathcal L(D)=Q\\\mathcal L(D')=Q'}}
E\!\left[d_{\mathrm H}(D,D')\right],
\qquad
d_{\mathrm H}(D,D')
=
\sum_{i=1}^{n}\mathbf 1\{D_i\ne D'_i\}.
\]
Then
\[
d_{\mathrm{TV}}(MQ,MQ')
\le 2\epsilon W_{\mathrm H}(Q,Q').
\]
Furthermore, for every everywhere probability-valued Markov kernel \(M\) from \(\mathcal O^n\) to \(\mathcal B\),
\[
d_{\mathrm{TV}}(MQ,MQ')\le d_{\mathrm{TV}}(Q,Q').
\]
If \(Q\ll Q'\) and the squared Radon--Nikodym deviation is integrable, that is,
\[
\int_{\mathcal O^n}
\left(\frac{dQ}{dQ'}(D)-1\right)^2\,Q'(dD)<\infty,
\]
then, with chi-square divergence defined by
\[
\chi^2(Q,Q')
=
\int_{\mathcal O^n}
\left(\frac{dQ}{dQ'}(D)-1\right)^2\,Q'(dD),
\]
one has
\[
d_{\mathrm{TV}}(Q,Q')
\le \frac12\sqrt{\chi^2(Q,Q')}.
\]
\end{lemmav}

\begin{proof}[Proof of \cref{obj:lem:private-kernel-comparison}]
\begin{enumerate}
\item For adjacent datasets \(D,D'\) and a measurable output event \(E\), \cref{obj:def:private-kernels,obj:ass:pure-privacy} give
\[
M(D,E)\le e^\epsilon M(D',E),
\qquad
M(D',E)\le e^\epsilon M(D,E).
\]
Since the event probabilities belong to \([0,1]\),
\[
\begin{aligned}
M(D,E)-M(D',E)
&\le (e^\epsilon-1)M(D',E)\le e^\epsilon-1,\\
M(D',E)-M(D,E)
&\le (e^\epsilon-1)M(D,E)\le e^\epsilon-1.
\end{aligned}
\]
Convexity of the exponential, together with \(0<\epsilon\le1\) and \(e<3\), yields
\[
e^\epsilon
\le (1-\epsilon)e^0+\epsilon e^1,
\qquad
e^\epsilon-1\le (e-1)\epsilon\le2\epsilon.
\]
Consequently,
\[
|M(D,E)-M(D',E)|\le2\epsilon
\qquad\text{if }d_{\mathrm H}(D,D')=1.
\]
% lean: CausalSmith.Stat.PrivateCateRoughdesign.exp_privacy_increment_le
% lean: CausalSmith.Stat.PrivateCateRoughdesign.private_adjacent_event_bound

\item Fix a measurable output event \(E\), and define
\[
\varphi_E:\mathcal O^n\longrightarrow[0,1],
\qquad
\varphi_E(D)=M(D,E).
\]
We prove by induction on a finite set of coordinate positions that
\[
\mathcal J\subseteq\{1,\ldots,n\},
\quad
D_i=D'_i\ \text{for all }i\notin\mathcal J
\quad\Longrightarrow\quad
|\varphi_E(D)-\varphi_E(D')|
\le2\epsilon|\mathcal J|.
\]
For \(\mathcal J=\varnothing\), the datasets coincide and both sides of the inequality are zero.

For the induction step, suppose \(D,D'\) agree outside
\(\mathcal J\cup\{\iota\}\), where \(\iota\notin\mathcal J\). Define the intermediate dataset by
\[
(D_{\mathrm{mid}})_{\iota'}
=
\begin{cases}
D'_\iota,&\iota'=\iota,\\
D_{\iota'},&\iota'\ne\iota,
\end{cases}
\qquad \iota'\in\{1,\ldots,n\}.
\]
Then \(D_{\mathrm{mid}}\) and \(D'\) agree outside \(\mathcal J\), so the induction hypothesis gives
\[
|\varphi_E(D_{\mathrm{mid}})-\varphi_E(D')|
\le2\epsilon|\mathcal J|.
\]
If \(D_\iota=D'_\iota\), then \(D_{\mathrm{mid}}=D\) and their event-probability difference is zero. Otherwise, their Hamming distance is one, and the adjacent bound applies. In either case,
\[
|\varphi_E(D)-\varphi_E(D_{\mathrm{mid}})|\le2\epsilon.
\]
The triangle inequality therefore gives
\[
\begin{aligned}
|\varphi_E(D)-\varphi_E(D')|
&\le
|\varphi_E(D)-\varphi_E(D_{\mathrm{mid}})|
+
|\varphi_E(D_{\mathrm{mid}})-\varphi_E(D')|\\
&\le2\epsilon+2\epsilon|\mathcal J|
=2\epsilon|\mathcal J\cup\{\iota\}|.
\end{aligned}
\]
Apply this conclusion to
\[
\mathcal J=\{i\in\{1,\ldots,n\}:D_i\ne D'_i\}.
\]
It follows that, for every pair of datasets,
\[
|M(D,E)-M(D',E)|
\le2\epsilon d_{\mathrm H}(D,D').
\]
The empty-set case also covers \(n=0\).
% lean: CausalSmith.Stat.PrivateCateRoughdesign.event_bound_of_hamming_path
% lean: CausalSmith.Stat.PrivateCateRoughdesign.private_hamming_event_bound

\item Let \(\gamma\) be any coupling of \(Q,Q'\); explicitly,
\[
\gamma\in\operatorname{Prob}(\mathcal O^n\times\mathcal O^n),
\qquad
\gamma(F\times\mathcal O^n)=Q(F),
\qquad
\gamma(\mathcal O^n\times F)=Q'(F)
\]
for every measurable dataset event \(F\). Since Hamming distance is bounded by the dataset size,
\[
0\le
\int d_{\mathrm H}(D,D')\,\gamma(dD,dD')
\le
\int n\,\gamma(dD,dD')
=n.
\]
Thus the coupling cost is finite. The measurable function \(\varphi_E\) is bounded by one, so its integrals against the marginals and against \(\gamma\) are well defined. The marginal identities give
\[
(MQ)(E)-(MQ')(E)
=
\int\bigl(\varphi_E(D)-\varphi_E(D')\bigr)\,
\gamma(dD,dD').
\]
Using the integral triangle inequality and the preceding event bound,
\[
\begin{aligned}
|(MQ)(E)-(MQ')(E)|
&\le
\int|\varphi_E(D)-\varphi_E(D')|\,\gamma(dD,dD')\\
&\le
\int2\epsilon d_{\mathrm H}(D,D')\,\gamma(dD,dD')\\
&=
2\epsilon\int d_{\mathrm H}(D,D')\,\gamma(dD,dD').
\end{aligned}
\]
Taking the supremum over measurable output events proves
\[
d_{\mathrm{TV}}(MQ,MQ')
\le
2\epsilon\int d_{\mathrm H}(D,D')\,\gamma(dD,dD').
\]
% lean: CausalSmith.Stat.PrivateCateRoughdesign.private_TV_le_coupling_cost
This holds for every coupling. Since \(2\epsilon\) is positive and finite, taking the infimum gives
\[
d_{\mathrm{TV}}(MQ,MQ')
\le
\inf_{\gamma\ \mathrm{couples}\ Q,Q'}
2\epsilon\int d_{\mathrm H}(D,D')\,d\gamma
=
2\epsilon W_{\mathrm H}(Q,Q').
\]
% lean: CausalSmith.Stat.PrivateCateRoughdesign.private_kernel_comparison

\item Let
\[
K:\mathcal O^n\rightsquigarrow\mathcal B
\]
be an everywhere probability-valued Markov kernel. For a measurable output event \(E\), define
\[
\varphi_{K,E}:\mathcal O^n\longrightarrow[0,1],
\qquad
\varphi_{K,E}(D)=K(D,E).
\]
This function is measurable and bounded. Its layer representation under each of the probability laws is
\[
\begin{aligned}
\int\varphi_{K,E}(D)\,Q(dD)
&=\int_{(0,1]}Q\{D:\eta_{\mathrm{layer}}\le\varphi_{K,E}(D)\}\,d\eta_{\mathrm{layer}},\\
\int\varphi_{K,E}(D)\,Q'(dD)
&=\int_{(0,1]}Q'\{D:\eta_{\mathrm{layer}}\le\varphi_{K,E}(D)\}\,d\eta_{\mathrm{layer}}.
\end{aligned}
\]
These identities follow by integrating
\(\varphi_{K,E}(D)=\int_{(0,1]}\mathbf 1\{\eta_{\mathrm{layer}}\le\varphi_{K,E}(D)\}\,d\eta_{\mathrm{layer}}\).
Every displayed level set is measurable, and hence
\[
\left|
Q\{D:\eta_{\mathrm{layer}}\le\varphi_{K,E}(D)\}
-
Q'\{D:\eta_{\mathrm{layer}}\le\varphi_{K,E}(D)\}
\right|
\le d_{\mathrm{TV}}(Q,Q').
\]
Subtracting the layer representations and bounding the integral therefore yields
\[
\begin{aligned}
|(KQ)(E)-(KQ')(E)|
&=
\left|
\int_{(0,1]}
\left[
Q\{D:\eta_{\mathrm{layer}}\le\varphi_{K,E}(D)\}
-
Q'\{D:\eta_{\mathrm{layer}}\le\varphi_{K,E}(D)\}
\right]d\eta_{\mathrm{layer}}
\right|\\
&\le \int_{(0,1]}d_{\mathrm{TV}}(Q,Q')\,d\eta_{\mathrm{layer}}
=d_{\mathrm{TV}}(Q,Q').
\end{aligned}
\]
Taking the supremum over \(E\) proves
\[
d_{\mathrm{TV}}(KQ,KQ')\le d_{\mathrm{TV}}(Q,Q').
\]
% lean: Causalean.Stat.tvDist_integral_range
% lean: CausalSmith.Stat.PrivateCateRoughdesign.kernel_TV_contraction

\item Assume \(Q\ll Q'\) and the stated squared-deviation integral is finite. Choose an extended nonnegative Radon--Nikodym derivative and define, on all of \(\mathcal O^n\),
\[
p_{\mathrm{RN}}(D)=
\begin{cases}
\dfrac{dQ}{dQ'}(D),&\dfrac{dQ}{dQ'}(D)<\infty,\\
0,&\dfrac{dQ}{dQ'}(D)=\infty,
\end{cases}
\qquad
g_{\mathrm{RN}}(D)=p_{\mathrm{RN}}(D)-1.
\]
The infinite-density set is \(Q'\)-null, so this convention preserves the Radon--Nikodym integral identities. Finiteness of the nonnegative squared-deviation integral gives
\[
g_{\mathrm{RN}}\in L^2(Q'),
\qquad
\int g_{\mathrm{RN}}(D)^2\,Q'(dD)=\chi^2(Q,Q').
\]
Both \(|g_{\mathrm{RN}}|\) and the constant function one belong to \(L^2(Q')\). Hölder's inequality with exponents two and two gives
\[
\begin{aligned}
\int|g_{\mathrm{RN}}(D)|\,Q'(dD)
&=\int|g_{\mathrm{RN}}(D)|\cdot1\,Q'(dD)\\
&\le
\left(\int|g_{\mathrm{RN}}(D)|^2\,Q'(dD)\right)^{1/2}
\left(\int1^2\,Q'(dD)\right)^{1/2}\\
&=\sqrt{\chi^2(Q,Q')}.
\end{aligned}
\]
% lean: CausalSmith.Stat.PrivateCateRoughdesign.TV_le_sqrt_chiSq_of_finite
% lean: Causalean.Stat.tvDist_le_half_sqrt_chiSqDiv

To relate this integral to total variation, the probability normalization and the Radon--Nikodym identity give
\[
\int g_{\mathrm{RN}}\,dQ'
=
\int p_{\mathrm{RN}}\,dQ'-1
=Q(\mathcal O^n)-1=0.
\]
For every measurable dataset event \(F\),
\[
Q(F)-Q'(F)=\int_F g_{\mathrm{RN}}\,dQ',
\qquad
\int_{F^c}g_{\mathrm{RN}}\,dQ'
=-\int_F g_{\mathrm{RN}}\,dQ'.
\]
The pointwise inequalities
\(g_{\mathrm{RN}}\le|g_{\mathrm{RN}}|\) and
\(-g_{\mathrm{RN}}\le|g_{\mathrm{RN}}|\), integrated over \(F\) and \(F^c\), imply
\[
\left|\int_F g_{\mathrm{RN}}\,dQ'\right|
\le\int_F|g_{\mathrm{RN}}|\,dQ',
\qquad
\left|\int_F g_{\mathrm{RN}}\,dQ'\right|
=
\left|\int_{F^c}g_{\mathrm{RN}}\,dQ'\right|
\le\int_{F^c}|g_{\mathrm{RN}}|\,dQ'.
\]
Adding these bounds yields
\[
2|Q(F)-Q'(F)|
\le
\int_F|g_{\mathrm{RN}}|\,dQ'
+\int_{F^c}|g_{\mathrm{RN}}|\,dQ'
=
\int|g_{\mathrm{RN}}|\,dQ'.
\]
Taking the supremum over \(F\) and using the square-integrability bound proves
\[
d_{\mathrm{TV}}(Q,Q')
\le\frac12\int|g_{\mathrm{RN}}|\,dQ'
\le\frac12\sqrt{\chi^2(Q,Q')}.
\]
% lean: Causalean.Stat.measureReal_sub_eq_setIntegral_rnDeriv_sub_one
% lean: Causalean.Stat.abs_setIntegral_le_half_integral_abs_of_integral_eq_zero
% lean: Causalean.Stat.tvDist_le_half_integral_abs_rnDeriv
% lean: Causalean.Stat.tvDist_le_half_sqrt_chiSqDiv
% lean: CausalSmith.Stat.PrivateCateRoughdesign.private_kernel_comparison
\end{enumerate}
\end{proof}

The first bound makes a coupling of datasets useful for every private release with the stated output space. The second expresses contraction under a randomized procedure, and the third converts a square-integrable likelihood comparison into total variation control. Thus \cref{obj:lem:private-kernel-comparison} accommodates both the sign-mixture likelihood bound in \cref{obj:lem:positive-family-certificate} and the measurable coupling in \cref{obj:lem:measurable-component-contraction}.

\subsection{Finite priors across the budget regimes}

For \(n\ge2\) and \(0<\epsilon\le1\), use the numerical benchmark \(r(n,\epsilon)\) of \cref{obj:thm:matched-risk-frontier} and choose the macro localization radius \(h_{\mathrm L}=r(n,\epsilon)/4\). The calibration in \cref{obj:def:cosine-family,obj:synth_5} then gives
\[
\delta_{\mathrm L}=h_{\mathrm L}^{5},
\qquad
t=2^{-12}h_{\mathrm L}
=2^{-14}r(n,\epsilon).
\]
The common target value of the sign alternatives is certified by \cref{obj:lem:positive-family-certificate}; the direct contrast is specified in \cref{obj:def:direct-family}.

When the sampling scale dominates, the sign-mixture chi-square bound and \cref{obj:lem:private-kernel-comparison} supply the output comparison. When the sparse privacy scale strictly exceeds both competing scales, the calibrated radius satisfies the sparse-occupancy restriction \(n\delta_{\mathrm L}\le1/128\), allowing the private comparison in \cref{obj:lem:measurable-component-contraction}. The direct alternative supplies the comparison in the direct privacy and saturated branches. The exact selection conditions, including ties, are recorded in the certificate below.

The sign-mixture exponential comparison uses the bounded-differences log-moment conclusion of \citet[Theorem 6.2, displayed log-mgf conclusion in its proof, p.~167]{BoucheronLugosiMassart2013}. The sparse component comparison also uses the enumeration of \(s^{s-2}\) undirected trees on a specified set of \(s\ge2\) vertex labels \citep[Section 5.6, Theorem 5.40]{KellerTrotterAppliedCombinatorics}. These published conclusions enter through \cref{obj:lem:positive-family-certificate,obj:lem:measurable-component-contraction} and support the common testing comparison stated next.

\begin{lemmav}{lem:frontier-testing-priors}[Finite testing priors]*
For every integer \(n\ge2\) and privacy budget \(0<\epsilon\le1\), define the numerical benchmark
\[
r(n,\epsilon)
=
\min\left\{1,\max\left\{n^{-1/4},
\max\left\{(n^2\epsilon)^{-1/7},(n\epsilon)^{-1/2}\right\}\right\}\right\}.
\]
There exist two nonempty finite families of laws in \(\mathcal P\), the model class of \cref{obj:def:complete-model}, whose uniformly averaged observed \(n\)-record product laws are \(Q^{(0)}\) and \(Q^{(1)}\), respectively, and a common target separation \(\Delta\) satisfying
\[
\Delta\ge2^{-14}r(n,\epsilon).
\]
Every member of the first family has point target \(\theta(P)=0\), and every member of the second has \(\theta(P)=\Delta\).

These families can be chosen with macro localization radius
\[
h_{\mathrm L}=\frac{r(n,\epsilon)}4,
\qquad 0<h_{\mathrm L}\le\frac14,
\]
and the following properties:
\begin{itemize}
\item \textbf{(Null family.)} Every member of the first family equals \(P_0\), the causal law with uniform covariate \(X\), conditionally independent fair binary treatment and potential outcomes given \(X\), and observed outcome \(Y=Y(A)\). Consequently,
\[
Q^{(0)}=(P_0^{\mathrm{obs}})^{\otimes n}.
\]
\item \textbf{(Sign family.)} If
\[
1<n\epsilon
\quad\text{and}\quad
\left[
\left((n^2\epsilon)^{-1/7}\le n^{-1/4}
\ \text{and}\ (n\epsilon)^{-1/2}\le n^{-1/4}\right)
\ \text{or}\
\left(n^{-1/4}<(n^2\epsilon)^{-1/7}
\ \text{and}\ (n\epsilon)^{-1/2}<(n^2\epsilon)^{-1/7}\right)
\right],
\]
every member of the second family is \(P_\lambda\) for some sign vector \(\lambda\), using the construction of \cref{obj:def:cosine-family} at radius \(h_{\mathrm L}\), and
\[
Q^{(1)}=Q_n,
\]
where \(Q_n\) is the uniform sign mixture at that radius defined in \cref{obj:def:cosine-family}.
\item \textbf{(Direct family.)} If the preceding condition fails, every member of the second family equals \(P_1\), the direct alternative of \cref{obj:def:direct-family} at radius \(h_{\mathrm L}\), and
\[
Q^{(1)}=(P_1^{\mathrm{obs}})^{\otimes n}.
\]
\end{itemize}
For every standard Borel output space \(\mathcal B\) and every private Markov kernel \(M\in\mathfrak M_{n,\epsilon}(\mathcal B)\) as defined in \cref{obj:def:private-kernels}, the induced output laws satisfy
\[
d_{\mathrm{TV}}\!\left(MQ^{(0)},MQ^{(1)}\right)\le\frac14,
\]
where \(d_{\mathrm{TV}}\) denotes total variation distance and \(MQ^{(0)},MQ^{(1)}\) are obtained by applying \(M\) to the respective dataset laws.
\end{lemmav}
% CAUSALSMITH-CITED-SCOPE-BEGIN lem:frontier-testing-priors
\verificationfootnotetext{\textbf{Formalization scope.} This result uses the cited conclusion \citep{BoucheronLugosiMassart2013} (Theorem 6.2, displayed log-mgf conclusion in its proof, printed page 167) and \citep{KellerTrotterAppliedCombinatorics} (Section 5.6, Theorem 5.40). The cited source proof is not formalized here: Lean verifies its use and all remaining steps. If that cited conclusion is false, the portions of this result that depend on it are not certified.}
% CAUSALSMITH-CITED-SCOPE-END lem:frontier-testing-priors

\begin{proof}[Proof of \cref{obj:lem:frontier-testing-priors}]
\begin{enumerate}
\item Define the three candidate scales and the separation constant by
\[
a_n:=n^{-1/4},\qquad
b_n:=(n^2\epsilon)^{-1/7},\qquad
c_n:=(n\epsilon)^{-1/2},\qquad
\kappa:=2^{-12}.
\]
All three scales are positive. Set
\[
h_{\mathrm L}:=\frac{r(n,\epsilon)}4,\qquad
t:=\kappa h_{\mathrm L},\qquad
\delta_{\mathrm L}:=h_{\mathrm L}^{5},\qquad
\Delta:=t.
\]
Since \(0<r(n,\epsilon)\le1\),
\[
0<h_{\mathrm L}\le\frac14,
\qquad
\Delta=\frac{\kappa}{4}r(n,\epsilon)
      =2^{-14}r(n,\epsilon).
\]
We now split according to the sign-family condition in the statement.
% lean: CausalSmith.Stat.PrivateCateRoughdesign.frontier_testing_priors

\item Suppose that
\[
1<n\epsilon
\quad\text{and}\quad
\bigl[(b_n\le a_n\ \text{and}\ c_n\le a_n)
\ \text{or}\
(a_n<b_n\ \text{and}\ c_n<b_n)\bigr].
\]
At radius \(h_{\mathrm L}\), let the active index set and its binary sign vectors be
\[
\begin{aligned}
\mathcal J_{h_{\mathrm L}}
:=\biggl\{j\in\mathbb Z:\;&
\left\lfloor\frac{x_0-h_{\mathrm L}}{\delta_{\mathrm L}}\right\rfloor-1
\le j\le
\left\lceil\frac{x_0+h_{\mathrm L}}{\delta_{\mathrm L}}\right\rceil+1,\\
&j\delta_{\mathrm L}-\delta_{\mathrm L}<x_0+h_{\mathrm L},
\quad
x_0-h_{\mathrm L}<j\delta_{\mathrm L}+\delta_{\mathrm L}
\biggr\},
\qquad x_0=\frac12,\\
\Lambda_{h_{\mathrm L}}&:=\{\lambda:\mathcal J_{h_{\mathrm L}}\to\{0,1\}\}.
\end{aligned}
\]
The integer interval in this definition is bounded, so \(\mathcal J_{h_{\mathrm L}}\) is finite. Consequently,
\[
|\Lambda_{h_{\mathrm L}}|=2^{|\mathcal J_{h_{\mathrm L}}|}>0.
\]
This also covers an empty active set, for which there is one empty sign vector.

Choose the singleton null family \((P_0)\) and the alternative family
\((P_\lambda)_{\lambda\in\Lambda_{h_{\mathrm L}}}\) from
\cref{obj:def:cosine-family}. By
\cref{obj:lem:positive-family-certificate},
\[
P_0\in\mathcal P,\qquad \theta(P_0)=0,\qquad
P_\lambda\in\mathcal P,\qquad \theta(P_\lambda)=t=\Delta
\quad(\lambda\in\Lambda_{h_{\mathrm L}}).
\]
Their uniform product-law averages are
\[
\begin{aligned}
Q^{(0)}&=(P_0^{\mathrm{obs}})^{\otimes n},\\
Q^{(1)}
&=\frac1{|\Lambda_{h_{\mathrm L}}|}
  \sum_{\lambda\in\Lambda_{h_{\mathrm L}}}
  (P_\lambda^{\mathrm{obs}})^{\otimes n}\\
&=2^{-|\mathcal J_{h_{\mathrm L}}|}
  \sum_{\lambda\in\Lambda_{h_{\mathrm L}}}
  (P_\lambda^{\mathrm{obs}})^{\otimes n}
 =Q_n.
\end{aligned}
\]
The equality follows by enumerating every binary sign vector exactly once.
% lean: CausalSmith.Stat.PrivateCateRoughdesign.finiteProductMixture_signs

\item First consider the sampling subcase
\[
b_n\le a_n,\qquad c_n\le a_n.
\]
Since \(n\ge2\), we have \(a_n\le1\), and hence
\[
r(n,\epsilon)=a_n,\qquad
h_{\mathrm L}=\frac{a_n}{4}.
\]
The power identity \(a_n^8=n^{-2}\) gives
\[
n^2h_{\mathrm L}^{8}=\frac1{4^8}.
\]
Define the nonnegative exponent
\[
\eta_{\mathrm{samp}}
:=160n^2t^2h_{\mathrm L}\delta_{\mathrm L}
 =160\kappa^2n^2h_{\mathrm L}^{8}.
\]
It satisfies
\[
0\le\eta_{\mathrm{samp}}
 =160\cdot2^{-40}\le\frac1{16}.
\]
The chi-square estimate of
\cref{obj:lem:positive-family-certificate} uses the bounded-differences
exponential-moment conclusion of
\citep[Theorem~6.2, displayed log-mgf conclusion in its proof, p.~167]{BoucheronLugosiMassart2013}.
It supplies
\[
\chi^2(Q_n,Q^{(0)})\le e^{\eta_{\mathrm{samp}}}-1,
\]
together with absolute continuity and square-integrability of the centered
likelihood ratio.

Convexity of the exponential between \(0\) and \(1\), and \(e<3\), imply
\[
e^{\eta_{\mathrm{samp}}}-1
\le\eta_{\mathrm{samp}}(e-1)
\le2\eta_{\mathrm{samp}}.
\]
Therefore,
\[
\chi^2(Q_n,Q^{(0)})\le2\eta_{\mathrm{samp}}\le\frac14,
\qquad
\sqrt{\chi^2(Q_n,Q^{(0)})}\le\frac12.
\]
For every allowed \(M\), symmetry and the two comparison inequalities in
\cref{obj:lem:private-kernel-comparison} now give
\[
d_{\mathrm{TV}}(MQ^{(0)},MQ^{(1)})
=d_{\mathrm{TV}}(MQ_n,MQ^{(0)})
\le d_{\mathrm{TV}}(Q_n,Q^{(0)})
\le\frac12\sqrt{\chi^2(Q_n,Q^{(0)})}
\le\frac14.
\]
% lean: CausalSmith.Stat.PrivateCateRoughdesign.sampling_sign_output_TV_le

\item Next consider the sparse subcase
\[
a_n<b_n,\qquad c_n<b_n.
\]
Here \(n\epsilon>1\), so \(n^2\epsilon>1\) and \(0<b_n\le1\). Thus
\[
r(n,\epsilon)=b_n,\qquad h_{\mathrm L}=\frac{b_n}{4}.
\]
To establish the occupancy condition, apply the increasing logarithm to
\(c_n\le b_n\). This gives
\[
-\frac12(\log n+\log\epsilon)
\le-\frac17(2\log n+\log\epsilon),
\qquad
3\log n+5\log\epsilon\ge0.
\]
Consequently,
\[
\log(nb_n^5)
=\log n-\frac57(2\log n+\log\epsilon)
=-\frac17(3\log n+5\log\epsilon)
\le0.
\]
It follows that
\[
nb_n^5\le1,\qquad
n\delta_{\mathrm L}
=\frac{nb_n^5}{4^5}
\le\frac1{1024}\le\frac1{128}.
\]

The component estimate in
\cref{obj:lem:measurable-component-contraction} uses Cayley's labeled-tree
enumeration from
\citep[Section~5.6, Theorem~5.40]{KellerTrotterAppliedCombinatorics}.
With the occupancy condition just established, that estimate supplies
\[
d_{\mathrm{TV}}(MQ_n,MQ^{(0)})
\le2048\epsilon t n^2h_{\mathrm L}\delta_{\mathrm L}.
\]
The identity
\[
b_n^7=(n^2\epsilon)^{-1}
\]
calibrates its right-hand side:
\[
\begin{aligned}
2048\epsilon t n^2h_{\mathrm L}\delta_{\mathrm L}
&=2048\epsilon\kappa n^2h_{\mathrm L}^{7}\\
&=\frac{2048\kappa}{4^7}(n^2\epsilon)b_n^7\\
&=\frac{2048\kappa}{4^7}
 =2^{-15}\le\frac14.
\end{aligned}
\]
Symmetry therefore yields the required bound for
\(d_{\mathrm{TV}}(MQ^{(0)},MQ^{(1)})\).
% lean: CausalSmith.Stat.PrivateCateRoughdesign.measurable_component_contraction

\item Suppose now that the sign-family condition fails. Choose the singleton
families \((P_0)\) and \((P_1)\), with \(P_1\) defined at radius
\(h_{\mathrm L}\) by \cref{obj:def:direct-family}. The null membership and
target follow from \cref{obj:lem:positive-family-certificate}. We verify
the corresponding properties of \(P_1\).

For the localized profile \(q\) in \cref{obj:def:direct-family},
\[
0\le q(x)\le1,\qquad q(x_0)=1,\qquad x_0=\frac12.
\]
The calibration gives
\[
0<t=\kappa h_{\mathrm L}\le2^{-14}<\frac12,
\qquad
\frac12\le\mu_{1,P_1}(x)=\frac12+tq(x)\le1.
\]
To check the contrast modulus, define, for \(x\in[0,1]\),
\[
g_{\mathrm{loc}}(x)
:=\cos\!\left(
\frac{\pi}{2h_{\mathrm L}}
\min\{h_{\mathrm L},\max\{-h_{\mathrm L},x-x_0\}\}
\right).
\]
At and beyond either support boundary the cosine is zero, so
\[
q(x)=g_{\mathrm{loc}}(x)^4,\qquad |g_{\mathrm{loc}}(x)|\le1.
\]
Clipping to an interval and the cosine function are both \(1\)-Lipschitz;
therefore
\[
|g_{\mathrm{loc}}(x)-g_{\mathrm{loc}}(x')|
\le\frac{\pi}{2h_{\mathrm L}}|x-x'|.
\]
Factoring a difference of squares twice gives
\[
\begin{aligned}
|g_{\mathrm{loc}}(x)^2-g_{\mathrm{loc}}(x')^2|
&\le2|g_{\mathrm{loc}}(x)-g_{\mathrm{loc}}(x')|,\\
|q(x)-q(x')|
&\le2|g_{\mathrm{loc}}(x)^2-g_{\mathrm{loc}}(x')^2|
 \le\frac{2\pi}{h_{\mathrm L}}|x-x'|.
\end{aligned}
\]
Hence
\[
|\tau_{P_1}(x)-\tau_{P_1}(x')|
=t|q(x)-q(x')|
\le2\pi\kappa|x-x'|
\le L|x-x'|,
\]
because \(\pi<4\), \(\kappa=2^{-12}\), and \(L=3\) by
\cref{obj:synth_4}.

The covariate density is one, and the propensity and control regression
are constant \(1/2\). In particular,
\[
|e_{P_1}(x)-e_{P_1}(x')|
=|\mu_{0,P_1}(x)-\mu_{0,P_1}(x')|
=0\le L|x-x'|^{1/10}.
\]
These facts give the density, overlap, and nuisance smoothness conditions.
The conditional independent generation in
\cref{obj:def:direct-family} gives exchangeability, and setting
\(Y=Y(A)\) gives consistency. Thus, using
\cref{obj:def:complete-model},
\[
P_1\in\mathcal P,\qquad
\theta(P_1)=tq(x_0)=t=\Delta.
\]
Uniform averaging over either singleton leaves its product law unchanged:
\[
Q^{(0)}=(P_0^{\mathrm{obs}})^{\otimes n},\qquad
Q^{(1)}=(P_1^{\mathrm{obs}})^{\otimes n}.
\]
% lean: CausalSmith.Stat.PrivateCateRoughdesign.directAlternative_completeModel

\item Construct the direct dataset coupling as follows. For each
\(x\in[0,1]\), define the crossing mass
\[
\omega_{\mathrm{cross}}(x):=\frac{tq(x)}2,
\qquad
0\le\omega_{\mathrm{cross}}(x)\le\frac14.
\]
The conditional coupling has the following five atoms; every other pair
has mass zero:
\[
\begin{array}{c|c|c}
\text{First observed record}&\text{Second observed record}&\text{Mass}\\ \hline
(x,0,0)&(x,0,0)&1/4\\
(x,0,1)&(x,0,1)&1/4\\
(x,1,0)&(x,1,0)&1/4-\omega_{\mathrm{cross}}(x)\\
(x,1,1)&(x,1,1)&1/4\\
(x,1,0)&(x,1,1)&\omega_{\mathrm{cross}}(x)
\end{array}
\]
The masses are nonnegative and sum to one:
\[
\frac14+\frac14+
\left(\frac14-\omega_{\mathrm{cross}}(x)\right)
+\frac14+\omega_{\mathrm{cross}}(x)=1.
\]
The first marginal has mass \(1/4\) at each binary treatment--outcome
configuration. In the second marginal the control-cell masses are
\(1/4,1/4\), while the treated-cell masses satisfy
\[
\frac14-\omega_{\mathrm{cross}}(x)
=\frac12\left(\frac12-tq(x)\right),\qquad
\frac14+\omega_{\mathrm{cross}}(x)
=\frac12\left(\frac12+tq(x)\right).
\]
Thus these marginals are exactly the conditional observed laws of
\(P_0\) and \(P_1\).

Integrate this measurable conditional coupling over uniform \(X\), and
take \(n\) independent copies. Group their first and second coordinates
into \(D,D'\in\mathcal O^n\), and define
\[
\gamma_{\mathrm{dir}}:=\mathcal L(D,D').
\]
The product construction gives
\[
\mathcal L(D)=Q^{(0)},\qquad \mathcal L(D')=Q^{(1)}.
\]
At a fixed covariate, precisely the crossing atom contributes to record
disagreement. Since \(q\le1\) and vanishes when
\(|x-x_0|\ge h_{\mathrm L}\),
\[
\int_0^1q(x)\,dx
\le\int_0^1\mathbf1\{|x-x_0|<h_{\mathrm L}\}\,dx
\le2h_{\mathrm L}.
\]
Adding the \(n\) record disagreement costs therefore yields
\[
\begin{aligned}
\int d_{\mathrm H}(D,D')\,\gamma_{\mathrm{dir}}(dD,dD')
&=\sum_{i=1}^n\Pr\{D_i\ne D'_i\}\\
&=n\int_0^1\omega_{\mathrm{cross}}(x)\,dx\\
&=\frac{nt}{2}\int_0^1q(x)\,dx
\le nth_{\mathrm L}.
\end{aligned}
\]
By \cref{obj:lem:private-kernel-comparison} and the definition of the
infimum coupling cost,
\[
\begin{aligned}
d_{\mathrm{TV}}(MQ^{(0)},MQ^{(1)})
&\le2\epsilon W_{\mathrm H}(Q^{(0)},Q^{(1)})\\
&\le2\epsilon
 \int d_{\mathrm H}(D,D')\,\gamma_{\mathrm{dir}}(dD,dD')\\
&\le2\epsilon nth_{\mathrm L}.
\end{aligned}
\]
% lean: CausalSmith.Stat.PrivateCateRoughdesign.directDatasetCoupling_cost_le

\item It remains to calibrate this direct privacy bound. We first show
that failure of the sign-family condition implies
\[
r(n,\epsilon)=\min\{1,c_n\}.
\]
If \(n\epsilon>1\), suppose that \(c_n<a_n\). When \(b_n\le a_n\), the
sampling sign condition holds. When \(a_n<b_n\), we also have
\(c_n<b_n\), so the sparse sign condition holds. Both contradict the
present branch. Thus \(a_n\le c_n\). If \(c_n<b_n\), then
\(a_n\le c_n<b_n\), again giving the sparse sign condition. Consequently,
\[
a_n\le c_n,\qquad b_n\le c_n,\qquad
r(n,\epsilon)=\min\{1,c_n\}.
\]
If \(n\epsilon\le1\), monotonicity of the negative power gives
\[
c_n=(n\epsilon)^{-1/2}\ge1,\qquad
r(n,\epsilon)=1=\min\{1,c_n\}.
\]
This proves the identity throughout the direct branch, including
\(n\epsilon=1\) and ties among the candidate scales.

Since all quantities are nonnegative,
\[
r(n,\epsilon)\le c_n,\qquad
c_n^2=(n\epsilon)^{-1},\qquad
r(n,\epsilon)^2\le(n\epsilon)^{-1},
\qquad
(n\epsilon)r(n,\epsilon)^2\le1.
\]
Using \(t=\kappa h_{\mathrm L}\) and
\(h_{\mathrm L}=r(n,\epsilon)/4\), we obtain
\[
2\epsilon nth_{\mathrm L}
=2\epsilon n\kappa h_{\mathrm L}^2
=\frac{\kappa}{8}(n\epsilon)r(n,\epsilon)^2
\le\frac{\kappa}{8}.
\]
For every standard Borel output space and every private kernel in the
statement, the direct comparison therefore satisfies
\[
d_{\mathrm{TV}}(MQ^{(0)},MQ^{(1)})
\le\frac{\kappa}{8}
=2^{-15}\le\frac14.
\]
Together with the two sign subcases, this completes the construction and
the claimed output comparison.
% lean: CausalSmith.Stat.PrivateCateRoughdesign.testing_direct_privacy_scale_le
\end{enumerate}
\end{proof}

The certificate places a common target separation of benchmark order alongside a uniform bound on released-output distinguishability. Uniform averaging preserves a fixed target within each family, even though the sign family varies the nuisance functions across its members. Its quantification over standard Borel output spaces permits the same priors to enter the scalar and interval decision problems. The experiment changes across budget regimes, while the separation and output-comparison guarantees retain a common form.

\subsection{Scalar absolute loss}

The scalar reduction underlying \cref{obj:thm:matched-risk-frontier} uses the two target values \(0\) and \(\Delta\) to associate an estimator with a test whose cutoff is \(\Delta/2\). Absolute error controls the probability of assigning an estimate to the opposite side of this cutoff. The separation and output comparison in \cref{obj:lem:frontier-testing-priors} therefore supply the testing obstruction for the maximal expected absolute loss of \cref{obj:def:risk-criterion}.

Finite-prior averaging connects this obstruction to the class-wide criterion: the maximal risk over \(\mathcal P\) bounds the average risk under each prior. The resulting lower bound, recorded in \cref{obj:thm:matched-risk-frontier}, is
\[
R_{n,\epsilon}\ge2^{-20}r(n,\epsilon).
\]
Together with the publicly tuned upper bound in that result, this characterizes the optimal order for every admissible sample size and privacy budget.

\subsection{Honest connected-interval length}

For the interval criterion, the relevant reduction concerns target containment. Uniform \(90\%\) coverage in \cref{obj:def:interval-criterion} applies to every member of each finite prior and hence to its averaged coverage probability. The common output comparison in \cref{obj:lem:frontier-testing-priors} is the ingredient used in \cref{obj:thm:sharp-interval-frontier} to obtain simultaneous containment under a common output law.

Connectedness supplies the geometric link to expected length: an interval containing both \(0\) and \(\Delta\) has Lebesgue length at least \(\Delta\). Thus the interval reduction uses the same target separation as scalar estimation, with honesty governing containment and connectedness governing length. The resulting lower bound in \cref{obj:thm:sharp-interval-frontier} is
\[
H_{n,\epsilon}\ge2^{-20}r(n,\epsilon).
\]
The publicly tuned interval in that result supplies the matching upper order and establishes the stated comparison between optimal honest expected length and minimax absolute risk.

\subsection{Budget balances and consistency}

The transition budgets in \cref{obj:prop:regime-and-sanity} are the pairwise balances of the benchmark scales:
\[
\begin{aligned}
n^{-1/4}=(n^2\epsilon)^{-1/7}
&\quad\Longleftrightarrow\quad \epsilon=n^{-1/4},\\
(n^2\epsilon)^{-1/7}=(n\epsilon)^{-1/2}
&\quad\Longleftrightarrow\quad \epsilon=n^{-3/5},\\
(n\epsilon)^{-1/2}=1
&\quad\Longleftrightarrow\quad \epsilon=n^{-1}.
\end{aligned}
\]
These balances delimit the sampling, sparse privacy, direct privacy, and saturated branches. The benchmark expressions agree at their shared endpoints. In the testing certificate, sampling dominance uses the sign family, strict sparse privacy dominance uses its component comparison, and the direct family covers the remaining branches, including the tie between the two privacy scales.

Along an admissible budget sequence \((\epsilon_n)\), the direct privacy scale identifies the consistency condition \(n\epsilon_n\to+\infty\). Under this condition, both privacy scales vanish, since \(n^2\epsilon_n=n(n\epsilon_n)\), and the sampling scale also vanishes. The lower bound in \cref{obj:thm:matched-risk-frontier} supplies the converse through the truncated direct privacy scale. Accordingly, \cref{obj:prop:regime-and-sanity} gives
\[
R_{n,\epsilon_n}\longrightarrow0
\quad\Longleftrightarrow\quad
n\epsilon_n\longrightarrow+\infty.
\]
The comparison of the two decision criteria in \cref{obj:thm:sharp-interval-frontier} gives the same condition for vanishing optimal honest expected length. The finite testing priors therefore support both finite-sample lower bounds and the common consistency criterion.


\section{Verification note}

This appendix records the scope of the formal links attached to statements in the interactive bundle. It does not assert that every presentation-level definition or cited source theorem is formalized. Each mapped environment links to the declarations used in its audit, and theorem-local footnotes distinguish imported published conclusions from the checked downstream derivations.

The three registered published premises are \texttt{EfronSteinReplacement} (\citet[Theorem 3.1, replacement-copy form, p.~54]{BoucheronLugosiMassart2013}), \texttt{BoundedDifferencesMGF} (\citet[Theorem 6.2, displayed log-mgf conclusion in its proof, p.~167]{BoucheronLugosiMassart2013}), and \texttt{CayleyLabeledTreeCount} (\citet[Section 5.6, Theorem 5.40]{KellerTrotterAppliedCombinatorics}). Their source proofs are premises, as are the statistical assumptions defining \cref{obj:def:complete-model}. The procedures in \cref{obj:def:private-ratio,obj:def:public-tuning,obj:def:interval-handle} are exact-real randomized constructions with the measurable release convention in \cref{obj:def:private-kernels}. Their mathematical privacy, risk, and coverage conclusions do not constitute a finite-precision software guarantee; arithmetic, noise generation, and numerical endpoint handling require a separate implementation analysis.


\section{Proofs of the main results}\label{sec:deferred-proofs}

\begin{proof}[Proof of \cref{obj:thm:uniform-private-upper}]
\begin{enumerate}
\item
We first establish privacy of the local releases. The public scales satisfy
\[
\delta=\frac{2h}{k}>0,
\qquad
\mathcal A=nh\min\{1,n\delta\}>0,
\qquad
d_0=\frac{3\mathcal A}{64}>0.
\]
The finite cell statistics are measurable, and the positive denominator floor makes the ratio and interval maps everywhere defined and measurable. For
\[
u=(u_N,u_D)\in\mathbb R^2,
\]
the conditional density of the joint release is
\[
p_D(u)
=
\left(\frac{\epsilon}{24}\right)^2
\exp\left\{
-\frac{\epsilon}{12}
\bigl(|u_N-S_N(D)|+|u_D-S_D(D)|\bigr)
\right\}.
\]
It integrates to one for every dataset. For replacement-adjacent datasets, \cref{obj:lem:all-count-stability} supplies
\[
|S_N(D)-S_N(D')|+|S_D(D)-S_D(D')|\le12.
\]
The triangle inequality therefore gives
\[
p_D(u)\le e^\epsilon p_{D'}(u)
\qquad\text{for every }u\in\mathbb R^2.
\]
Integration proves the required privacy inequality for every Borel output event. Measurable postprocessing preserves both the Markov property and this inequality, proving the assertions for the scalar and interval outputs in \cref{obj:def:private-ratio}.
% lean: CausalSmith.Stat.PrivateCateRoughdesign.privateRatioRelease_private; CausalSmith.Stat.PrivateCateRoughdesign.Thk_private; CausalSmith.Stat.PrivateCateRoughdesign.Ihk_private

\item
The public tuning also preserves these properties. When \(r(n,\epsilon)\ge1/8\), the outputs in \cref{obj:def:public-tuning} are constant probability kernels and satisfy the privacy inequality.

On the branch \(r(n,\epsilon)<1/8\), define
\[
q_{\mathrm{dyad}}
=
\left\lceil\log_2\frac1{r(n,\epsilon)}\right\rceil,
\qquad
h_* =2^{-q_{\mathrm{dyad}}},
\qquad
k_* =2\cdot2^{4q_{\mathrm{dyad}}},
\qquad
\delta_*=\frac{2h_*}{k_*}.
\]
Here \(r(n,\epsilon)>0\), and \(q_{\mathrm{dyad}}\) is a nonnegative integer. The ceiling inequalities imply
\[
\frac1{r(n,\epsilon)}
\le2^{q_{\mathrm{dyad}}}
\le\frac2{r(n,\epsilon)},
\]
hence
\[
0<\frac{r(n,\epsilon)}2
\le h_*
\le r(n,\epsilon)<\frac18.
\]
Moreover,
\[
k_*=2h_*^{-4}\in\mathbb N,
\qquad
k_*\ge2,
\qquad
\delta_*=h_*^5.
\]
Thus the active branch uses admissible local parameters. Applying the local privacy conclusions proves the Markov and privacy assertions for both publicly tuned outputs.
% lean: CausalSmith.Stat.PrivateCateRoughdesign.tunedH_bounds; CausalSmith.Stat.PrivateCateRoughdesign.public_tuning_parameters; CausalSmith.Stat.PrivateCateRoughdesign.publicTunedRelease_private; CausalSmith.Stat.PrivateCateRoughdesign.publicTunedInterval_private

\item
It remains to establish the statistical guarantees. We reduce them to the following uniform squared-error claim: for every
\[
k'\in\mathbb N,\quad k'\ge2,
\qquad
h'\in(0,1/4],
\qquad
\delta'=\frac{2h'}{k'},
\qquad
P\in\mathcal P,
\]
one has
\[
E_{(P^{\mathrm{obs}})^{\otimes n},\widehat T_{h',k'}}
\bigl[(\widehat T_{h',k'}-\theta(P))^2\bigr]
\le V(h',\delta').
\]
We first derive the remaining conclusions from this claim, and then prove it.
% lean: CausalSmith.Stat.PrivateCateRoughdesign.uniform_private_upper

\item
For the original local parameters, define the interval radius
\[
\rho=\sqrt{10V(h,\delta)}>0.
\]
By \cref{obj:lem:point-version}, the target belongs to \([-1,1]\). Consequently, the endpoint formulas in \cref{obj:def:private-ratio} give
\[
\{\theta(P)\notin\widehat I_{h,k}\}
\subseteq
\{|\widehat T_{h,k}-\theta(P)|>\rho\}
\subseteq
\{(\widehat T_{h,k}-\theta(P))^2\ge10V(h,\delta)\}.
\]
Markov's inequality and the squared-error claim yield
\[
\Pr\{\theta(P)\notin\widehat I_{h,k}\}
\le
\frac{
E[(\widehat T_{h,k}-\theta(P))^2]
}{10V(h,\delta)}
\le\frac1{10}.
\]
All probabilities and expectations here integrate both the sample and release randomization. This proves the local coverage assertion, including targets at the endpoints of \([-1,1]\).

The interval contains its clipped center. Its endpoints lie in \([-1,1]\) and within distance \(\rho\) of that center, so, on every release,
\[
0\le\operatorname{Leb}(\widehat I_{h,k})
\le2,
\qquad
\operatorname{Leb}(\widehat I_{h,k})\le2\rho.
\]
Integrating gives
\[
E\operatorname{Leb}(\widehat I_{h,k})
\le\min\{2,2\sqrt{10V(h,\delta)}\}.
\]
% lean: CausalSmith.Stat.PrivateCateRoughdesign.Ihk_coverage_of_squared_error_le; CausalSmith.Stat.PrivateCateRoughdesign.Ihk_expectedLength_le

\item
We next prove the uniform absolute-error bound for the public estimator. If \(r(n,\epsilon)\ge1/8\), \cref{obj:lem:point-version} gives
\[
E|\widehat T^*-\theta(P)|
=|\theta(P)|
\le1
\le2^{18}r(n,\epsilon).
\]

For the active branch, the cap in the benchmark definition is inactive. Raising its three constituent lower bounds to their respective positive integer powers gives
\[
1\le nr(n,\epsilon)^4,
\qquad
1\le n^2\epsilon r(n,\epsilon)^7,
\qquad
1\le n\epsilon r(n,\epsilon)^2.
\]
In particular, since \(0<r(n,\epsilon)<1\),
\[
nr(n,\epsilon)^3\ge nr(n,\epsilon)^4\ge1,
\qquad
n^2r(n,\epsilon)^8
=\bigl(nr(n,\epsilon)^4\bigr)^2\ge1.
\]
Define the tuned occupancy scale by
\[
\mathcal A_*
=nh_*\min\{1,n\delta_*\}
=\min\{nh_*,n^2h_*^6\}.
\]
The dyadic bounds imply
\[
h_*\ge\frac{r(n,\epsilon)}{64},
\qquad
h_*^6\ge\frac{r(n,\epsilon)^6}{64},
\]
and therefore
\[
\mathcal A_*
\ge
\frac1{64}
\min\{nr(n,\epsilon),n^2r(n,\epsilon)^6\}.
\]
Multiplying by the appropriate positive scales yields
\[
\begin{aligned}
r(n,\epsilon)^2\mathcal A_*
&\ge
\frac1{64}
\min\{nr(n,\epsilon)^3,n^2r(n,\epsilon)^8\}
\ge\frac1{64},\\
r(n,\epsilon)\epsilon\mathcal A_*
&\ge
\frac1{64}
\min\{n\epsilon r(n,\epsilon)^2,
      n^2\epsilon r(n,\epsilon)^7\}
\ge\frac1{64}.
\end{aligned}
\]
Thus
\[
\mathcal A_*^{-1}\le64r(n,\epsilon)^2,
\qquad
(\epsilon\mathcal A_*)^{-1}\le64r(n,\epsilon),
\qquad
(\epsilon\mathcal A_*)^{-2}\le4096r(n,\epsilon)^2.
\]
Since \(\delta_*=h_*^5\) and \(h_*>0\),
\[
\delta_*^{2/5}=h_*^2\le r(n,\epsilon)^2.
\]
Substitution into the public envelope gives
\[
V(h_*,\delta_*)
\le
2^{20}(2+64+4096)r(n,\epsilon)^2
\le2^{33}r(n,\epsilon)^2
\le\bigl(2^{18}r(n,\epsilon)\bigr)^2.
\]

Define the positive absolute-error scale
\[
\beta_{\mathrm{abs}}=2^{18}r(n,\epsilon)>0.
\]
For every scalar output \(u_{\mathrm{sc}}\in\mathbb R\), expanding the nonnegative square
\((|u_{\mathrm{sc}}-\theta(P)|-\beta_{\mathrm{abs}})^2\) gives
\[
|u_{\mathrm{sc}}-\theta(P)|
\le
\frac{(u_{\mathrm{sc}}-\theta(P))^2}{2\beta_{\mathrm{abs}}}
+\frac{\beta_{\mathrm{abs}}}{2}.
\]
The uniform squared-error claim applies at \((h_*,k_*)\). Integrating the last inequality therefore proves
\[
E|\widehat T^*-\theta(P)|
\le
\frac{V(h_*,\delta_*)}{2\beta_{\mathrm{abs}}}
+\frac{\beta_{\mathrm{abs}}}{2}
\le\beta_{\mathrm{abs}}.
\]
Taking the supremum over \(P\in\mathcal P\) proves the stated risk bound.
% lean: CausalSmith.Stat.PrivateCateRoughdesign.active_rate_constraints; CausalSmith.Stat.PrivateCateRoughdesign.tuned_info_lower; CausalSmith.Stat.PrivateCateRoughdesign.tuned_scaled_info_lower; CausalSmith.Stat.PrivateCateRoughdesign.tuned_inverse_info_upper; CausalSmith.Stat.PrivateCateRoughdesign.tuned_cell_roughness; CausalSmith.Stat.PrivateCateRoughdesign.tuned_Vbound_le; CausalSmith.Stat.PrivateCateRoughdesign.absolute_moment_le_of_squared_moment; CausalSmith.Stat.PrivateCateRoughdesign.publicTunedRelease_worstRisk_of_squared_error_le

\item
For the public interval, the fallback branch has deterministic length two, and
\[
E\operatorname{Leb}(\widehat I^*)=2
\le2^{22}r(n,\epsilon)
\qquad\text{when }r(n,\epsilon)\ge1/8.
\]
On the active branch, the envelope estimate just proved gives
\[
10V(h_*,\delta_*)
\le10\cdot2^{33}r(n,\epsilon)^2
\le\bigl(2^{21}r(n,\epsilon)\bigr)^2.
\]
The right-hand square root is nonnegative, so
\[
\sqrt{10V(h_*,\delta_*)}\le2^{21}r(n,\epsilon).
\]
The local expected-length bound consequently yields
\[
E\operatorname{Leb}(\widehat I^*)
\le2\sqrt{10V(h_*,\delta_*)}
\le2^{22}r(n,\epsilon).
\]
Taking the supremum over the model proves the public expected-length assertion.
% lean: CausalSmith.Stat.PrivateCateRoughdesign.tuned_interval_radius_le; CausalSmith.Stat.PrivateCateRoughdesign.publicTunedInterval_expectedLength_le; CausalSmith.Stat.PrivateCateRoughdesign.publicTunedInterval_worstLength_le

\item
Public coverage follows through the same two tuning branches. On the fallback branch, \cref{obj:lem:point-version} gives
\[
\Pr\{\theta(P)\in[-1,1]\}=1.
\]
On the active branch, the admissibility of \(h_*,k_*\) allows application of the uniform squared-error claim:
\[
E[(\widehat T_{h_*,k_*}-\theta(P))^2]
\le V(h_*,\delta_*).
\]
The local coverage argument then gives
\[
\Pr\{\theta(P)\in\widehat I^*\}
=
\Pr\{\theta(P)\in\widehat I_{h_*,k_*}\}
\ge\frac9{10}.
\]
% lean: CausalSmith.Stat.PrivateCateRoughdesign.publicTunedInterval_coverage_of_squared_error_le

\item
We now prove the uniform squared-error claim. Fix arbitrary admissible local parameters \(h,k\) and \(P\in\mathcal P\). For measurable integrands, write
\[
\mathbb E_P[F(D,\widetilde S)]
:=
\int_{\mathcal O^n}\int_{\mathbb R^2}
F(D,u)\,\widetilde S(D,du)\,
(P^{\mathrm{obs}})^{\otimes n}(dD).
\]
Thus this expectation includes the independent Laplace randomization. Define the population means, bias envelope, and comparison ratio by
\[
\overline N=\mathbb E_P[S_N(D)],
\qquad
\overline D=\mathbb E_P[S_D(D)],
\qquad
b=3h+24\delta^{1/5},
\qquad
\theta^\circ=\frac{\overline N}{\overline D}.
\]
By \cref{obj:lem:occupancy-cancellation},
\[
\overline N=\sum_{j=1}^k w_j\nu_j,
\qquad
\overline D=\sum_{j=1}^k w_jd_j,
\qquad
\overline D\ge d_0>0,
\qquad
|\theta^\circ-\theta(P)|\le b.
\]

For binary values
\(a_{\mathrm{obs}},a'_{\mathrm{obs}},
y_{\mathrm{obs}},y'_{\mathrm{obs}}\in\{0,1\}\),
\[
\left|
(a_{\mathrm{obs}}-a'_{\mathrm{obs}})
(y_{\mathrm{obs}}-y'_{\mathrm{obs}})
\right|
\le
(a_{\mathrm{obs}}-a'_{\mathrm{obs}})^2.
\]
Taking expectations under two independent observations conditional on cell \(j\) proves
\[
|\nu_j|\le d_j.
\]
The cell conditioning is well defined by the positive cell masses under the density bounds. The occupancy weights in \cref{obj:lem:occupancy-cancellation} are nonnegative, hence
\[
|\overline N|
\le\sum_{j=1}^k w_j|\nu_j|
\le\sum_{j=1}^k w_jd_j
=\overline D,
\qquad
|\theta^\circ|\le1.
\]
Also, \cref{obj:lem:point-version} supplies
\[
\theta(P)\in[-1,1].
\]
% lean: CausalSmith.Stat.PrivateCateRoughdesign.conditional_cell_numerator_abs_le_denominator; CausalSmith.Stat.PrivateCateRoughdesign.population_numerator_abs_le_denominator; CausalSmith.Stat.PrivateCateRoughdesign.population_ratio_abs_le_one

\item
For every possible joint output \(u=(u_N,u_D)\in\mathbb R^2\), define
\[
B_D(u)=\max\{d_0,u_D\},
\qquad
T_{\mathrm{loc}}(u)
=
\operatorname{clip}_{[-1,1]}
\left(\frac{u_N}{B_D(u)}\right).
\]
These maps are defined on all of \(\mathbb R^2\), with \(B_D(u)\ge d_0>0\). Since \(\overline D\ge d_0\), the denominator floor fixes \(\overline D\) and satisfies
\[
|B_D(u)-\overline D|\le|u_D-\overline D|.
\]
The exact perturbation identity is
\[
\frac{u_N}{B_D(u)}-\theta^\circ
=
\frac{
(u_N-\overline N)+
\theta^\circ(\overline D-B_D(u))
}{B_D(u)}.
\]
Using \(|\theta^\circ|\le1\) gives
\[
\left|\frac{u_N}{B_D(u)}-\theta^\circ\right|
\le
\frac{|u_N-\overline N|+|u_D-\overline D|}{d_0}.
\]
Clipping cannot increase distance to \(\theta(P)\in[-1,1]\). Combining this fact with the population bias bound yields
\[
|T_{\mathrm{loc}}(u)-\theta(P)|
\le
b+\frac{|u_N-\overline N|+|u_D-\overline D|}{d_0}.
\]
Two applications of the quadratic sum inequality now give
\[
\begin{aligned}
(T_{\mathrm{loc}}(u)-\theta(P))^2
&\le
2b^2+
\frac{2}{d_0^2}
\bigl(|u_N-\overline N|+|u_D-\overline D|\bigr)^2\\
&\le
2b^2+\frac4{d_0^2}
\left\{
(u_N-\overline N)^2+(u_D-\overline D)^2
\right\}.
\end{aligned}
\]
% lean: CausalSmith.Stat.PrivateCateRoughdesign.floored_ratio_perturbation; CausalSmith.Stat.PrivateCateRoughdesign.ratioMap_error_le; CausalSmith.Stat.PrivateCateRoughdesign.ratioMap_squared_error_le; CausalSmith.Stat.PrivateCateRoughdesign.ratioMap_model_squared_error_le

\item
For the variance calculation, invoke the replacement-copy Efron--Stein inequality \citep[Theorem~3.1, replacement-copy form, p.~54]{BoucheronLugosiMassart2013}. Explicitly, for the independent observed records and their independent replacement copies, define
\[
D=(O_1,\ldots,O_n),
\qquad
D^{(i)}=(O_1,\ldots,O_{i-1},O_i',O_{i+1},\ldots,O_n),
\]
where \(O_1,\ldots,O_n,O_i'\) have law \(P^{\mathrm{obs}}\), and \(O_i'\) is independent of \(D\). The published inequality gives
\[
\operatorname{Var}(\mathsf S(D))
\le
\frac12\sum_{i=1}^n
E\bigl[(\mathsf S(D)-\mathsf S(D^{(i)}))^2\bigr],
\qquad
\mathsf S\in\{S_N,S_D\}.
\]
The statistics are measurable and square integrable. With this inequality, \cref{obj:lem:all-count-stability} supplies the bound
\[
\max\{\operatorname{Var}(S_N(D)),\operatorname{Var}(S_D(D))\}
\le18W(P),
\qquad
W(P)=\sum_{j=1}^k w_j.
\]

Each independent Laplace noise has
\[
E[\xi_N]=E[\xi_D]=0,
\qquad
E[\xi_N^2]=E[\xi_D^2]
=2\left(\frac{12}{\epsilon}\right)^2
=\frac{288}{\epsilon^2}.
\]
Expanding the centered squares, independence and the zero noise means eliminate the cross terms. Thus the exact released-moment identity is
\[
\begin{aligned}
&\mathbb E_P\left[
(\widetilde S_N-\overline N)^2+
(\widetilde S_D-\overline D)^2
\right]\\
&\qquad=
\operatorname{Var}(S_N(D))
+\operatorname{Var}(S_D(D))
+\frac{576}{\epsilon^2}
\le36W(P)+\frac{576}{\epsilon^2}.
\end{aligned}
\]
Integrating the pointwise ratio bound consequently gives
\[
\begin{aligned}
\mathbb E_P[(\widehat T_{h,k}-\theta(P))^2]
&\le
2b^2+\frac4{d_0^2}
\mathbb E_P\left[
(\widetilde S_N-\overline N)^2+
(\widetilde S_D-\overline D)^2
\right]\\
&\le
2b^2+\frac4{d_0^2}
\left(36W(P)+\frac{576}{\epsilon^2}\right)\\
&=
2b^2+\frac{144W(P)}{d_0^2}
+\frac{2304}{\epsilon^2d_0^2}.
\end{aligned}
\]
% lean: CausalSmith.Stat.PrivateCateRoughdesign.privateRatioRelease_centered_moments; CausalSmith.Stat.PrivateCateRoughdesign.Thk_model_squared_error_le

\item
Finally, the quadratic sum inequality and positivity of \(\delta\) imply
\[
(\delta^{1/5})^2=\delta^{2/5},
\qquad
2b^2
=2(3h+24\delta^{1/5})^2
\le36h^2+2304\delta^{2/5}.
\]
By \cref{obj:lem:occupancy-cancellation},
\[
W(P)\le\frac92\mathcal A.
\]
Substituting \(d_0=3\mathcal A/64\) gives
\[
\frac{144W(P)}{d_0^2}
\le
\frac{144(9\mathcal A/2)}{(3\mathcal A/64)^2}
=294912\mathcal A^{-1},
\]
and
\[
\frac{2304}{\epsilon^2d_0^2}
=
\frac{2304}{\epsilon^2(3\mathcal A/64)^2}
=1048576(\epsilon\mathcal A)^{-2}.
\]
All four terms in the public envelope are nonnegative. Since
\(36,2304,294912\le2^{20}\) and \(1048576=2^{20}\), we obtain
\[
\begin{aligned}
\mathbb E_P[(\widehat T_{h,k}-\theta(P))^2]
&\le
36h^2+2304\delta^{2/5}
+294912\mathcal A^{-1}
+1048576(\epsilon\mathcal A)^{-2}\\
&\le
2^{20}\left\{
h^2+\delta^{2/5}+\mathcal A^{-1}
+(\epsilon\mathcal A)^{-2}
\right\}
=V(h,\delta).
\end{aligned}
\]
The local parameters and model law were arbitrary, proving the uniform squared-error claim and completing all the preceding deductions.
% lean: CausalSmith.Stat.PrivateCateRoughdesign.bias_squared_le; CausalSmith.Stat.PrivateCateRoughdesign.variance_floor_envelope_le; CausalSmith.Stat.PrivateCateRoughdesign.model_variance_envelope_le; CausalSmith.Stat.PrivateCateRoughdesign.Thk_model_squared_error_le; CausalSmith.Stat.PrivateCateRoughdesign.uniform_private_upper
\end{enumerate}
\end{proof}

\begin{proof}[Proof of \cref{obj:thm:matched-risk-frontier}]
\begin{enumerate}
\item Fix \(n\ge2\) and \(0<\epsilon\le1\). The numerical benchmark satisfies
\[
0<r(n,\epsilon)\le1,
\]
because \(n^{-1/4}>0\) and the defining minimum is taken with \(1\).

Apply \cref{obj:lem:frontier-testing-priors}. Its probabilistic and combinatorial inputs are the bounded-differences exponential-moment bound from \citep[Theorem 6.2, displayed log-mgf conclusion in its proof, p.~167]{BoucheronLugosiMassart2013} and the labeled-tree count from \citep[Section 5.6, Theorem 5.40]{KellerTrotterAppliedCombinatorics}. Denote the resulting nonempty finite families and their uniformly averaged dataset laws by
\[
\begin{gathered}
N_{\mathrm{prior},0},N_{\mathrm{prior},1}\in\mathbb N_{\ge1},
\qquad
\mathcal F_\eta
=
(P_{\eta,\iota})_{\iota=1}^{N_{\mathrm{prior},\eta}},
\qquad \eta\in\{0,1\},\\
Q^{(\eta)}
=
\frac{1}{N_{\mathrm{prior},\eta}}
\sum_{\iota=1}^{N_{\mathrm{prior},\eta}}
(P_{\eta,\iota}^{\mathrm{obs}})^{\otimes n}.
\end{gathered}
\]
The supplied target separation \(\Delta\) obeys
\[
\Delta\ge2^{-14}r(n,\epsilon)>0,
\qquad
P_{\eta,\iota}\in\mathcal P,
\qquad
\theta(P_{\eta,\iota})=\eta\Delta.
\]
Moreover, for every \(M\in\mathfrak M_{n,\epsilon}(\mathbb R)\),
\[
d_{\mathrm{TV}}\!\left(MQ^{(0)},MQ^{(1)}\right)\le\frac14.
\]
% lean: CausalSmith.Stat.PrivateCateRoughdesign.frontier_testing_priors
% lean: CausalSmith.Stat.PrivateCateRoughdesign.rate_pos

\item Fix any \(M\in\mathfrak M_{n,\epsilon}(\mathbb R)\), and define its maximal absolute risk, its two output laws, and their error events by
\[
\begin{aligned}
\mathcal R(M)
&=
\sup_{P\in\mathcal P}
E_{(P^{\mathrm{obs}})^{\otimes n},M}
|M(D)-\theta(P)|,\\
\mathsf L_\eta&=MQ^{(\eta)},\\
\mathcal E_\eta
&=
\{u\in\mathbb R:\Delta/2\le |u-\eta\Delta|\},
\qquad \eta\in\{0,1\}.
\end{aligned}
\]
Each \(\mathsf L_\eta\) is a probability measure, since its dataset prior is a normalized nonempty finite mixture and \(M\) is a Markov kernel.

The target distance satisfies
\[
2(\Delta/2)=|0-\Delta|.
\]
Consequently, the triangle inequality gives
\(\mathcal E_0^c\subseteq\mathcal E_1\). The two-error testing inequality therefore yields
\[
\begin{aligned}
\mathsf L_0(\mathcal E_0)+\mathsf L_1(\mathcal E_1)
&\ge
\mathsf L_0(\mathcal E_0)+\mathsf L_1(\mathcal E_0^c)\\
&=
1+\mathsf L_0(\mathcal E_0)-\mathsf L_1(\mathcal E_0)\\
&\ge
1-d_{\mathrm{TV}}(\mathsf L_0,\mathsf L_1)
\ge\frac34.
\end{aligned}
\]
Thus
\[
\mathsf L_0(\mathcal E_0)\ge\frac38
\quad\text{or}\quad
\mathsf L_1(\mathcal E_1)\ge\frac38.
\]

In the first case take \(\eta=0\), and in the second take \(\eta=1\). The absolute loss is at least \(\Delta/2\) on the corresponding error event, so
\[
\int_{\mathbb R}|u-\eta\Delta|\,d\mathsf L_\eta(u)
\ge
\frac{\Delta}{2}\mathsf L_\eta(\mathcal E_\eta)
\ge
\frac{\Delta}{2}\frac38
=
\frac{3\Delta}{16}.
\]
% lean: CausalSmith.Stat.PrivateCateRoughdesign.absolute_loss_of_error_probability

Finite-prior averaging and the common target within each family give
\[
\begin{aligned}
\int_{\mathbb R}|u-\eta\Delta|\,d\mathsf L_\eta(u)
&=
\frac{1}{N_{\mathrm{prior},\eta}}
\sum_{\iota=1}^{N_{\mathrm{prior},\eta}}
E_{(P_{\eta,\iota}^{\mathrm{obs}})^{\otimes n},M}
|M(D)-\theta(P_{\eta,\iota})|\\
&\le
\frac{1}{N_{\mathrm{prior},\eta}}
\sum_{\iota=1}^{N_{\mathrm{prior},\eta}}\mathcal R(M)
=
\mathcal R(M).
\end{aligned}
\]
These identities and inequalities concern nonnegative extended integrals and remain valid when a risk is infinite. In either case,
\[
\frac{3\Delta}{16}\le\mathcal R(M).
\]
% lean: CausalSmith.Stat.PrivateCateRoughdesign.finite_prior_loss_le_worstRisk
% lean: CausalSmith.Stat.PrivateCateRoughdesign.finite_priors_scalar_lower

\item The numerical calibration is
\[
2^{-20}\le\frac{3}{16}\,2^{-14}.
\]
Multiplying by the nonnegative benchmark and then using the separation bound gives
\[
2^{-20}r(n,\epsilon)
\le
\frac{3}{16}\bigl(2^{-14}r(n,\epsilon)\bigr)
\le
\frac{3\Delta}{16}
\le
\mathcal R(M).
\]
Since this holds for every admissible scalar kernel, \cref{obj:def:risk-criterion} implies
\[
2^{-20}r(n,\epsilon)
\le
\inf_{M\in\mathfrak M_{n,\epsilon}(\mathbb R)}\mathcal R(M)
=
R_{n,\epsilon}.
\]
% lean: CausalSmith.Stat.PrivateCateRoughdesign.matched_risk_frontier

\item For the upper bound, use the replacement-copy Efron--Stein inequality from \citep[Theorem 3.1, p.~54]{BoucheronLugosiMassart2013} as the variance input to \cref{obj:thm:uniform-private-upper}. Apply that result with the admissible auxiliary parameters \(k=2\) and \(h=1/4\). It supplies, for the publicly tuned scalar estimator in \cref{obj:def:public-tuning},
\[
\widehat T^*\in\mathfrak M_{n,\epsilon}(\mathbb R),
\qquad
\mathcal R(\widehat T^*)\le2^{18}r(n,\epsilon).
\]
In particular, this estimator is an everywhere-defined Markov kernel satisfying pure replacement \(\epsilon\)-differential privacy. Its risk bound takes the supremum over the full class in \cref{obj:def:complete-model}, including every admissible measurable covariate density.

Using this kernel as a candidate in the minimax infimum gives
\[
R_{n,\epsilon}
=
\inf_{M\in\mathfrak M_{n,\epsilon}(\mathbb R)}\mathcal R(M)
\le
\mathcal R(\widehat T^*)
\le
2^{18}r(n,\epsilon).
\]
Together with the lower bound, this proves all the assertions.
% lean: CausalSmith.Stat.PrivateCateRoughdesign.uniform_private_upper
% lean: CausalSmith.Stat.PrivateCateRoughdesign.matched_risk_frontier
\end{enumerate}
\end{proof}

\begin{proof}[Proof of \cref{obj:thm:sharp-interval-frontier}]
\begin{enumerate}
\item Since \(n\ge2\), the sampling term is strictly positive. Hence
\[
0<n^{-1/4}
\le
\max\left\{
n^{-1/4},(n^2\epsilon)^{-1/7},(n\epsilon)^{-1/2}
\right\},
\qquad
r(n,\epsilon)>0.
\]
The definitions of the benchmark and constants therefore give
\[
\ell(n,\epsilon)=r(n,\epsilon)>0,
\qquad
c_I=2^{-20}>0,
\qquad
C_I=2^{22}>0.
\]
% lean: CausalSmith.Stat.PrivateCateRoughdesign.rate_pos

\item We first identify the inversion kernel with the publicly tuned interval kernel. If \(r(n,\epsilon)\ge1/8\), both constructions in
\cref{obj:def:interval-handle,obj:def:public-tuning} return \([-1,1]\).

Suppose instead that \(r(n,\epsilon)<1/8\). Define the integer rounding index and public parameters by
\[
\jmath_h
=
\left\lceil\log_2\!\left(\frac1{r(n,\epsilon)}\right)\right\rceil
\in\mathbb Z_{\ge0},
\qquad
h=2^{-\jmath_h},
\qquad
k=2\,2^{4\jmath_h}\in\mathbb N.
\]
The ceiling inequalities imply
\[
\log_2\!\left(\frac1{r(n,\epsilon)}\right)
\le\jmath_h
\le
\log_2\!\left(\frac1{r(n,\epsilon)}\right)+1,
\]
and thus
\[
\frac1{r(n,\epsilon)}
\le2^{\jmath_h}
\le\frac2{r(n,\epsilon)},
\qquad
0<\frac{r(n,\epsilon)}2\le h\le r(n,\epsilon)<\frac18.
\]
In particular,
\[
0<h\le\frac14,
\qquad
k=2h^{-4}\ge2,
\qquad
\delta=\frac{2h}{k}=h^5>0.
\]
This also proves that the integer conversion in the public cell count is exact.

For every possible value of the joint release in
\cref{obj:def:private-ratio}, define
\[
\mathcal A=nh\min\{1,n\delta\},
\qquad
d_0=\frac{3\mathcal A}{64},
\qquad
B_D=\max\{d_0,\widetilde S_D\},
\qquad
B_N=\operatorname{clip}_{[-B_D,B_D]}(\widetilde S_N),
\qquad
\rho=\sqrt{10V(h,\delta)}.
\]
Because \(n,h,\delta>0\),
\[
\mathcal A>0,
\qquad
B_D\ge d_0>0,
\qquad
\rho\ge0.
\]
Division by the positive \(B_D\) preserves minima and maxima, so
\[
\frac{B_N}{B_D}
=
\frac{\max\{-B_D,\min\{B_D,\widetilde S_N\}\}}{B_D}
=
\max\left\{-1,\min\left\{1,\frac{\widetilde S_N}{B_D}\right\}\right\}
=
\widehat T_{h,k}.
\]
Consequently, the accepted candidate set is
\[
\begin{aligned}
\mathcal V_{\mathrm{acc}}
&=
\left\{\vartheta\in[-1,1]:
|B_N-\vartheta B_D|\le\rho B_D\right\}\\
&=
\left\{\vartheta\in[-1,1]:
\left|\frac{B_N}{B_D}-\vartheta\right|\le\rho\right\}\\
&=
\left[
\max\{-1,\widehat T_{h,k}-\rho\},
\min\{1,\widehat T_{h,k}+\rho\}
\right].
\end{aligned}
\]
Indeed, the absolute-value inequality is equivalent to
\(\widehat T_{h,k}-\rho\le\vartheta\le\widehat T_{h,k}+\rho\).
Moreover,
\[
-1\le\widehat T_{h,k}\le1,
\qquad
\max\{-1,\widehat T_{h,k}-\rho\}
\le\widehat T_{h,k}
\le
\min\{1,\widehat T_{h,k}+\rho\}.
\]
Thus the accepted set is a nonempty closed interval, and its hull has exactly these endpoints. The two constructions apply equal measurable interval maps to the same joint release. Together with the constant branch, this proves
\[
I^*_{n,\epsilon}=\widehat I^*.
\]
% lean: CausalSmith.Stat.PrivateCateRoughdesign.optimalIntervalHandle_eq_publicTunedInterval

\item By \cref{obj:lem:frontier-testing-priors}, choose positive integers \(s_0,s_1\), families
\[
(P^{(0)}_\iota)_{\iota=1}^{s_0}\subset\mathcal P,
\qquad
(P^{(1)}_\iota)_{\iota=1}^{s_1}\subset\mathcal P,
\]
and a separation satisfying
\[
\theta(P^{(0)}_\iota)=0,
\qquad
\theta(P^{(1)}_\iota)=\Delta,
\qquad
\Delta\ge2^{-14}r(n,\epsilon)>0.
\]
Their dataset mixture laws are
\[
Q^{(0)}
=
\frac1{s_0}\sum_{\iota=1}^{s_0}
\bigl((P^{(0)}_\iota)^{\mathrm{obs}}\bigr)^{\otimes n},
\qquad
Q^{(1)}
=
\frac1{s_1}\sum_{\iota=1}^{s_1}
\bigl((P^{(1)}_\iota)^{\mathrm{obs}}\bigr)^{\otimes n}.
\]
The testing-prior construction uses the bounded-differences log-mgf conclusion of
\citep[Theorem 6.2, proof, printed p.~167]{BoucheronLugosiMassart2013}
and the labelled-tree count of
\citep[Section 5.6, Theorem 5.40]{KellerTrotterAppliedCombinatorics}.
Its conclusion gives, for every private Markov interval kernel \(I\),
\[
d_{\mathrm{TV}}(IQ^{(0)},IQ^{(1)})\le\frac14.
\]

Fix any kernel \(I\) feasible for \cref{obj:def:interval-criterion}. Define its two output probability laws and the measurable containment events by
\[
\mathsf J_0=IQ^{(0)},
\qquad
\mathsf J_1=IQ^{(1)},
\qquad
\mathsf E_0=\{\mathsf i\in\mathcal I:0\in\mathsf i\},
\qquad
\mathsf E_1=\{\mathsf i\in\mathcal I:\Delta\in\mathsf i\}.
\]
The finite families are nonempty, so honesty at every constituent law yields
\[
\begin{aligned}
\mathsf J_0(\mathsf E_0)
&=
\frac1{s_0}\sum_{\iota=1}^{s_0}
\Pr_{((P^{(0)}_\iota)^{\mathrm{obs}})^{\otimes n},I}
\{0\in I(D)\}
\ge\frac9{10},\\
\mathsf J_1(\mathsf E_1)
&=
\frac1{s_1}\sum_{\iota=1}^{s_1}
\Pr_{((P^{(1)}_\iota)^{\mathrm{obs}})^{\otimes n},I}
\{\Delta\in I(D)\}
\ge\frac9{10}.
\end{aligned}
\]
Total variation comparison on \(\mathsf E_1\) gives
\[
\mathsf J_1(\mathsf E_1)-\mathsf J_0(\mathsf E_1)
\le d_{\mathrm{TV}}(\mathsf J_0,\mathsf J_1)
\le\frac14.
\]
Using the union--intersection identity and
\(\mathsf J_0(\mathsf E_0\cup\mathsf E_1)\le1\), we obtain
\[
\begin{aligned}
\mathsf J_0(\mathsf E_0\cap\mathsf E_1)
&=
\mathsf J_0(\mathsf E_0)+\mathsf J_0(\mathsf E_1)
-\mathsf J_0(\mathsf E_0\cup\mathsf E_1)\\
&\ge\frac9{10}+\left(\frac9{10}-\frac14\right)-1
=\frac{11}{20}.
\end{aligned}
\]

For every interval \(\mathsf i\in\mathsf E_0\cap\mathsf E_1\), its endpoint inequalities imply
\[
(0,\Delta)\subseteq\mathsf i,
\qquad
\operatorname{Leb}(\mathsf i)
\ge\operatorname{Leb}((0,\Delta))=\Delta.
\]
This argument applies to open and closed endpoints and to unbounded intervals; an empty interval cannot belong to the intersection. Therefore, with all length integrals understood as nonnegative extended integrals,
\[
\begin{aligned}
\frac{11}{20}\Delta
&\le\Delta\,\mathsf J_0(\mathsf E_0\cap\mathsf E_1)\\
&=
\int_{\mathcal I}
\Delta\,\mathbf 1_{\mathsf E_0\cap\mathsf E_1}(\mathsf i)
\,\mathsf J_0(d\mathsf i)\\
&\le
\int_{\mathcal I}\operatorname{Leb}(\mathsf i)\,
\mathsf J_0(d\mathsf i)\\
&=
\frac1{s_0}\sum_{\iota=1}^{s_0}
E_{((P^{(0)}_\iota)^{\mathrm{obs}})^{\otimes n},I}
\operatorname{Leb}(I(D))\\
&\le
\sup_{P\in\mathcal P}
E_{(P^{\mathrm{obs}})^{\otimes n},I}
\operatorname{Leb}(I(D)).
\end{aligned}
\]
The last inequality follows by bounding each constituent expectation by the same supremum.

The numerical calibration is
\[
2^{-20}\le\frac{11}{20}\,2^{-14}.
\]
Multiplying by the nonnegative benchmark and using the separation bound gives
\[
2^{-20}r(n,\epsilon)
\le
\frac{11}{20}\,2^{-14}r(n,\epsilon)
\le
\frac{11}{20}\Delta
\le
\sup_{P\in\mathcal P}
E_{(P^{\mathrm{obs}})^{\otimes n},I}
\operatorname{Leb}(I(D)).
\]
Since this holds for every feasible \(I\), taking the defining infimum proves
\[
2^{-20}r(n,\epsilon)\le H_{n,\epsilon}.
\]
% lean: CausalSmith.Stat.PrivateCateRoughdesign.finite_priors_interval_lower

\item Apply \cref{obj:thm:uniform-private-upper} with auxiliary cell count two and localization radius \(1/4\), which satisfy its parameter conditions. Its variance argument uses the replacement-copy Efron--Stein inequality from
\citep[Theorem 3.1, printed p.~54]{BoucheronLugosiMassart2013}.
The publicly tuned conclusions establish that \(\widehat I^*\) is an everywhere-defined Markov kernel satisfying pure replacement \(\epsilon\)-privacy, and that
\[
\Pr_{(P^{\mathrm{obs}})^{\otimes n},\widehat I^*}
\{\theta(P)\in\widehat I^*(D)\}
\ge\frac9{10}
\qquad(P\in\mathcal P),
\]
\[
\sup_{P\in\mathcal P}
E_{(P^{\mathrm{obs}})^{\otimes n},\widehat I^*}
\operatorname{Leb}(\widehat I^*(D))
\le2^{22}r(n,\epsilon).
\]
Thus \(\widehat I^*\) is feasible for the infimum in
\cref{obj:def:interval-criterion}, and
\[
H_{n,\epsilon}
\le
\sup_{P\in\mathcal P}
E_{(P^{\mathrm{obs}})^{\otimes n},\widehat I^*}
\operatorname{Leb}(\widehat I^*(D)).
\]
The kernel identity proved above transfers privacy, coverage, and length to \(I^*_{n,\epsilon}\), yielding
\[
2^{-20}r(n,\epsilon)
\le H_{n,\epsilon}
\le
\sup_{P\in\mathcal P}
E_{(P^{\mathrm{obs}})^{\otimes n},I^*_{n,\epsilon}}
\operatorname{Leb}(I^*_{n,\epsilon}(D))
\le2^{22}r(n,\epsilon).
\]
% lean: CausalSmith.Stat.PrivateCateRoughdesign.uniform_private_upper

\item Finally, \cref{obj:thm:matched-risk-frontier} supplies
\[
2^{-20}r(n,\epsilon)
\le R_{n,\epsilon}
\le2^{18}r(n,\epsilon).
\]
Multiplication by nonnegative finite constants preserves order in the extended nonnegative reals. The risk upper bound and interval lower bound give
\[
2^{-38}R_{n,\epsilon}
\le
2^{-38}\,2^{18}r(n,\epsilon)
=
2^{-20}r(n,\epsilon)
\le H_{n,\epsilon}.
\]
The interval upper bound and risk lower bound give
\[
H_{n,\epsilon}
\le2^{22}r(n,\epsilon)
=
2^{42}\bigl(2^{-20}r(n,\epsilon)\bigr)
\le2^{42}R_{n,\epsilon}.
\]
These are the asserted comparisons of the two frontiers.
% lean: CausalSmith.Stat.PrivateCateRoughdesign.sharp_interval_frontier
\end{enumerate}
\end{proof}

\begin{proof}[Proof of \cref{obj:prop:regime-and-sanity}]
\begin{enumerate}
\item Fix \(n\ge2\) and \(0<\epsilon\le1\), and define the three benchmark terms by
\[
r_{\mathrm{samp}}=n^{-1/4},
\qquad
r_{\mathrm{sparse}}=(n^2\epsilon)^{-1/7},
\qquad
r_{\mathrm{dense}}=(n\epsilon)^{-1/2}.
\]
All three terms are positive, and
\[
r(n,\epsilon)
=\min\{1,\max\{r_{\mathrm{samp}},
                    \max\{r_{\mathrm{sparse}},r_{\mathrm{dense}}\}\}\}.
\]
Because logarithms preserve order on positive numbers,
\[
\begin{aligned}
r_{\mathrm{sparse}}\le r_{\mathrm{samp}}
&\iff
-\frac17(2\log n+\log\epsilon)\le-\frac14\log n\\
&\iff \log\epsilon\ge-\frac14\log n
\iff \epsilon\ge n^{-1/4},
\\
r_{\mathrm{dense}}\le r_{\mathrm{samp}}
&\iff
-\frac12(\log n+\log\epsilon)\le-\frac14\log n\\
&\iff \log\epsilon\ge-\frac12\log n
\iff \epsilon\ge n^{-1/2}.
\end{aligned}
\]
If \(n^{-1/4}\le\epsilon\), then
\[
n^{-1/2}\le n^{-1/4}\le\epsilon,
\qquad
r_{\mathrm{sparse}}\le r_{\mathrm{samp}},
\qquad
r_{\mathrm{dense}}\le r_{\mathrm{samp}}.
\]
Here the first comparison follows from \(n\ge1\) and
\(-1/2\le-1/4\). Also \(r_{\mathrm{samp}}\le1\), since its exponent is nonpositive. Consequently,
\[
r(n,\epsilon)
=\min\{1,r_{\mathrm{samp}}\}
=r_{\mathrm{samp}}
=n^{-1/4}.
\]
% lean: CausalSmith.Stat.PrivateCateRoughdesign.rate_sampling_regime

\item Suppose that \(n^{-3/5}\le\epsilon\le n^{-1/4}\).
The upper budget bound gives
\[
\log\epsilon\le-\frac14\log n,
\qquad
-\frac14\log n
\le-\frac17(2\log n+\log\epsilon),
\]
and hence
\[
r_{\mathrm{samp}}\le r_{\mathrm{sparse}}.
\]
For the other comparison, logarithms give
\[
\begin{aligned}
r_{\mathrm{dense}}\le r_{\mathrm{sparse}}
&\iff
-\frac12(\log n+\log\epsilon)
\le-\frac17(2\log n+\log\epsilon)\\
&\iff \log\epsilon\ge-\frac35\log n
\iff \epsilon\ge n^{-3/5}.
\end{aligned}
\]
Thus \(r_{\mathrm{dense}}\le r_{\mathrm{sparse}}\).
Moreover, monotonicity in the exponent for a base at least one yields
\[
n^{-1}\le n^{-3/5}\le\epsilon.
\]
Multiplying by \(n>0\), and then using \(n\ge1\), gives
\[
1\le n\epsilon\le n^2\epsilon.
\]
Since the exponent \(-1/7\) is nonpositive,
\[
r_{\mathrm{sparse}}=(n^2\epsilon)^{-1/7}\le1.
\]
The maximum and the cap therefore simplify to
\[
r(n,\epsilon)
=\min\{1,r_{\mathrm{sparse}}\}
=r_{\mathrm{sparse}}
=(n^2\epsilon)^{-1/7}.
\]
% lean: CausalSmith.Stat.PrivateCateRoughdesign.rate_sparse_regime

\item Suppose that \(n^{-1}\le\epsilon\le n^{-3/5}\).
Reversing the preceding logarithmic comparison gives
\[
\begin{aligned}
r_{\mathrm{sparse}}\le r_{\mathrm{dense}}
&\iff
-\frac17(2\log n+\log\epsilon)
\le-\frac12(\log n+\log\epsilon)\\
&\iff \log\epsilon\le-\frac35\log n
\iff \epsilon\le n^{-3/5}.
\end{aligned}
\]
Also,
\[
\epsilon\le n^{-3/5}\le n^{-1/2},
\]
where the second inequality follows from \(n\ge1\).
Taking logarithms and rearranging gives
\[
\log\epsilon\le-\frac12\log n,
\qquad
-\frac14\log n
\le-\frac12(\log n+\log\epsilon),
\]
so that
\[
r_{\mathrm{samp}}\le r_{\mathrm{dense}}.
\]
The lower budget bound gives
\[
n\epsilon\ge1,
\qquad
r_{\mathrm{dense}}=(n\epsilon)^{-1/2}\le1.
\]
It follows that
\[
r(n,\epsilon)
=\min\{1,r_{\mathrm{dense}}\}
=r_{\mathrm{dense}}
=(n\epsilon)^{-1/2}.
\]
% lean: CausalSmith.Stat.PrivateCateRoughdesign.rate_dense_regime

\item If \(\epsilon\le n^{-1}\), then positivity and multiplication by \(n\) give
\[
0<n\epsilon\le1.
\]
A nonpositive power reverses order on positive bases, so
\[
1\le(n\epsilon)^{-1/2}
=r_{\mathrm{dense}}
\le\max\{r_{\mathrm{samp}},
             \max\{r_{\mathrm{sparse}},r_{\mathrm{dense}}\}\}.
\]
Therefore the cap is active:
\[
r(n,\epsilon)=1.
\]
% lean: CausalSmith.Stat.PrivateCateRoughdesign.rate_capped_regime

\item Since \(n\ge1\),
\[
n^{-1/4}\le1.
\]
Applying the sampling-regime identity at \(\epsilon=1\) yields
\[
r(n,1)=n^{-1/4}.
\]
% lean: CausalSmith.Stat.PrivateCateRoughdesign.rate_unit_budget

\item Fix a budget sequence satisfying \(0<\epsilon_n\le1\) for every \(n\ge2\).
The replacement-copy variance inequality
\citep[Theorem 3.1, printed p.~54]{BoucheronLugosiMassart2013},
the bounded-differences exponential-moment bound
\citep[Theorem 6.2, displayed log-mgf conclusion in its proof, printed p.~167]{BoucheronLugosiMassart2013},
and Cayley's labeled-tree formula
\citep[Section 5.6, Theorem 5.40]{KellerTrotterAppliedCombinatorics}
supply the variance, concentration, and counting inputs used in
\cref{obj:thm:matched-risk-frontier}. Applying that result at each \(n\ge2\) gives
\[
2^{-20}r(n,\epsilon_n)
\le R_{n,\epsilon_n}
\le2^{18}r(n,\epsilon_n).
\]
The finite real bounds are regarded here as elements of \([0,\infty]\).
We first establish the numerical equivalence
\[
r(n,\epsilon_n)\longrightarrow0
\quad\Longleftrightarrow\quad
n\epsilon_n\longrightarrow+\infty.
\]
% lean: CausalSmith.Stat.PrivateCateRoughdesign.matched_risk_frontier

\item Suppose that \(r(n,\epsilon_n)\to0\).
For every \(n\ge2\), the dense term is bounded by the uncapped maximum, and monotonicity of the minimum therefore gives
\[
\min\{1,(n\epsilon_n)^{-1/2}\}
\le r(n,\epsilon_n).
\]
For an arbitrary real threshold, set
\[
K_{\mathrm{test}}\in\mathbb R,
\qquad
B_{\mathrm{test}}=\max\{1,K_{\mathrm{test}}\}.
\]
Then
\[
B_{\mathrm{test}}\ge1,
\qquad
\min\{1,B_{\mathrm{test}}^{-1/2}\}>0.
\]
Convergence of the benchmark implies that, eventually,
\[
\min\{1,(n\epsilon_n)^{-1/2}\}
\le r(n,\epsilon_n)
<\min\{1,B_{\mathrm{test}}^{-1/2}\}.
\]
This forces
\[
(n\epsilon_n)^{-1/2}<B_{\mathrm{test}}^{-1/2}.
\]
Indeed, the opposite weak inequality would, by monotonicity of the minimum, imply
\[
\min\{1,B_{\mathrm{test}}^{-1/2}\}
\le\min\{1,(n\epsilon_n)^{-1/2}\},
\]
contradicting the strict bound above.
Both bases are positive, and the power \(-1/2\) is strictly decreasing on positive bases. Hence, eventually,
\[
n\epsilon_n>B_{\mathrm{test}}\ge K_{\mathrm{test}}.
\]
Since the threshold was arbitrary,
\[
n\epsilon_n\longrightarrow+\infty.
\]
% lean: CausalSmith.Stat.PrivateCateRoughdesign.budget_diverges_of_rate_tendsto_zero

\item Conversely, suppose that \(n\epsilon_n\to+\infty\).
First,
\[
n^{-1/4}\longrightarrow0.
\]
For \(n\ge2\), positivity of the budget and \(n\le n^2\) give
\[
n\epsilon_n\le n^2\epsilon_n.
\]
Thus
\[
n^2\epsilon_n\longrightarrow+\infty.
\]
Negative powers tend to zero as their positive bases tend to infinity, so
\[
(n^2\epsilon_n)^{-1/7}\longrightarrow0,
\qquad
(n\epsilon_n)^{-1/2}\longrightarrow0.
\]
Continuity of the finite maximum and minimum now gives
\[
\begin{aligned}
r(n,\epsilon_n)
&=\min\left\{1,\max\left\{
n^{-1/4},
\max\{(n^2\epsilon_n)^{-1/7},(n\epsilon_n)^{-1/2}\}
\right\}\right\}\\
&\longrightarrow
\min\{1,\max\{0,\max\{0,0\}\}\}=0.
\end{aligned}
\]
This completes the numerical equivalence.
% lean: CausalSmith.Stat.PrivateCateRoughdesign.rate_tendsto_zero_of_budget_diverges

\item It remains to transfer convergence between the benchmark and the risk.
Suppose first that \(R_{n,\epsilon_n}\to0\) in \([0,\infty]\).
For \(n\ge2\), positivity of the sampling term and of the cap gives
\[
0<r(n,\epsilon_n).
\]
Fix a tolerance
\[
\eta\in(0,\infty).
\]
Since \(2^{-20}\eta>0\), convergence of the risk implies that, eventually,
\[
R_{n,\epsilon_n}<2^{-20}\eta.
\]
Combining this with the lower frontier bound gives
\[
2^{-20}r(n,\epsilon_n)
\le R_{n,\epsilon_n}
<2^{-20}\eta.
\]
Cancellation of the positive constant yields
\[
0<r(n,\epsilon_n)<\eta.
\]
Therefore \(r(n,\epsilon_n)\to0\), and the numerical equivalence implies
\(n\epsilon_n\to+\infty\).

Conversely, if \(n\epsilon_n\to+\infty\), the numerical equivalence gives
\(r(n,\epsilon_n)\to0\). Consequently,
\[
2^{18}r(n,\epsilon_n)\longrightarrow0
\]
also when these finite values are viewed in \([0,\infty]\).
The upper frontier bound and nonnegativity give, for every \(n\ge2\),
\[
0\le R_{n,\epsilon_n}\le2^{18}r(n,\epsilon_n).
\]
The squeeze principle in \([0,\infty]\) yields
\[
R_{n,\epsilon_n}\longrightarrow0.
\]
Together these implications prove the asserted consistency equivalence.
% lean: CausalSmith.Stat.PrivateCateRoughdesign.private_risk_consistency_iff
\end{enumerate}
\end{proof}

\bibliographystyle{plainnat}

\bibliography{references}

\end{document}
